\input{preamble}

% OK, start here.
%
\begin{document}

\title{Cohomologie locale}


\maketitle

\phantomsection
\label{section-phantom}

\tableofcontents

\section{Introduction}
\label{section-introduction}

\noindent
Ce chapitre poursuit l'étude de la cohomologie locale. On pourra consulter \cite{SGA2}. La définition de la cohomologie locale se trouve dans Complexes dualisants, section \ref{dualizing-section-local-cohomology}. Pour les anneaux noethériens, prendre la cohomologie locale revient à dériver un foncteur de torsion convenable, comme le montre Complexes dualisants, section \ref{dualizing-section-local-cohomology-noetherian}. Le lien avec la profondeur est exposé dans Complexes dualisants, section \ref{dualizing-section-depth}.

\medskip\noindent
Nous étudions les propriétés de finitude de la cohomologie locale, ce qui conduit à une démonstration d'une version assez générale du théorème de finitude de Grothendieck ; voir le théorème \ref{theorem-finiteness} et le lemme \ref{lemma-finiteness-Rjstar} (images directes supérieures de modules cohérents par les immersions ouvertes). Nos méthodes font intervenir quelques arguments particulièrement élégants que le lecteur trouvera dans des articles de Faltings ; voir \cite{Faltings-annulators} et \cite{Faltings-finiteness}.

\medskip\noindent
Comme applications, nous présentons le théorème d'annulation de Hartshorne-Lichtenbaum. Nous étudions aussi l'action du Frobenius et des opérateurs différentiels sur la cohomologie locale.




\section{Généralités}
\label{section-generalities}

\noindent
Le lemme suivant relie le foncteur $R\Gamma_Z$ à la cohomologie à supports.

\begin{lemma}
\label{lemma-local-cohomology-is-local-cohomology}
Soient $A$ un anneau et $I$ un idéal de type fini. Posons $Z = V(I) \subset X = \Spec(A)$. Pour $K \in D(A)$ correspondant à $\widetilde{K} \in D_\QCoh(\mathcal{O}_X)$ par Catégories dérivées des schémas, lemme \ref{perfect-lemma-affine-compare-bounded}, il existe un isomorphisme fonctoriel
$$
R\Gamma_Z(K) = R\Gamma_Z(X, \widetilde{K})
$$
où le membre de gauche est donné par Complexes dualisants, équation (\ref{dualizing-equation-local-cohomology}), et celui de droite par le foncteur de Cohomologie, section \ref{cohomology-section-cohomology-support-bis}.
\end{lemma}

\begin{proof}
D'après Cohomologie, lemme \ref{cohomology-lemma-triangle-sections-with-support}, il existe un triangle distingué
$$
R\Gamma_Z(X, \widetilde{K})
\to R\Gamma(X, \widetilde{K})
\to R\Gamma(U, \widetilde{K})
\to R\Gamma_Z(X, \widetilde{K})[1]
$$
où $U = X \setminus Z$. On sait que $R\Gamma(X, \widetilde{K}) = K$ d'après Catégories dérivées des schémas, lemme \ref{perfect-lemma-affine-compare-bounded}. Écrivons $I = (f_1, \ldots, f_r)$. On obtient alors un recouvrement ouvert affine fini $\mathcal{U} : U = D(f_1) \cup \ldots \cup D(f_r)$. D'après Catégories dérivées des schémas, lemme \ref{perfect-lemma-alternating-cech-complex-complex-computes-cohomology}, le complexe de {\v C}ech alterné $\text{Tot}(\check{\mathcal{C}}_{alt}^\bullet(\mathcal{U},
\widetilde{K^\bullet}))$ calcule $R\Gamma(U, \widetilde{K})$, où $K^\bullet$ est un complexe quelconque de $A$-modules représentant $K$. En explicitant les définitions, on trouve
$$
R\Gamma(U, \widetilde{K}) =
\text{Tot}\left(
K^\bullet \otimes_A
(\prod\nolimits_{i_0} A_{f_{i_0}} \to
\prod\nolimits_{i_0 < i_1} A_{f_{i_0}f_{i_1}} \to
\ldots \to A_{f_1\ldots f_r})\right)
$$
Il est clair que $K^\bullet = R\Gamma(X, \widetilde{K^\bullet}) \to
R\Gamma(U, \widetilde{K}^\bullet)$ est induit par l'application diagonale de $A$ dans $\prod A_{f_i}$. On en déduit donc
$$
R\Gamma_Z(X, \widetilde{K}) =
\text{Tot}\left(
K^\bullet \otimes_A
(A \to \prod\nolimits_{i_0} A_{f_{i_0}} \to
\prod\nolimits_{i_0 < i_1} A_{f_{i_0}f_{i_1}} \to
\ldots \to A_{f_1\ldots f_r})\right)
$$
D'après Complexes dualisants, lemme \ref{dualizing-lemma-local-cohomology-adjoint}, le complexe du membre de droite calcule $R\Gamma_Z(K)$, ce qui prouve le lemme pour un $K$ donné.

\medskip\noindent
Pour terminer, montrons que l'isomorphisme peut être rendu fonctoriel en $K$. Considérons le morphisme de complexes de $A$-modules
$$
M^\bullet = \text{Tot}\left(
K^\bullet \otimes_A
(A \to \prod\nolimits_{i_0} A_{f_{i_0}} \to
\prod\nolimits_{i_0 < i_1} A_{f_{i_0}f_{i_1}} \to
\ldots \to A_{f_1\ldots f_r})\right)
\longrightarrow
K^\bullet
$$
Les arguments précédents montrent que cette flèche devient un isomorphisme après application du foncteur $R\Gamma_Z(X, \widetilde{\ })$ et que $R\Gamma(U, \widetilde{M}^\bullet) = 0$ (détails omis). On a alors
$$
R\Gamma_Z(X, \widetilde{K}^\bullet)
\leftarrow
R\Gamma_Z(X, \widetilde{M}^\bullet)
\rightarrow
R\Gamma(X, \widetilde{M}^\bullet) = M^\bullet = R\Gamma_Z(K^\bullet)
$$
et toutes les flèches sont des isomorphismes dans $D(A)$, fonctoriels en $K^\bullet$, comme voulu.
\end{proof}

\begin{lemma}
\label{lemma-local-cohomology}
Soient $A$ un anneau et $I \subset A$ un idéal de type fini. Posons $X = \Spec(A)$, $Z = V(I)$, $U = X \setminus Z$, et notons $j : U \to X$ le morphisme d'inclusion. Soit $\mathcal{F}$ un $\mathcal{O}_U$-module quasi-cohérent. Alors
\begin{enumerate}
\item il existe un $A$-module $M$ tel que $\mathcal{F}$ soit la restriction de $\widetilde{M}$ à $U$,
\item pour un tel $M$, on a une suite exacte
$$
0 \to H^0_Z(M) \to M \to H^0(U, \mathcal{F}) \to H^1_Z(M) \to 0
$$
et des isomorphismes $H^p(U, \mathcal{F}) = H^{p + 1}_Z(M)$ pour $p \geq 1$,
\item on peut prendre $M = H^0(U, \mathcal{F})$, auquel cas on a $H^0_Z(M) = H^1_Z(M) = 0$.
\end{enumerate}
\end{lemma}

\begin{proof}
L'existence de $M$ résulte de Propriétés, lemme \ref{properties-lemma-extend-trivial}, et du fait que les faisceaux quasi-cohérents sur $X$ correspondent aux $A$-modules (Schémas, lemme \ref{schemes-lemma-equivalence-quasi-coherent}). Considérons alors le triangle distingué
$$
R\Gamma_Z(X, \widetilde{M}) \to R\Gamma(X, \widetilde{M}) \to
R\Gamma(U, \widetilde{M}|_U) \to R\Gamma_Z(X, \widetilde{M})[1]
$$
de Cohomologie, lemme \ref{cohomology-lemma-triangle-sections-with-support}. Puisque $X$ est affine, on a $R\Gamma(X, \widetilde{M}) = M$ d'après Cohomologie des schémas, lemme \ref{coherent-lemma-quasi-coherent-affine-cohomology-zero}. Par le choix de $M$, on a $\mathcal{F} = \widetilde{M}|_U$, d'où une suite exacte
$$
0 \to H^0_Z(X, \widetilde{M}) \to M \to H^0(U, \mathcal{F}) \to
H^1_Z(X, \widetilde{M}) \to 0
$$
et des isomorphismes $H^p(U, \mathcal{F}) = H^{p + 1}_Z(X, \widetilde{M})$ pour $p \geq 1$. D'après le lemme \ref{lemma-local-cohomology-is-local-cohomology}, on a $H^i_Z(M) = H^i_Z(X, \widetilde{M})$ pour tout $i$. Ainsi (1) et (2) sont vérifiés. Enfin, en posant $M' = H^0(U, \mathcal{F})$, on voit que le noyau et le conoyau de $M \to M'$ sont de torsion par rapport aux puissances de $I$. Par conséquent, $\widetilde{M}|_U \to \widetilde{M'}|_U$ est un isomorphisme et l'on peut bien utiliser $M'$ comme indiqué en (3). Il va de soi que $H^0_Z(M')$ et $H^0_Z(M')$ sont tous deux nuls.
\end{proof}

\begin{lemma}
\label{lemma-already-torsion}
Soient $I, J \subset A$ des idéaux de type fini d'un anneau $A$. Si $M$ est un module de torsion par rapport aux puissances de $I$, alors l'application canonique
$$
H^i_{V(I) \cap V(J)}(M) \to H^i_{V(J)}(M)
$$
est un isomorphisme pour tout $i$.
\end{lemma}

\begin{proof}
Utiliser la suite spectrale de Complexes dualisants, lemme \ref{dualizing-lemma-local-cohomology-ss}, pour se ramener à l'énoncé $R\Gamma_I(M) = M$, qui résulte immédiatement de la construction de la cohomologie locale dans Complexes dualisants, section \ref{dualizing-section-local-cohomology}.
\end{proof}

\begin{lemma}
\label{lemma-multiplicative}
Soit $S \subset A$ une partie multiplicative d'un anneau $A$. Soit $M$ un $A$-module tel que $S^{-1}M = 0$. Alors $\colim_{f \in S} H^0_{V(f)}(M) = M$ et $\colim_{f \in S} H^1_{V(f)}(M) = 0$.
\end{lemma}

\begin{proof}
{\emergencystretch=2em
L'énoncé concernant $H^0$ résulte directement des définitions. Pour celui concernant $H^1$, observer que $R\Gamma_{V(f)}$ et $H^1_{V(f)}$ commutent aux limites inductives. On peut donc supposer $M$ annulé par un certain $f \in S$. Alors $H^1_{V(ff')}(M) = 0$ pour tout $f' \in S$ (par exemple d'après le lemme \ref{lemma-already-torsion}).
\par}
\end{proof}

\begin{lemma}
\label{lemma-elements-come-from-bigger}
Soit $I \subset A$ un idéal de type fini d'un anneau $A$. Soit $\mathfrak p$ un idéal premier. Soit $M$ un $A$-module. Soit $i \geq 0$ un entier et considérons l'application
$$
\Psi :
\colim_{f \in A, f \not \in \mathfrak p} H^i_{V((I, f))}(M)
\longrightarrow
H^i_{V(I)}(M)
$$
Alors
\begin{enumerate}
\item $\Im(\Psi)$ est l'ensemble des éléments dont l'image dans $H^i_{V(I)}(M)_\mathfrak p$ est nulle,
\item si $H^{i - 1}_{V(I)}(M)_\mathfrak p = 0$, alors $\Psi$ est injective,
\item si $H^{i - 1}_{V(I)}(M)_\mathfrak p = H^i_{V(I)}(M)_\mathfrak p = 0$, alors $\Psi$ est un isomorphisme.
\end{enumerate}
\end{lemma}

\begin{proof}
Pour $f \in A$, $f \not \in \mathfrak p$, la suite spectrale de Complexes dualisants, lemme \ref{dualizing-lemma-local-cohomology-ss}, dégénère et fournit des suites exactes courtes
$$
0 \to H^1_{V(f)}(H^{i - 1}_{V(I)}(M)) \to
H^i_{V((I, f))}(M) \to H^0_{V(f)}(H^i_{V(I)}(M)) \to 0
$$
Cela prouve (1) ; le point (2) en résulte à l'aide du lemme \ref{lemma-multiplicative}. Le point (3) en est une conséquence formelle.
\end{proof}

\begin{lemma}
\label{lemma-isomorphism}
Soient $I \subset I' \subset A$ des idéaux de type fini d'un anneau noethérien $A$. Soit $M$ un $A$-module. Soit $i \geq 0$ un entier. Considérons l'application
$$
\Psi : H^i_{V(I')}(M) \to H^i_{V(I)}(M)
$$
On a les assertions suivantes :
\begin{enumerate}
\item si $H^i_{\mathfrak pA_\mathfrak p}(M_\mathfrak p) = 0$ pour tout $\mathfrak p \in V(I) \setminus V(I')$, alors $\Psi$ est surjective,
\item si $H^{i - 1}_{\mathfrak pA_\mathfrak p}(M_\mathfrak p) = 0$ pour tout $\mathfrak p \in V(I) \setminus V(I')$, alors $\Psi$ est injective,
\item si $H^i_{\mathfrak pA_\mathfrak p}(M_\mathfrak p) =
H^{i - 1}_{\mathfrak pA_\mathfrak p}(M_\mathfrak p) = 0$ pour tout $\mathfrak p \in V(I) \setminus V(I')$, alors $\Psi$ est un isomorphisme.
\end{enumerate}
\end{lemma}

\begin{proof}
Démonstration de (1). Soit $\xi \in H^i_{V(I)}(M)$. Puisque $A$ est noethérien, il existe un plus grand idéal $I \subset I'' \subset I'$ tel que $\xi$ soit l'image d'un certain $\xi'' \in H^i_{V(I'')}(M)$. Si $V(I'') = V(I')$, la démonstration est terminée. Sinon, choisissons un point générique $\mathfrak p \in V(I'')$ n'appartenant pas à $V(I')$. On a alors $H^i_{V(I'')}(M)_\mathfrak p =
H^i_{\mathfrak pA_\mathfrak p}(M_\mathfrak p) = 0$ par hypothèse. Le lemme \ref{lemma-elements-come-from-bigger} permet d'agrandir $I''$, ce qui contredit sa maximalité.

\medskip\noindent
Démonstration de (2). Soit $\xi' \in H^i_{V(I')}(M)$ dans le noyau de $\Psi$. Puisque $A$ est noethérien, il existe un plus grand idéal $I \subset I'' \subset I'$ tel que $\xi'$ ait une image nulle dans $H^i_{V(I'')}(M)$. Si $V(I'') = V(I')$, la démonstration est terminée. Sinon, choisissons un point générique $\mathfrak p  \in V(I'')$ n'appartenant pas à $V(I')$. On a alors $H^{i - 1}_{V(I'')}(M)_\mathfrak p =
H^{i - 1}_{\mathfrak pA_\mathfrak p}(M_\mathfrak p) = 0$ par hypothèse. Le lemme \ref{lemma-elements-come-from-bigger} permet d'agrandir $I''$, ce qui contredit sa maximalité.

\medskip\noindent
Le point (3) résulte formellement de (1) et (2).
\end{proof}





\section{Lemme de connexité de Hartshorne}
\label{section-hartshorne-connectedness}

\noindent
Le titre de cette section renvoie au résultat suivant.

\begin{lemma}
\label{lemma-depth-2-connected-punctured-spectrum}
\begin{reference}
\cite[Proposition 2.1]{Hartshorne-connectedness}
\end{reference}
\begin{slogan}
Connexité de Hartshorne
\end{slogan}
Soit $A$ un anneau local noethérien de profondeur $\geq 2$. Alors les spectres épointés de $A$, $A^h$ et $A^{sh}$ sont connexes.
\end{lemma}

\begin{proof}
Soit $U$ le spectre épointé de $A$. Si $U$ n'est pas connexe, alors $\Gamma(U, \mathcal{O}_U)$ possède un idempotent non trivial. Or $A$, étant local, n'en possède pas. Par conséquent, $A \to \Gamma(U, \mathcal{O}_U)$ n'est pas un isomorphisme. Le lemme \ref{lemma-local-cohomology} montre que l'un au moins de $H^0_\mathfrak m(A)$ et $H^1_\mathfrak m(A)$ est non nul. Ainsi $\text{depth}(A) \leq 1$ d'après Complexes dualisants, lemme \ref{dualizing-lemma-depth}. Pour obtenir le résultat pour $A^h$ et $A^{sh}$, utiliser Compléments d'algèbre, lemme \ref{more-algebra-lemma-henselization-depth}.
\end{proof}

\begin{lemma}
\label{lemma-catenary-S2-equidimensional}
\begin{reference}
\cite[IV Corollary 5.10.9]{EGA}
\end{reference}
Soit $A$ un anneau local noethérien caténaire vérifiant $(S_2)$. Alors $\Spec(A)$ est équidimensionnel.
\end{lemma}

\begin{proof}
Posons $X = \Spec(A)$. Écrivons $d = \dim(A) = \dim(X)$. Dans $X$, considérons la réunion $X_1$ des composantes irréductibles de dimension $d$ et la réunion $X_2$ des composantes irréductibles de dimension $< d$. Bien entendu, $X = X_1 \cup X_2$. Si $X_2 = \emptyset$, le lemme est établi. Sinon, $Z = X_1 \cap X_2$ est une partie fermée non vide de $X$, puisqu'elle contient au moins le point fermé de $X$. On peut donc choisir un point générique $z \in Z$ d'une composante irréductible de $Z$. Rappelons que le spectre de $\mathcal{O}_{Z, z}$ est l'ensemble des points de $X$ qui se spécialisent en $z$. Puisque $z$ appartient à la fois à une composante irréductible de dimension $d$ et à une composante irréductible de dimension $< d$, on obtient des spécialisations non triviales $x_1 \leadsto z$ et $x_2 \leadsto z$ telles que les adhérences de $x_1$ et $x_2$ aient des dimensions différentes. Puisque $X$ est caténaire, cela n'est possible que si l'une au moins des spécialisations $x_1 \leadsto z$ et $x_2 \leadsto z$ n'est pas immédiate ! Ainsi $\dim(\mathcal{O}_{Z, z}) \geq 2$. Par conséquent, $\text{depth}(\mathcal{O}_{Z, z}) \geq 2$ puisque $A$ vérifie $(S_2)$. Cependant, le spectre épointé $U$ de $\mathcal{O}_{Z, z}$ n'est pas connexe, car les parties fermées $U \cap X_1$ et $U \cap X_2$ sont disjointes (par le choix de $z$) et recouvrent $U$. Cela contredit le lemme \ref{lemma-depth-2-connected-punctured-spectrum} et achève la démonstration.
\end{proof}



\section{Dimension cohomologique}
\label{section-cd}

\noindent
Une brève section consacrée à la dimension cohomologique.

\begin{lemma}
\label{lemma-cd}
Soit $I \subset A$ un idéal de type fini d'un anneau $A$. Posons $Y = V(I) \subset X = \Spec(A)$. Soit $d \geq -1$ un entier. Les conditions suivantes sont équivalentes :
\begin{enumerate}
\item $H^i_Y(A) = 0$ pour $i > d$,
\item $H^i_Y(M) = 0$ pour $i > d$ et pour tout $A$-module $M$,
\item si $d = -1$, alors $Y = \emptyset$ ; si $d = 0$, alors $Y$ est ouvert et fermé dans $X$ ; et si $d > 0$, alors $H^i(X \setminus Y, \mathcal{F}) = 0$ pour $i \geq d$ et pour tout $\mathcal{O}_{X \setminus Y}$-module quasi-cohérent $\mathcal{F}$.
\end{enumerate}
\end{lemma}

\begin{proof}
Observons que $R\Gamma_Y(-)$ est de dimension cohomologique finie, par exemple d'après Complexes dualisants, lemme \ref{dualizing-lemma-local-cohomology-adjoint}. Il existe donc un entier $i_0$ tel que $H^i_Y(M) = 0$ pour tous les $A$-modules $M$ et tout $i \geq i_0$.

\medskip\noindent
Montrons que (1) et (2) sont équivalents. Il est immédiat que (2) implique (1). Supposons (1). Par récurrence descendante sur $i > d$, montrons que $H^i_Y(M) = 0$ pour tous les $A$-modules $M$. Pour $i \geq i_0$, cela a été vu ci-dessus. Pour effectuer l'étape de récurrence, soit $i_0 > i > d$. Choisissons un $A$-module quelconque $M$ et insérons-le dans une suite exacte courte $0 \to N \to F \to M \to 0$, où $F$ est un $A$-module libre. Puisque $R\Gamma_Y$ est un adjoint à droite, on voit que $H^i_Y(-)$ commute aux sommes directes. Ainsi $H^i_Y(F) = 0$ puisque $i > d$, par l'hypothèse (1). On voit alors que $H^i_Y(M) = H^{i + 1}_Y(N) = 0$, comme voulu.

\medskip\noindent
Supposons $d = -1$ et (2) vérifié. Alors $0 = H^0_Y(A/I) = A/I \Rightarrow A = I
\Rightarrow Y = \emptyset$. Ainsi (3) est vérifié. Nous omettons la démonstration de la réciproque.

\medskip\noindent
{\emergencystretch=4em
Supposons $d = 0$ et (2) vérifié. Posons $J = H^0_I(A) = \{x \in A \mid I^nx = 0 \text{ pour un certain }n > 0\}$. Alors
\par}
$$
H^1_Y(A) = \Coker(A \to \Gamma(X \setminus Y, \mathcal{O}_{X \setminus Y}))
\quad\text{et}\quad
H^1_Y(I) = \Coker(I \to \Gamma(X \setminus Y, \mathcal{O}_{X \setminus Y}))
$$
et le noyau de la première application est égal à $J$. Voir le lemme \ref{lemma-local-cohomology}. On déduit de (2) que $I(A/J) = A/J$. On peut donc choisir $f \in I$ d'image $1$ dans $A/J$. Alors $1 - f \in J$, donc $I^n(1 - f) = 0$ pour un certain $n > 0$. Par suite, $f^n = f^{n + 1}$. Ainsi $e = f^n \in I$ est un idempotent. Considérons l'idempotent complémentaire $e' = 1 - f^n \in J$. Pour tout élément $g \in I$, on a $g^m e' = 0$ pour un certain $m > 0$. Ainsi $I$ est contenu dans le radical de l'idéal $(e) \subset I$. Cela signifie que $Y = V(I) = V(e)$ est ouvert et fermé dans $X$, comme prévu en (3). Réciproquement, si $Y = V(I)$ est ouvert et fermé, alors le foncteur $H^0_Y(-)$ est exact et ses foncteurs dérivés supérieurs sont nuls.

\medskip\noindent
Si $d > 0$, on voit immédiatement, d'après le lemme \ref{lemma-local-cohomology}, que (2) équivaut à (3).
\end{proof}

\begin{definition}
\label{definition-cd}
Soit $I \subset A$ un idéal de type fini d'un anneau $A$. Le plus petit entier $d \geq -1$ vérifiant les conditions équivalentes du lemme \ref{lemma-cd} est appelé la {\it dimension cohomologique de $I$ dans $A$} et se note $\text{cd}(A, I)$.
\end{definition}

\noindent
On a donc $\text{cd}(A, I) = -1$ si $I = A$ et $\text{cd}(A, I) = 0$ si $I$ est localement nilpotent ou engendré par un idempotent. Le lemme suivant montre que $\text{cd}(A, I)$ existe.

\begin{lemma}
\label{lemma-bound-cd}
Soit $I \subset A$ un idéal de type fini d'un anneau $A$. Alors
\begin{enumerate}
\item $\text{cd}(A, I)$ est au plus égal au nombre de générateurs de $I$,
\item $\text{cd}(A, I) \leq r$ s'il existe $f_1, \ldots, f_r \in A$ tels que $V(f_1, \ldots, f_r) = V(I)$,
\item $\text{cd}(A, I) \leq c$ si $\Spec(A) \setminus V(I)$ peut être recouvert par $c$ ouverts affines.
\end{enumerate}
\end{lemma}

\begin{proof}
La description explicite de $R\Gamma_Y(-)$ donnée dans Complexes dualisants, lemme \ref{dualizing-lemma-local-cohomology-adjoint}, prouve (1). On peut déduire (2) de (1) en utilisant le fait que $R\Gamma_Z$ ne dépend que de la partie fermée $Z$ et non du choix de l'idéal de type fini $I \subset A$ tel que $V(I) = Z$. Cela résulte soit de la construction de la cohomologie locale dans Complexes dualisants, section \ref{dualizing-section-local-cohomology}, combinée à Compléments d'algèbre, lemme \ref{more-algebra-lemma-local-cohomology-closed}, soit du lemme \ref{lemma-local-cohomology-is-local-cohomology}. Pour (3), on utilise le lemme \ref{lemma-cd} et le résultat d'annulation de Cohomologie des schémas, lemme \ref{coherent-lemma-vanishing-nr-affines}.
\end{proof}

\begin{lemma}
\label{lemma-cd-sum}
Soient $I, J \subset A$ des idéaux de type fini d'un anneau $A$. Alors $\text{cd}(A, I + J) \leq \text{cd}(A, I) + \text{cd}(A, J)$.
\end{lemma}

\begin{proof}
Utiliser la définition et Complexes dualisants, lemme \ref{dualizing-lemma-local-cohomology-ss}.
\end{proof}

\begin{lemma}
\label{lemma-cd-change-rings}
Soit $A \to B$ un homomorphisme d'anneaux. Soit $I \subset A$ un idéal de type fini. Alors $\text{cd}(B, IB) \leq \text{cd}(A, I)$. Si $A \to B$ est fidèlement plat, on a l'égalité.
\end{lemma}

\begin{proof}
Utiliser la définition et Complexes dualisants, lemme \ref{dualizing-lemma-torsion-change-rings}.
\end{proof}

\begin{lemma}
\label{lemma-cd-local}
Soit $I \subset A$ un idéal de type fini d'un anneau $A$. Alors $\text{cd}(A, I) = \max \text{cd}(A_\mathfrak p, I_\mathfrak p)$.
\end{lemma}

\begin{proof}
Soient $Y = V(I)$ et $Y' = V(I_\mathfrak p) \subset \Spec(A_\mathfrak p)$. Rappelons que $R\Gamma_Y(A) \otimes_A A_\mathfrak p = R\Gamma_{Y'}(A_\mathfrak p)$ d'après Complexes dualisants, lemme \ref{dualizing-lemma-torsion-change-rings}. On conclut donc par Algèbre, lemme \ref{algebra-lemma-characterize-zero-local}.
\end{proof}

\begin{lemma}
\label{lemma-cd-dimension}
Soit $I \subset A$ un idéal de type fini d'un anneau $A$. Si $M$ est un $A$-module de type fini, alors $H^i_{V(I)}(M) = 0$ pour $i > \dim(\text{Supp}(M))$. En particulier, on a $\text{cd}(A, I) \leq \dim(A)$.
\end{lemma}

\begin{proof}
Prouvons d'abord la seconde assertion. Rappelons que $\dim(A)$ désigne la dimension de Krull. Le lemme \ref{lemma-cd-local} permet de supposer $A$ local. Si $V(I) = \emptyset$, le résultat est vrai. Si $V(I) \not = \emptyset$, alors $\dim(\Spec(A) \setminus V(I)) < \dim(A)$, puisque le point fermé n'y figure pas. Observons que $U = \Spec(A) \setminus V(I)$ est un ouvert quasi-compact de l'espace spectral $\Spec(A)$, donc est lui-même un espace spectral. Voir Algèbre, lemme \ref{algebra-lemma-spec-spectral}, et Topologie, lemme \ref{topology-lemma-spectral-sub}. Ainsi Cohomologie, proposition \ref{cohomology-proposition-cohomological-dimension-spectral}, donne $H^i(U, \mathcal{F}) = 0$ pour $i \geq \dim(A)$, d'où le résultat cherché par le lemme \ref{lemma-cd}. Dans le cas noethérien, le lecteur peut utiliser le théorème de Grothendieck de Cohomologie, proposition \ref{cohomology-proposition-vanishing-Noetherian}.

\medskip\noindent
Déduisons la première assertion de la seconde. Soit $\mathfrak a$ l'annulateur du $A$-module de type fini $M$. Posons $B = A/\mathfrak a$. Rappelons que $\Spec(B) = \text{Supp}(M)$ ; voir Algèbre, lemme \ref{algebra-lemma-support-closed}. Posons $J = IB$. Alors $M$ est un $B$-module et $H^i_{V(I)}(M) = H^i_{V(J)}(M)$ ; voir Complexes dualisants, lemme \ref{dualizing-lemma-local-cohomology-and-restriction}. Puisque $\text{cd}(B, J) \leq \dim(B) = \dim(\text{Supp}(M))$ d'après la première partie de la démonstration, on conclut.
\end{proof}

\begin{lemma}
\label{lemma-cd-is-one}
Soit $I \subset A$ un idéal de type fini d'un anneau $A$. Si $\text{cd}(A, I) = 1$, alors $\Spec(A) \setminus V(I)$ est affine et non vide.
\end{lemma}

\begin{proof}
Cela résulte du lemme \ref{lemma-cd} et de Cohomologie des schémas, lemme \ref{coherent-lemma-quasi-compact-h1-zero-covering}.
\end{proof}

\begin{lemma}
\label{lemma-cd-maximal}
Soit $(A, \mathfrak m)$ un anneau local noethérien de dimension $d$. Alors $H^d_\mathfrak m(A)$ est non nul et $\text{cd}(A, \mathfrak m) = d$.
\end{lemma}

\begin{proof}
D'après l'une des caractérisations de la dimension, il existe un idéal de définition de $A$ engendré par $d$ éléments ; voir Algèbre, proposition \ref{algebra-proposition-dimension}. Ainsi $\text{cd}(A, \mathfrak m) \leq d$ d'après le lemme \ref{lemma-bound-cd}. Par conséquent, $H^d_\mathfrak m(A)$ est non nul si et seulement si $\text{cd}(A, \mathfrak m) = d$, ou encore si et seulement si $\text{cd}(A, \mathfrak m) \geq d$.

\medskip\noindent
Soit $A \to A^\wedge$ l'application de $A$ dans son complété. Observons que $A^\wedge$ est un anneau local noethérien de même dimension que $A$, d'idéal maximal $\mathfrak m A^\wedge$. Voir Algèbre, lemmes \ref{algebra-lemma-completion-Noetherian-Noetherian}, \ref{algebra-lemma-completion-complete} et \ref{algebra-lemma-completion-faithfully-flat}, et Compléments d'algèbre, lemme \ref{more-algebra-lemma-completion-dimension}. D'après le lemme \ref{lemma-cd-change-rings}, il suffit de démontrer le lemme pour $A^\wedge$.

\medskip\noindent
D'après le paragraphe précédent, on peut supposer que $A$ est un anneau local complet. Alors $A$ possède un complexe dualisant normalisé $\omega_A^\bullet$ (Complexes dualisants, lemme \ref{dualizing-lemma-ubiquity-dualizing}). Le théorème de dualité locale (sous la forme de Complexes dualisants, lemme \ref{dualizing-lemma-special-case-local-duality}) dit que $H^d_\mathfrak m(A)$ est le dual de Matlis de $\text{Ext}^{-d}(A, \omega_A^\bullet) = H^{-d}(\omega_A^\bullet)$, qui est non nul, par exemple d'après Complexes dualisants, lemme \ref{dualizing-lemma-nonvanishing-generically-local}.
\end{proof}

\begin{lemma}
\label{lemma-cd-bound-dim-local}
Soit $(A, \mathfrak m)$ un anneau local noethérien. Soit $I \subset A$ un idéal propre. Soit $\mathfrak p \subset A$ un idéal premier tel que $V(\mathfrak p) \cap V(I) = \{\mathfrak m\}$. Alors $\dim(A/\mathfrak p) \leq \text{cd}(A, I)$.
\end{lemma}

\begin{proof}
D'après le lemme \ref{lemma-cd-change-rings}, on a $\text{cd}(A, I) \geq \text{cd}(A/\mathfrak p, I(A/\mathfrak p))$. Puisque $V(I) \cap V(\mathfrak p) = \{\mathfrak m\}$, on a $\text{cd}(A/\mathfrak p, I(A/\mathfrak p)) =
\text{cd}(A/\mathfrak p, \mathfrak m/\mathfrak p)$. Le lemme \ref{lemma-cd-maximal} montre que cette quantité est égale à $\dim(A/\mathfrak p)$.
\end{proof}

\begin{lemma}
\label{lemma-cd-blowup}
Soit $A$ un anneau noethérien. Soit $I \subset A$ un idéal. Soit $b : X' \to X = \Spec(A)$ l'éclatement de $I$. Si les fibres de $b$ sont de dimension $\leq d - 1$, alors $\text{cd}(A, I) \leq d$.
\end{lemma}

\begin{proof}
Posons $U = X \setminus V(I)$. Notons $j : U \to X'$ l'immersion ouverte canonique ; voir Diviseurs, section \ref{divisors-section-blowing-up}. Puisque le diviseur exceptionnel est un diviseur de Cartier effectif (Diviseurs, lemme \ref{divisors-lemma-blowing-up-gives-effective-Cartier-divisor}), $j$ est affine ; voir Diviseurs, lemme \ref{divisors-lemma-complement-locally-principal-closed-subscheme}. Soit $\mathcal{F}$ un $\mathcal{O}_U$-module quasi-cohérent. Alors $R^pj_*\mathcal{F} = 0$ pour $p > 0$ ; voir Cohomologie des schémas, lemme \ref{coherent-lemma-relative-affine-vanishing}. D'autre part, on a $R^qb_*(j_*\mathcal{F}) = 0$ pour $q \geq d$ d'après Limites, lemme \ref{limits-lemma-higher-direct-images-zero-above-dimension-fibre}. Par la suite spectrale de Leray (Cohomologie, lemme \ref{cohomology-lemma-relative-Leray}), on en déduit $R^n(b \circ j)_*\mathcal{F} = 0$ pour $n \geq d$. Ainsi $H^n(U, \mathcal{F}) = 0$ pour $n \geq d$ (d'après Cohomologie, lemme \ref{cohomology-lemma-apply-Leray}). Par définition, cela signifie que $\text{cd}(A, I) \leq d$.
\end{proof}







\section{Supports plus généraux}
\label{section-supports}

\noindent
Soit $A$ un anneau noethérien. Soit $M$ un $A$-module. Soit $T \subset \Spec(A)$ une partie stable par spécialisation (Topologie, définition \ref{topology-definition-specialization}). Définissons
$$
H^0_T(M) = \colim_{Z \subset T} H^0_Z(M)
$$
où la limite inductive porte sur l'ensemble ordonné filtrant des parties fermées $Z$ de $\Spec(A)$ contenues dans $T$\footnote{Puisque $T$ est stable par spécialisation, on a $T = \bigcup_{Z \subset T} Z$ ; voir Topologie, lemme \ref{topology-lemma-stable-specialization}.}. Autrement dit, un élément $m$ de $M$ appartient à $H^0_T(M) \subset M$ si et seulement si le support $V(\text{Ann}_R(m))$ de $m$ est contenu dans $T$.

\begin{lemma}
\label{lemma-support}
Soit $A$ un anneau noethérien. Soit $T \subset \Spec(A)$ une partie stable par spécialisation. Pour un $A$-module $M$, les conditions suivantes sont équivalentes :
\begin{enumerate}
\item $H^0_T(M) = M$,
\item $\text{Supp}(M) \subset T$.
\end{enumerate}
La catégorie de ces $A$-modules est une sous-catégorie de Serre de la catégorie des $A$-modules, stable par sommes directes.
\end{lemma}

\begin{proof}
L'équivalence tient au fait que le support d'un élément de $M$ est contenu dans le support de $M$ et que, réciproquement, le support de $M$ est la réunion des supports de ses éléments. La catégorie de ces modules est une sous-catégorie de Serre (Homologie, définition \ref{homology-definition-serre-subcategory}) de $\text{Mod}_A$ d'après Algèbre, lemme \ref{algebra-lemma-support-quotient}. Nous omettons la démonstration de l'assertion concernant les sommes directes.
\end{proof}

\noindent
Soit $A$ un anneau noethérien. Soit $T \subset \Spec(A)$ une partie stable par spécialisation. Notons $\text{Mod}_{A, T} \subset \text{Mod}_A$ la sous-catégorie de Serre décrite au lemme \ref{lemma-support}. Notons $D_T(A) \subset D(A)$ la sous-catégorie triangulée strictement pleine et saturée de $D(A)$ (Catégories dérivées, lemme \ref{derived-lemma-cohomology-in-serre-subcategory}) constituée des complexes de $A$-modules dont les modules de cohomologie appartiennent à $\text{Mod}_{A, T}$. On obtient des foncteurs
$$
D(\text{Mod}_{A, T}) \to D_T(A) \to D(A)
$$
Voir la discussion dans Catégories dérivées, section \ref{derived-section-triangulated-sub}. Notons $RH^0_T : D(A) \to D(\text{Mod}_{A, T})$ le prolongement dérivé à droite de $H^0_T$. Nous noterons
$$
R\Gamma_T : D^+(A) \to D^+_T(A),
$$
le composé de $RH^0_T : D^+(A) \to D^+(\text{Mod}_{A, T})$ et de $D^+(\text{Mod}_{A, T}) \to D^+_T(A)$. Si la dimension de $A$ est finie\footnote{Si $\dim(A) = \infty$, la construction peut présenter des propriétés inattendues sur les complexes non bornés.}, nous noterons
$$
R\Gamma_T : D(A) \to D_T(A)
$$
le composé de $RH^0_T$ et de $D(\text{Mod}_{A, T}) \to D_T(A)$.

\begin{lemma}
\label{lemma-adjoint}
Soit $A$ un anneau noethérien. Soit $T \subset \Spec(A)$ une partie stable par spécialisation. Le foncteur $RH^0_T$ est l'adjoint à droite du foncteur $D(\text{Mod}_{A, T}) \to D(A)$.
\end{lemma}

\begin{proof}
Cela résulte du fait que le foncteur $H^0_T(-)$ est l'adjoint à droite du foncteur d'inclusion $\text{Mod}_{A, T} \to \text{Mod}_A$ ; voir Catégories dérivées, lemme \ref{derived-lemma-derived-adjoint-functors}.
\end{proof}

\begin{lemma}
\label{lemma-adjoint-ext}
Soit $A$ un anneau noethérien. Soit $T \subset \Spec(A)$ une partie stable par spécialisation. Pour tout objet $K$ de $D(A)$, on a
$$
H^i(RH^0_T(K)) = \colim_{Z \subset T\text{ fermé}} H^i_Z(K)
$$
\end{lemma}

\begin{proof}
Soit $J^\bullet$ un complexe K-injectif représentant $K$. Par définition, $RH^0_T$ est représenté par le complexe
$$
H^0_T(J^\bullet) = \colim H^0_Z(J^\bullet)
$$
où l'égalité résulte de notre définition de $H^0_T$. Puisque les limites inductives filtrantes sont exactes, la cohomologie de ce complexe en degré $i$ est $\colim H^i(H^0_Z(J^\bullet)) = \colim H^i_Z(K)$, comme voulu.
\end{proof}

\begin{lemma}
\label{lemma-equal-plus}
Soit $A$ un anneau noethérien. Soit $T \subset \Spec(A)$ une partie stable par spécialisation. Le foncteur $D^+(\text{Mod}_{A, T}) \to D^+_T(A)$ est une équivalence.
\end{lemma}

\begin{proof}
Soit $M$ un objet de $\text{Mod}_{A, T}$. Choisissons un plongement $M \to J$ dans un $A$-module injectif. D'après Complexes dualisants, proposition \ref{dualizing-proposition-structure-injectives-noetherian}, le module $J$ est une somme directe d'enveloppes injectives de corps résiduels. Soit $E$ une enveloppe injective du corps résiduel de $\mathfrak p$. Puisque $E$ est de torsion par rapport aux puissances de $\mathfrak p$, on voit que $H^0_T(E) = 0$ si $\mathfrak p \not \in T$ et $H^0_T(E) = E$ si $\mathfrak p \in T$. Ainsi $H^0_T(J)$ est injectif, comme somme directe d'enveloppes injectives (d'après la proposition), et l'on a un plongement $M \to H^0_T(J)$. Tout objet $M$ de $\text{Mod}_{A, T}$ possède donc une résolution injective $M \to J^\bullet$ avec $J^n$ appartenant aussi à $\text{Mod}_{A, T}$. Il s'ensuit que $RH^0_T(M) = M$.

\medskip\noindent
Supposons maintenant que $K \in D_T^+(A)$. Alors la suite spectrale
$$
R^qH^0_T(H^p(K)) \Rightarrow R^{p + q}H^0_T(K)
$$
(Catégories dérivées, lemme \ref{derived-lemma-two-ss-complex-functor}) converge, et nous avons vu ci-dessus que seuls les termes avec $q = 0$ sont non nuls. Ainsi $RH^0_T(K) \to K$ est un isomorphisme. Par conséquent, le foncteur $D^+(\text{Mod}_{A, T}) \to D^+_T(A)$ est une équivalence de quasi-inverse $RH^0_T$.
\end{proof}

\begin{lemma}
\label{lemma-equal-full}
Soit $A$ un anneau noethérien. Soit $T \subset \Spec(A)$ une partie stable par spécialisation. Si $\dim(A) < \infty$, alors le foncteur $D(\text{Mod}_{A, T}) \to D_T(A)$ est une équivalence.
\end{lemma}

\begin{proof}
Écrivons $\dim(A) = d$. Alors $H^i_Z(M) = 0$ pour $i > d$, pour toute partie fermée $Z$ de $\Spec(A)$ ; voir le lemme \ref{lemma-cd-dimension}. D'après le lemme \ref{lemma-adjoint-ext}, le foncteur $H^0_T$ est de dimension cohomologique bornée.

\medskip\noindent
Soit $K \in D_T(A)$. Montrons que $RH^0_T(K) \to K$ est un isomorphisme. Nous le savons lorsque $K$ est borné inférieurement ; voir le lemme \ref{lemma-equal-plus}. Or, puisque $H^0_T$ est de dimension cohomologique bornée, la cohomologie en degré $i$ de $RH_T^0(K)$ ne dépend que de $\tau_{\geq -d + i}K$, ce qui permet de conclure. Ainsi $D(\text{Mod}_{A, T}) \to D_T(A)$ est une équivalence de quasi-inverse $RH^0_T$.
\end{proof}

\begin{remark}
\label{remark-upshot}
Soit $A$ un anneau noethérien. Soit $T \subset \Spec(A)$ une partie stable par spécialisation. La discussion précédente montre que $R\Gamma_T : D^+(A) \to D_T^+(A)$ est l'adjoint à droite du foncteur d'inclusion $D_T^+(A) \to D^+(A)$. Si $\dim(A) < \infty$, alors $R\Gamma_T : D(A) \to D_T(A)$ est l'adjoint à droite du foncteur d'inclusion $D_T(A) \to D(A)$. Dans les deux cas, on a
$$
H^i_T(K) = H^i(R\Gamma_T(K)) = R^iH^0_T(K) =
\colim_{Z \subset T\text{ fermé}} H^i_Z(K)
$$
Cela résulte de la combinaison des lemmes \ref{lemma-adjoint}, \ref{lemma-adjoint-ext}, \ref{lemma-equal-plus} et \ref{lemma-equal-full}.
\end{remark}

\begin{lemma}
\label{lemma-torsion-change-rings}
Soit $A \to B$ un homomorphisme plat d'anneaux noethériens. Soit $T \subset \Spec(A)$ une partie stable par spécialisation. Soit $T' \subset \Spec(B)$ l'image réciproque de $T$. Alors l'application canonique
$$
R\Gamma_T(K) \otimes_A^\mathbf{L} B
\longrightarrow
R\Gamma_{T'}(K \otimes_A^\mathbf{L} B)
$$
est un isomorphisme pour $K \in D^+(A)$. Si $A$ et $B$ sont de dimension finie, cela est vrai pour $K \in D(A)$.
\end{lemma}

\begin{proof}
L'application $R\Gamma_T(K) \to K$ fournit une application $R\Gamma_T(K) \otimes_A^\mathbf{L} B \to K \otimes_A^\mathbf{L} B$. Les modules de cohomologie de $R\Gamma_T(K) \otimes_A^\mathbf{L} B$ ont leur support dans $T'$, d'où la flèche du lemme. Cette flèche est un isomorphisme lorsque $T$ est une partie fermée de $\Spec(A)$, d'après Complexes dualisants, lemme \ref{dualizing-lemma-torsion-change-rings}. Rappelons que $H^i_T(K)$ est la limite inductive des $H^i_Z(K)$, où $Z$ parcourt l'ensemble filtrant des parties fermées de $T$ ; voir le lemme \ref{lemma-adjoint-ext}. De même, $H^i_{T'}(K \otimes_A^\mathbf{L} B) =
\colim H^i_{Z'}(K \otimes_A^\mathbf{L} B)$, où $Z'$ est l'image réciproque de $Z$. Le résultat s'ensuit, car $\otimes_A B$ commute aux limites inductives filtrantes et les Tor supérieurs sont nuls.
\end{proof}

\begin{lemma}
\label{lemma-local-cohomology-ss}
Soient $A$ un anneau et $T, T' \subset \Spec(A)$ des parties stables par spécialisation. Pour $K \in D^+(A)$, il existe une suite spectrale
$$
E_2^{p, q} = H^p_T(H^p_{T'}(K)) \Rightarrow H^{p + q}_{T \cap T'}(K)
$$
comme dans Catégories dérivées, lemme \ref{derived-lemma-grothendieck-spectral-sequence}.
\end{lemma}

\begin{proof}
Soit $E$ un objet de $D_{T \cap T'}(A)$. Alors on a
$$
\Hom(E, R\Gamma_T(R\Gamma_{T'}(K))) =
\Hom(E, R\Gamma_{T'}(K)) =
\Hom(E, K)
$$
La première égalité résulte de la propriété d'adjonction de $R\Gamma_T$ et la seconde de celle de $R\Gamma_{T'}$. D'autre part, si $J^\bullet$ est un complexe d'injectifs borné inférieurement représentant $K$, alors $H^0_{T'}(J^\bullet)$ est un complexe de $A$-modules injectifs représentant $R\Gamma_{T'}(K)$ ; par conséquent, $H^0_T(H^0_{T'}(J^\bullet))$ est un complexe représentant $R\Gamma_T(R\Gamma_{T'}(K))$. Ainsi $R\Gamma_T(R\Gamma_{T'}(K))$ est un objet de $D^+_{T \cap T'}(A)$. En combinant ces deux faits, on trouve $R\Gamma_{T \cap T'} = R\Gamma_T \circ R\Gamma_{T'}$. Cela donne la suite spectrale, d'après le lemme cité dans l'énoncé.
\end{proof}

\begin{lemma}
\label{lemma-torsion-tensor-product}
Soit $A$ un anneau noethérien. Soit $T \subset \Spec(A)$ une partie stable par spécialisation. Supposons $A$ de dimension finie. Alors
$$
R\Gamma_T(K) = R\Gamma_T(A) \otimes_A^\mathbf{L} K
$$
pour $K \in D(A)$. Pour $K, L \in D(A)$, on a
$$
R\Gamma_T(K \otimes_A^\mathbf{L} L) =
K \otimes_A^\mathbf{L} R\Gamma_T(L) =
R\Gamma_T(K) \otimes_A^\mathbf{L} L =
R\Gamma_T(K) \otimes_A^\mathbf{L} R\Gamma_T(L)
$$
Si $K$ ou $L$ appartient à $D_T(A)$, alors $K \otimes_A^\mathbf{L} L$ y appartient aussi.
\end{lemma}

\begin{proof}
Par construction, on peut représenter $R\Gamma_T(A)$ par un complexe $J^\bullet$ dans $\text{Mod}_{A, T}$. Si l'on représente $K$ par un complexe K-plat $K^\bullet$, alors $R\Gamma_T(A) \otimes_A^\mathbf{L} K$ est représenté par le complexe $\text{Tot}(J^\bullet \otimes_A K^\bullet)$ dans $\text{Mod}_{A, T}$. L'application $R\Gamma_T(A) \to A$ fournit une application $R\Gamma_T(A) \otimes_A^\mathbf{L} K\to K$. Par la propriété d'adjonction de $R\Gamma_T$, on obtient donc une application canonique
$$
R\Gamma_T(A) \otimes_A^\mathbf{L} K \longrightarrow R\Gamma_T(K)
$$
par laquelle se factorise l'application que l'on vient de construire. Observons que $R\Gamma_T$ commute aux sommes directes dans $D(A)$, par exemple d'après le lemme \ref{lemma-adjoint-ext}, le fait que les limites inductives filtrantes commutent aux sommes directes, et le fait que la cohomologie locale usuelle commute aux sommes directes (par exemple d'après Complexes dualisants, lemme \ref{dualizing-lemma-local-cohomology-adjoint}). Ainsi, d'après Compléments d'algèbre, remarque \ref{more-algebra-remark-P-resolution}, il suffit de vérifier que l'application est un isomorphisme pour $K = A[k]$, où $k \in \mathbf{Z}$. Cela est clair.

\medskip\noindent
Les dernières assertions résultent de ce que nous venons de montrer, par transitivité des produits tensoriels dérivés.
\end{proof}





\section{Filtrations sur la cohomologie locale}
\label{section-filter-local-cohomology}

\noindent
Quelques procédés liés à la suite spectrale du lemme \ref{lemma-local-cohomology-ss}.

\begin{lemma}
\label{lemma-filter-local-cohomology}
Soit $A$ un anneau noethérien. Soit $T \subset \Spec(A)$ une partie stable par spécialisation. Soit $T' \subset T$ l'ensemble des idéaux premiers non minimaux dans $T$. Alors $T'$ est une partie de $\Spec(A)$ stable par spécialisation et, pour tout $A$-module $M$, il existe une suite exacte
$$
0 \to
\colim_{Z, f} H^1_f(H^{i - 1}_Z(M)) \to
H^i_{T'}(M) \to H^i_T(M) \to
\bigoplus\nolimits_{\mathfrak p \in T \setminus T'}
H^i_{\mathfrak p A_\mathfrak p}(M_\mathfrak p)
$$
où la limite inductive porte sur les parties fermées $Z \subset T$ et les $f \in A$ tels que $V(f) \cap Z \subset T'$.
\end{lemma}

\begin{proof}
Pour tous $Z$ et $f$, la suite spectrale de Complexes dualisants, lemme \ref{dualizing-lemma-local-cohomology-ss}, dégénère et fournit des suites exactes courtes
$$
0 \to H^1_f(H^{i - 1}_Z(M)) \to
H^i_{Z \cap V(f)}(M) \to H^0_f(H^i_Z(M)) \to 0
$$
Nous utiliserons ce fait ci-dessous sans le mentionner à nouveau.

\medskip\noindent
Soit $\xi \in H^i_T(M)$ d'image nulle dans la somme directe. Écrivons d'abord $\xi$ comme l'image d'un certain $\xi' \in H^i_Z(M)$ pour une partie fermée $Z \subset T$ ; voir le lemme \ref{lemma-adjoint-ext}. Alors $\xi'$ a une image nulle dans $H^i_{\mathfrak p A_\mathfrak p}(M_\mathfrak p)$ pour tout $\mathfrak p \in Z$, $\mathfrak p \not \in T'$. Puisque ces idéaux premiers sont en nombre fini, on peut choisir $f \in A$ n'appartenant à aucun d'entre eux et tel que $f$ annule $\xi'$. Alors $\xi'$ est l'image d'un certain $\xi'' \in H^i_{Z'}(M)$, où $Z' = Z \cap V(f)$. Le choix de $f$ assure que $Z' \subset T'$, d'où l'exactitude à l'avant-dernier terme.

\medskip\noindent
Soit $\xi \in H^i_{T'}(M)$ d'image nulle dans $H^i_T(M)$. Choisissons des parties fermées $Z' \subset Z$, avec $Z' \subset T'$ et $Z \subset T$, telles que $\xi$ provienne d'un élément $\xi' \in H^i_{Z'}(M)$ dont l'image dans $H^i_Z(M)$ est nulle. On peut alors trouver $f \in A$ tel que $V(f) \cap Z = Z'$, et l'on conclut.
\end{proof}

\begin{lemma}
\label{lemma-zero}
{\emergencystretch=2em
Soit $A$ un anneau noethérien de dimension finie. Soit $T \subset \Spec(A)$ une partie stable par spécialisation. Soit $\{M_n\}_{n \geq 0}$ un système projectif de $A$-modules. Soit $i \geq 0$ un entier. Supposons que pour tout $m$, il existe un entier $m'(m) \geq m$ tel que, pour tout $\mathfrak p \in T$, l'application induite
\par}
$$
H^i_{\mathfrak p A_\mathfrak p}(M_{k, \mathfrak p})
\longrightarrow
H^i_{\mathfrak p A_\mathfrak p}(M_{m, \mathfrak p})
$$
soit nulle pour $k \geq m'(m)$. Soit $m'' : \mathbf{N} \to \mathbf{N}$ la composée, $2^{\dim(T)}$ fois avec elle-même, de $m'$. Alors l'application $H^i_T(M_k) \to H^i_T(M_m)$ est nulle pour tout $k \geq m''(m)$.
\end{lemma}

\begin{proof}
Commençons par une remarque générale : supposons donnée une suite exacte
$$
(A_n) \to (B_n) \to (C_n)
$$
de systèmes projectifs de groupes abéliens. Supposons que pour tout $m$, il existe un entier $m'(m) \geq m$ tel que
$$
A_k \to A_m
\quad\text{et}\quad
C_k \to C_m
$$
soient nulles pour $k \geq m'(m)$. Alors, pour $k \geq m'(m'(m))$, l'application $B_k \to B_m$ est nulle.

\medskip\noindent
Démontrons le lemme par récurrence sur $\dim(T)$, qui est fini puisque $\dim(A)$ est fini. Soit $T' \subset T$ l'ensemble des idéaux premiers non minimaux dans $T$. Alors $T'$ est une partie de $\Spec(A)$ stable par spécialisation, et les hypothèses du lemme s'appliquent à $T'$. Puisque $\dim(T') < \dim(T)$, on sait que le lemme est vrai pour $T'$. Pour tout $A$-module $M$, il existe une suite exacte
$$
H^i_{T'}(M) \to H^i_T(M) \to
\bigoplus\nolimits_{\mathfrak p \in T \setminus T'}
H^i_{\mathfrak p A_\mathfrak p}(M_\mathfrak p)
$$
d'après le lemme \ref{lemma-filter-local-cohomology}. On conclut donc par la remarque initiale de la démonstration.
\end{proof}

\begin{lemma}
\label{lemma-essential-image}
Soit $A$ un anneau noethérien. Soit $T \subset \Spec(A)$ une partie stable par spécialisation. Soit $\{M_n\}_{n \geq 0}$ un système projectif de $A$-modules. Soit $i \geq 0$ un entier. Supposons que la dimension de $A$ soit finie et que, pour tout $m$, il existe un entier $m'(m) \geq m$ tel que, pour tout $\mathfrak p \in T$, les propriétés suivantes soient vérifiées :
\begin{enumerate}
\item $H^{i - 1}_{\mathfrak p A_\mathfrak p}(M_{k, \mathfrak p})
\to H^{i - 1}_{\mathfrak p A_\mathfrak p}(M_{m, \mathfrak p})$ est nulle pour $k \geq m'(m)$,
\item $ H^i_{\mathfrak p A_\mathfrak p}(M_{k, \mathfrak p}) \to
H^i_{\mathfrak p A_\mathfrak p}(M_{m, \mathfrak p})$ a pour image $G(\mathfrak p, m)$, indépendante de $k \geq m'(m)$ ; de plus, $G(\mathfrak p, m)$ s'envoie injectivement dans $H^i_{\mathfrak p A_\mathfrak p}(M_{0, \mathfrak p})$.
\end{enumerate}
Alors il existe un entier $m_0$ tel que, pour tout $m \geq m_0$, il existe un entier $m''(m) \geq m$ tel que, pour $k \geq m''(m)$, l'image de $H^i_T(M_k) \to H^i_T(M_m)$ s'envoie injectivement dans $H^i_T(M_{m_0})$.
\end{lemma}

\begin{proof}
Commençons par une remarque générale : supposons donnée une suite exacte
$$
(A_n) \to (B_n) \to (C_n) \to (D_n)
$$
de systèmes projectifs de groupes abéliens. Supposons qu'il existe un entier $m_0$ tel que, pour tout $m \geq m_0$, il existe un entier $m'(m) \geq m$ tel que les applications
$$
\Im(B_k \to B_m) \longrightarrow B_{m_0}
\quad\text{et}\quad
\Im(D_k \to D_m) \longrightarrow D_{m_0}
$$
soient injectives pour $k \geq m'(m)$ et que $A_k \to A_m$ soit nulle pour $k \geq m'(m)$. Alors, pour $m \geq m'(m_0)$ et $k \geq m'(m'(m))$, l'application
$$
\Im(C_k \to C_m) \to C_{m'(m_0)}
$$
est injective. En effet, soit $c_0 \in C_m$ l'image de $c_3 \in C_k$, et supposons que $c_0$ ait une image nulle dans $C_{m'(m_0)}$. Représentons la situation par
$$
C_k \to C_{m'(m'(m))} \to C_{m'(m)} \to C_m \to C_{m'(m_0)},\quad
c_3 \mapsto c_2 \mapsto c_1 \mapsto c_0 \mapsto 0
$$
Il faut montrer que $c_0 = 0$. L'image $d_3$ de $c_3$ a une image nulle dans $C_{m_0}$ ; on en déduit que l'image $d_1 \in D_{m'(m)}$ est nulle. On peut donc choisir $b_1 \in B_{m'(m)}$ s'envoyant sur l'image $c_1$. Puisque $c_3$ a une image nulle dans $C_{m'(m_0)}$, on trouve un élément $a_{-1} \in A_{m'(m_0)}$ qui s'envoie sur l'image $b_{-1} \in B_{m'(m_0)}$ de $b_1$. Puisque $a_{-1}$ a une image nulle dans $A_{m_0}$, on en déduit que $b_1$ a une image nulle dans $B_{m_0}$. Ainsi l'image $b_0 \in B_m$ est nulle, ce qui entraîne bien entendu $c_0 = 0$, comme voulu.

\medskip\noindent
Démontrons le lemme par récurrence sur $\dim(T)$, qui est fini puisque $\dim(A)$ est fini. Soit $T' \subset T$ l'ensemble des idéaux premiers non minimaux dans $T$. Alors $T'$ est une partie de $\Spec(A)$ stable par spécialisation, et les hypothèses du lemme s'appliquent à $T'$. Puisque $\dim(T') < \dim(T)$, on sait que le lemme est vrai pour $T'$. Pour tout $A$-module $M$, il existe une suite exacte
$$
0 \to \colim_{Z, f} H^1_f(H^{i - 1}_Z(M)) \to
H^i_{T'}(M) \to H^i_T(M) \to
\bigoplus\nolimits_{\mathfrak p \in T \setminus T'}
H^i_{\mathfrak p A_\mathfrak p}(M_\mathfrak p)
$$
d'après le lemme \ref{lemma-filter-local-cohomology}. On conclut donc par la remarque initiale de la démonstration et par le fait que le système de groupes
$$
\left\{\colim_{Z, f} H^1_f(H^{i - 1}_Z(M_n))\right\}_{n \geq 0}
$$
est pro-nul d'après le lemme \ref{lemma-zero} ; on utilise ici le fait que la fonction $m''(m)$ de ce lemme, appliqué à $H^{i - 1}_Z(M)$, est indépendante de $Z$.
\end{proof}










\section{Finitude de la cohomologie locale, I}
\label{section-finiteness}

\noindent
Nous suivrons l'approche de Faltings pour la finitude des modules de cohomologie locale ; voir \cite{Faltings-annulators} et \cite{Faltings-finiteness}. Le lemme suivant montre que, pour prouver que les modules de cohomologie locale sont de type fini, il suffit de montrer qu'ils possèdent un annulateur.

\begin{lemma}
\label{lemma-check-finiteness-local-cohomology-by-annihilator}
\begin{reference}
\cite[Lemma 3]{Faltings-annulators}
\end{reference}
Soit $A$ un anneau noethérien. Soit $T \subset \Spec(A)$ une partie stable par spécialisation. Soit $M$ un $A$-module de type fini. Soit $n \geq 0$. Les conditions suivantes sont équivalentes :
\begin{enumerate}
\item $H^i_T(M)$ est de type fini pour $i \leq n$,
\item il existe un idéal $J \subset A$ tel que $V(J) \subset T$ et tel que $J$ annule $H^i_T(M)$ pour $i \leq n$.
\end{enumerate}
Si $T = V(I) = Z$ pour un idéal $I \subset A$, ces conditions sont également équivalentes à
\begin{enumerate}
\item[(3)] il existe $e \geq 0$ tel que $I^e$ annule $H^i_Z(M)$ pour $i \leq n$.
\end{enumerate}
\end{lemma}

\begin{proof}
Démontrons l'équivalence de (1) et (2) par récurrence sur $n$. Pour $n = 0$, le module $H^0_T(M) \subset M$ est de type fini. Ainsi (1) est vrai. Puisque $H^0_T(M) = \colim H^0_{V(J)}(M)$, avec $J$ comme dans (2), on voit que (2) est vrai. Supposons que $n > 0$.

\medskip\noindent
Supposons (1) vrai. Rappelons que $H^i_J(M) = H^i_{V(J)}(M)$ ; voir Complexes dualisants, lemme \ref{dualizing-lemma-local-cohomology-noetherian}. Ainsi $H^i_T(M) = \colim H^i_J(M)$, où la limite inductive porte sur les idéaux $J \subset A$ tels que $V(J) \subset T$ ; voir le lemme \ref{lemma-adjoint-ext}. Puisque $H^i_T(M)$ est de type fini pour $i \leq n$, on peut trouver $J \subset A$ comme dans (2) tel que $H^i_J(M) \to H^i_T(M)$ soit surjective pour $i \leq n$. Les éléments de la liste finie de générateurs sont donc de torsion par rapport aux puissances de $J$, et (2) est vérifié en remplaçant $J$ par une puissance convenable.

\medskip\noindent
Supposons donné $J$ comme dans (2). Posons $N = H^0_T(M)$ et $M' = M/N$. La construction de $R\Gamma_T$ donne $H^i_T(N) = 0$ pour $i > 0$ et $H^0_T(N) = N$ ; voir la remarque \ref{remark-upshot}. On obtient donc $H^0_T(M') = 0$ et $H^i_T(M') = H^i_T(M)$ pour $i > 0$. On peut par conséquent remplacer $M$ par $M'$ et supposer ainsi que $H^0_T(M) = 0$. Cela signifie qu'aucun des idéaux premiers associés à $M$, qui sont en nombre fini, n'appartient à $T$. Par évitement des idéaux premiers (Algèbre, lemme \ref{algebra-lemma-silly}), on peut trouver $f \in J$ n'appartenant à aucun des idéaux premiers associés à $M$. Alors la suite exacte longue de cohomologie locale associée à la suite exacte courte
$$
0 \to M \to M \to M/fM \to 0
$$
se décompose en suites exactes courtes
$$
0 \to H^i_T(M) \to H^i_T(M/fM) \to H^{i + 1}_T(M) \to 0
$$
pour $i < n$. On en déduit que $J^2$ annule $H^i_T(M/fM)$ pour $i < n$. Par l'hypothèse de récurrence, $H^i_T(M/fM)$ est de type fini pour $i < n$. En utilisant une nouvelle fois la suite exacte courte, on voit que $H^{i + 1}_T(M)$ est de type fini pour $i < n$, comme voulu.

\medskip\noindent
Nous omettons la démonstration de l'équivalence de (2) et (3) dans le cas $T = V(I)$.
\end{proof}

\noindent
Le résultat suivant de Faltings permet de prouver la finitude de la cohomologie locale au niveau des anneaux locaux.

\begin{lemma}
\label{lemma-check-finiteness-local-cohomology-locally}
\begin{reference}
C'est un cas particulier de \cite[Satz 1]{Faltings-finiteness}.
\end{reference}
Soient $A$ un anneau noethérien, $I \subset A$ un idéal, $M$ un $A$-module de type fini et $n \geq 0$ un entier. Posons $Z = V(I)$. Les conditions suivantes sont équivalentes :
\begin{enumerate}
\item les modules $H^i_Z(M)$ sont de type fini pour $i \leq n$,
\item pour tout $\mathfrak p \in \Spec(A)$, les modules $H^i_Z(M)_\mathfrak p$, $i \leq n$ sont des $A_\mathfrak p$-modules de type fini.
\end{enumerate}
\end{lemma}

\begin{proof}
L'implication (1) $\Rightarrow$ (2) est immédiate. Démontrons la réciproque par récurrence sur $n$. Le cas $n = 0$ est clair, car (1) et (2) sont toujours vrais dans ce cas.

\medskip\noindent
Supposons $n > 0$ et (2) vrai. Posons $N = H^0_Z(M)$ et $M' = M/N$. D'après Complexes dualisants, lemme \ref{dualizing-lemma-divide-by-torsion}, on peut remplacer $M$ par $M'$. On peut donc supposer $H^0_Z(M) = 0$. Cela signifie que $\text{depth}_I(M) > 0$ (Complexes dualisants, lemme \ref{dualizing-lemma-depth}). Choisissons $f \in I$ non diviseur de zéro sur $M$ et considérons la suite exacte courte
$$
0 \to M \to M \to M/fM \to 0
$$
qui fournit une suite exacte longue
$$
0 \to H^0_Z(M/fM) \to H^1_Z(M) \to H^1_Z(M) \to H^1_Z(M/fM) \to
H^2_Z(M) \to \ldots
$$
et de même après localisation. L'hypothèse (2) implique donc que les modules $H^i_Z(M/fM)_\mathfrak p$ sont de type fini pour $i < n$. Par l'hypothèse de récurrence, les $H^i_Z(M/fM)$ sont alors de type fini pour $i < n$.

\medskip\noindent
Soit $\mathfrak p$ un idéal premier de $A$ associé à $H^i_Z(M)$ pour un certain $i \leq n$. Écrivons $\mathfrak p$ comme l'annulateur d'un élément $x \in H^i_Z(M)$. Alors $\mathfrak p \in Z$, donc $f \in \mathfrak p$. Ainsi $fx = 0$, de sorte que $x$ provient d'un élément de $H^{i - 1}_Z(M/fM)$ par l'homomorphisme de bord $\delta$ de la suite exacte longue ci-dessus. Par conséquent, $\mathfrak p$ est un idéal premier associé au module de type fini $\Im(\delta)$. On conclut que $\text{Ass}(H^i_Z(M))$ est fini pour $i \leq n$ ; voir Algèbre, lemme \ref{algebra-lemma-finite-ass}.

\medskip\noindent
Rappelons que
$$
H^i_Z(M) \subset
\prod\nolimits_{\mathfrak p \in \text{Ass}(H^i_Z(M))}
H^i_Z(M)_\mathfrak p
$$
d'après Algèbre, lemme \ref{algebra-lemma-zero-at-ass-zero}. Puisque, par hypothèse, les modules du membre de droite sont de type fini et de torsion par rapport aux puissances de $I$, on peut trouver des entiers $e_{\mathfrak p, i} \geq 0$, $i \leq n$, $\mathfrak p \in \text{Ass}(H^i_Z(M))$ tels que $I^{e_{\mathfrak p, i}}$ annule $H^i_Z(M)_\mathfrak p$. On en déduit que $I^e$, avec $e = \max\{e_{\mathfrak p, i}\}$, annule $H^i_Z(M)$ pour $i \leq n$. Le lemme \ref{lemma-check-finiteness-local-cohomology-by-annihilator} montre alors que $H^i_Z(M)$ est de type fini pour $i \leq n$.
\end{proof}

\begin{lemma}
\label{lemma-annihilate-local-cohomology}
Soient $A$ un anneau et $J \subset I \subset A$ des idéaux de type fini. Soit $i \geq 0$ un entier. Posons $Z = V(I)$. Si $H^i_Z(A)$ est annulé par $J^n$ pour un certain $n$, alors $H^i_Z(M)$ est annulé par $J^m$ pour un certain $m = m(M)$, pour tout $A$-module de présentation finie $M$ tel que $M_f$ soit un $A_f$-module localement libre de type fini pour tout $f \in I$.
\end{lemma}

\begin{proof}
Considérons l'annulateur $\mathfrak a$ de $H^i_Z(M)$. Soit $\mathfrak p \subset A$ tel que $\mathfrak p \not \in Z$. Par hypothèse, il existe $f \in I$, $f \not \in \mathfrak p$ et un isomorphisme $\varphi : A_f^{\oplus r} \to M_f$ de $A_f$-modules. En chassant les dénominateurs (et en utilisant le fait que $M$ est de présentation finie), on trouve des applications
$$
a : A^{\oplus r} \longrightarrow M
\quad\text{et}\quad
b : M \longrightarrow A^{\oplus r}
$$
avec $a_f = f^N \varphi$ et $b_f = f^N \varphi^{-1}$ pour un certain $N$. De plus, on peut supposer que $a \circ b$ et $b \circ a$ sont égales à la multiplication par $f^{2N}$. On voit alors que $H^i_Z(M)$ est annulé par $f^{2N}J^n$, c'est-à-dire que $f^{2N}J^n \subset \mathfrak a$.

\medskip\noindent
Puisque $U = \Spec(A) \setminus Z$ est quasi-compact, on peut trouver un nombre fini de $f_1, \ldots, f_t$ et de $N_1, \ldots, N_t$ tels que $U = \bigcup D(f_j)$ et $f_j^{2N_j}J^n \subset \mathfrak a$. Alors $V(I) = V(f_1, \ldots, f_t)$ et, puisque $I$ est de type fini, on en déduit $I^M \subset (f_1, \ldots, f_t)$ pour un certain $M$. Finalement, on voit que $J^m \subset \mathfrak a$ pour $m \gg 0$ ; par exemple, $m = M (2N_1 + \ldots + 2N_t) n$ convient.
\end{proof}

\begin{lemma}
\label{lemma-local-finiteness-for-finite-locally-free}
Soit $A$ un anneau noethérien. Soit $I \subset A$ un idéal. Posons $Z = V(I)$. Soit $n \geq 0$ un entier. Si $H^i_Z(A)$ est de type fini pour $0 \leq i \leq n$, il en est de même de $H^i_Z(M)$, $0 \leq i \leq n$, pour tout $A$-module de type fini $M$ tel que $M_f$ soit un $A_f$-module localement libre de type fini pour tout $f \in I$.
\end{lemma}

\begin{proof}
L'hypothèse selon laquelle $H^i_Z(A)$ est de type fini pour $0 \leq i \leq n$ implique l'existence de $e \geq 0$ tel que $I^e$ annule $H^i_Z(A)$ pour $0 \leq i \leq n$ ; voir le lemme \ref{lemma-check-finiteness-local-cohomology-by-annihilator}. Le lemme \ref{lemma-annihilate-local-cohomology} implique alors que $H^i_Z(M)$, $0 \leq i \leq n$ est annulé par $I^m$ pour un certain $m = m(M, i)$. On peut prendre le même $m$ pour tous les $0 \leq i \leq n$. Le lemme \ref{lemma-check-finiteness-local-cohomology-by-annihilator} implique alors que $H^i_Z(M)$ est de type fini pour $0 \leq i \leq n$, comme voulu.
\end{proof}





\section{Finitude des images directes, I}
\label{section-finiteness-pushforward}

\noindent
Dans cette section, nous étudions le cas non trivial le plus simple du théorème de finitude, à savoir la finitude de la cohomologie locale de degré un, ou, de manière équivalente, celle de $j_*\mathcal{F}$, où $j : U \to X$ est une immersion ouverte, $X$ est localement noethérien et $\mathcal{F}$ est un faisceau cohérent sur $U$. En suivant une méthode de Koll\'ar (\cite{Kollar-variants} et \cite{Kollar-local-global-hulls}), nous obtenons une condition nécessaire et suffisante ; voir la proposition \ref{proposition-kollar}. Le lecteur qui s'intéresse aux images directes supérieures ou aux groupes de cohomologie locale de degré supérieur peut passer directement à la section \ref{section-finiteness-pushforward-II} ou à la section \ref{section-finiteness-II} (développées indépendamment du reste de la présente section).

\begin{lemma}
\label{lemma-check-finiteness-pushforward-on-associated-points}
Soit $X$ un schéma localement noethérien. Soit $j : U \to X$ l'inclusion d'un sous-schéma ouvert de complémentaire $Z$. Pour $x \in U$, soit $i_x : W_x \to U$ le sous-schéma fermé intègre de point générique $x$. Soit $\mathcal{F}$ un $\mathcal{O}_U$-module cohérent. Les conditions suivantes sont équivalentes :
\begin{enumerate}
\item pour tout $x \in \text{Ass}(\mathcal{F})$, le $\mathcal{O}_X$-module $j_*i_{x, *}\mathcal{O}_{W_x}$ est cohérent,
\item $j_*\mathcal{F}$ est cohérent.
\end{enumerate}
\end{lemma}

\begin{proof}
Montrons d'abord que (1) implique (2). Supposons (1) vérifié. L'énoncé est local sur $X$ ; on peut donc supposer $X$ affine. Alors $U$ est quasi-compact, donc $\text{Ass}(\mathcal{F})$ est fini (Diviseurs, lemme \ref{divisors-lemma-finite-ass}). On peut ainsi raisonner par récurrence sur le nombre de points associés. Soit $x \in U$ un point générique d'une composante irréductible du support de $\mathcal{F}$. D'après Diviseurs, lemme \ref{divisors-lemma-finite-ass}, on a $x \in \text{Ass}(\mathcal{F})$. Le choix de $x$ assure que $\dim(\mathcal{F}_x) = 0$ en tant que $\mathcal{O}_{X, x}$-module. Ainsi $\mathcal{F}_x$ est de longueur finie comme $\mathcal{O}_{X, x}$-module (Algèbre, lemme \ref{algebra-lemma-support-point}). On peut donc raisonner par récurrence sur cette longueur.

\medskip\noindent
Posons $\mathcal{G} = j_*i_{x, *}\mathcal{O}_{W_x}$. C'est un $\mathcal{O}_X$-module cohérent par hypothèse. On a $\mathcal{G}_x = \kappa(x)$. Choisissons une application non nulle $\varphi_x : \mathcal{F}_x \to \kappa(x) = \mathcal{G}_x$. D'après Cohomologie des schémas, lemme \ref{coherent-lemma-map-stalks-local-map}, il existe un ouvert $x \in V \subset U$ et une application $\varphi_V : \mathcal{F}|_V \to \mathcal{G}|_V$ dont le germe en $x$ est $\varphi_x$. Choisissons $f \in \Gamma(X, \mathcal{O}_X)$ ne s'annulant pas en $x$ et tel que $D(f) \subset V$. D'après Cohomologie des schémas, lemme \ref{coherent-lemma-homs-over-open} (par exemple), on voit que $\varphi_V$ se prolonge en $f^n\mathcal{F} \to \mathcal{G}|_U$ pour un certain $n$. En précomposant avec la multiplication par $f^n$, on obtient une application $\mathcal{F} \to \mathcal{G}|_U$ dont le germe en $x$ est non nul. Soit $\mathcal{F}' \subset \mathcal{F}$ son noyau. Notons que $\text{Ass}(\mathcal{F}') \subset \text{Ass}(\mathcal{F})$ ; voir Diviseurs, lemme \ref{divisors-lemma-ses-ass}. Puisque $\text{length}_{\mathcal{O}_{X, x}}(\mathcal{F}'_x) =
\text{length}_{\mathcal{O}_{X, x}}(\mathcal{F}_x) - 1$, l'hypothèse de récurrence permet de conclure que $j_*\mathcal{F}'$ est cohérent. Puisque $\mathcal{G} = j_*(\mathcal{G}|_U) = j_*i_{x, *}\mathcal{O}_{W_x}$ est cohérent, on peut considérer la suite exacte
$$
0 \to j_*\mathcal{F}' \to j_*\mathcal{F} \to \mathcal{G}
$$
D'après Schémas, lemme \ref{schemes-lemma-push-forward-quasi-coherent}, le faisceau $j_*\mathcal{F}$ est quasi-cohérent. L'image de $j_*\mathcal{F}$ dans $j_*(\mathcal{G}|_U)$ est donc cohérente, d'après Cohomologie des schémas, lemme \ref{coherent-lemma-coherent-Noetherian-quasi-coherent-sub-quotient}. Enfin, $j_*\mathcal{F}$ est cohérent d'après Cohomologie des schémas, lemme \ref{coherent-lemma-coherent-abelian-Noetherian}.

\medskip\noindent
Supposons (2) vérifié. Exactement comme ci-dessus, on se ramène au cas où $X$ est affine. Choisissons $x \in \text{Ass}(\mathcal{F})$ et posons $\mathcal{G} = j_*i_{x, *}\mathcal{O}_{W_x}$. Choisissons ensuite une application non nulle $\varphi_x : \mathcal{G}_x = \kappa(x) \to \mathcal{F}_x$, qui existe précisément parce que $x$ est un point associé de $\mathcal{F}$. En raisonnant exactement comme ci-dessus, on peut supposer que $\varphi_x$ se prolonge en un homomorphisme de $\mathcal{O}_U$-modules $\varphi : \mathcal{G}|_U \to \mathcal{F}$. Alors $\varphi$ est injectif (par exemple d'après Diviseurs, lemme \ref{divisors-lemma-check-injective-on-ass}), et l'on obtient une application injective $\mathcal{G} = j_*(\mathcal{G}|_V) \to j_*\mathcal{F}$. Ainsi (1) est vérifié.
\end{proof}

\begin{lemma}
\label{lemma-finiteness-pushforwards-and-H1-local}
Soient $A$ un anneau noethérien et $I \subset A$ un idéal. Posons $X = \Spec(A)$, $Z = V(I)$, $U = X \setminus Z$, et notons $j : U \to X$ le morphisme d'inclusion. Soit $\mathcal{F}$ un $\mathcal{O}_U$-module cohérent. Alors
\begin{enumerate}
\item il existe un $A$-module de type fini $M$ tel que $\mathcal{F}$ soit la restriction de $\widetilde{M}$ à $U$,
\item pour un tel $M$, on a une suite exacte
$$
0 \to H^0_Z(M) \to M \to H^0(U, \mathcal{F}) \to H^1_Z(M) \to 0
$$
et des isomorphismes $H^p(U, \mathcal{F}) = H^{p + 1}_Z(M)$ pour $p \geq 1$,
\item pour $M$ et $p \geq 0$ donnés, les conditions suivantes sont équivalentes :
\begin{enumerate}
\item $R^pj_*\mathcal{F}$ est cohérent,
\item $H^p(U, \mathcal{F})$ est un $A$-module de type fini,
\item $H^{p + 1}_Z(M)$ est un $A$-module de type fini,
\end{enumerate}
\item si les conditions équivalentes de (3) sont vérifiées pour $p = 0$, on peut prendre $M = \Gamma(U, \mathcal{F})$, auquel cas on a $H^0_Z(M) = H^1_Z(M) = 0$.
\end{enumerate}
\end{lemma}

\begin{proof}
D'après Propriétés, lemme \ref{properties-lemma-lift-finite-presentation}, il existe un $\mathcal{O}_X$-module cohérent $\mathcal{F}'$ dont la restriction à $U$ est isomorphe à $\mathcal{F}$. Supposons que $\mathcal{F}'$ corresponde au $A$-module de type fini $M$ comme en (1). Notons que $R^pj_*\mathcal{F}$ est quasi-cohérent (Cohomologie des schémas, lemme \ref{coherent-lemma-quasi-coherence-higher-direct-images}) et correspond au $A$-module $H^p(U, \mathcal{F})$. D'après le lemme \ref{lemma-local-cohomology-is-local-cohomology} et la discussion de Cohomologie, sections \ref{cohomology-section-cohomology-support} et \ref{cohomology-section-cohomology-support-bis}, on obtient une suite exacte
$$
0 \to H^0_Z(M) \to M \to H^0(U, \mathcal{F}) \to H^1_Z(M) \to 0
$$
et des isomorphismes $H^p(U, \mathcal{F}) = H^{p + 1}_Z(M)$ pour $p \geq 1$. On utilise ici le fait que $H^j(X, \mathcal{F}') = 0$ pour $j > 0$, puisque $X$ est affine et $\mathcal{F}'$ est quasi-cohérent (Cohomologie des schémas, lemme \ref{coherent-lemma-quasi-coherent-affine-cohomology-zero}). Cela prouve (2). Les points (3) et (4) résultent directement de (2) ; voir aussi le lemme \ref{lemma-local-cohomology}.
\end{proof}

\begin{lemma}
\label{lemma-finiteness-pushforward}
Soit $X$ un schéma localement noethérien. Soit $j : U \to X$ l'inclusion d'un sous-schéma ouvert de complémentaire $Z$. Soit $\mathcal{F}$ un $\mathcal{O}_U$-module cohérent. Supposons que
\begin{enumerate}
\item $X$ soit de Nagata,
\item $X$ soit universellement caténaire,
\item pour $x \in \text{Ass}(\mathcal{F})$ et $z \in Z \cap \overline{\{x\}}$, on ait $\dim(\mathcal{O}_{\overline{\{x\}}, z}) \geq 2$.
\end{enumerate}
Alors $j_*\mathcal{F}$ est cohérent.
\end{lemma}

\begin{proof}
D'après le lemme \ref{lemma-check-finiteness-pushforward-on-associated-points}, il suffit de montrer que $j_*i_{x, *}\mathcal{O}_{W_x}$ est cohérent pour $x \in \text{Ass}(\mathcal{F})$. Soit $\pi : Y \to X$ la normalisation de $X$ dans $\Spec(\kappa(x))$ ; voir Morphismes, section \ref{morphisms-section-normalization}. D'après Morphismes, lemme \ref{morphisms-lemma-nagata-normalization-finite-general}, le morphisme $\pi$ est fini. Puisque $\pi$ est fini, $\mathcal{G} = \pi_*\mathcal{O}_Y$ est un $\mathcal{O}_X$-module cohérent d'après Cohomologie des schémas, lemme \ref{coherent-lemma-finite-pushforward-coherent}. Observons que $W_x = U \cap \pi(Y)$. Ainsi $\pi|_{\pi^{-1}(U)} : \pi^{-1}(U) \to U$ se factorise par $i_x : W_x \to U$, et l'on obtient une application canonique
$$
i_{x, *}\mathcal{O}_{W_x}
\longrightarrow
(\pi|_{\pi^{-1}(U)})_*(\mathcal{O}_{\pi^{-1}(U)}) =
(\pi_*\mathcal{O}_Y)|_U = \mathcal{G}|_U
$$
Cette application est injective (par exemple d'après Diviseurs, lemme \ref{divisors-lemma-check-injective-on-ass}). On a donc $j_*i_{x, *}\mathcal{O}_{W_x} \subset j_*\mathcal{G}|_U$, et il suffit de montrer que $j_*\mathcal{G}|_U$ est cohérent.

\medskip\noindent
Il reste à montrer que $j_*(\mathcal{G}|_U)$ est cohérent. Montrons que Diviseurs, lemme \ref{divisors-lemma-depth-2-hartog}, s'applique à
$$
\mathcal{G} \longrightarrow j_*(\mathcal{G}|_U)
$$
ce qui achève la démonstration. Il suffit de montrer que $\text{depth}(\mathcal{G}_z) \geq 2$ pour $z \in Z$. Soient $y_1, \ldots, y_n \in Y$ les points s'envoyant sur $z$. D'après Algèbre, lemme \ref{algebra-lemma-depth-goes-down-finite}, il suffit de montrer que $\text{depth}(\mathcal{O}_{Y, y_i}) \geq 2$ pour $i = 1, \ldots, n$. Sinon, Propriétés, lemme \ref{properties-lemma-criterion-normal}, montre que $\dim(\mathcal{O}_{Y, y_i}) = 1$ pour un certain $i$. C'est impossible d'après la formule des dimensions (Morphismes, lemme \ref{morphisms-lemma-dimension-formula}) pour $\pi : Y \to \overline{\{x\}}$ et l'hypothèse (3).
\end{proof}

\begin{lemma}
\label{lemma-sharp-finiteness-pushforward}
Soit $X$ un schéma intègre localement noethérien. Soit $j : U \to X$ l'inclusion d'un sous-schéma ouvert non vide de complémentaire $Z$. Supposons que, pour tout $z \in Z$ et tout idéal premier associé $\mathfrak p$ du complété $\mathcal{O}_{X, z}^\wedge$, on ait $\dim(\mathcal{O}_{X, z}^\wedge/\mathfrak p) \geq 2$. Alors $j_*\mathcal{O}_U$ est cohérent.
\end{lemma}

\begin{proof}
On peut supposer $X$ affine. Les lemmes \ref{lemma-check-finiteness-local-cohomology-locally} et \ref{lemma-finiteness-pushforwards-and-H1-local} permettent de se ramener à $X = \Spec(A)$, où $(A, \mathfrak m)$ est un anneau local noethérien intègre et $\mathfrak m \in Z$. On peut alors raisonner par récurrence sur $d = \dim(A)$. (Les cas initiaux $d = 0, 1$ ne se présentent pas en vertu de l'hypothèse sur les anneaux locaux.) Posons $V = \Spec(A) \setminus \{\mathfrak m\}$. Observons que les anneaux locaux de $V$ ont une dimension strictement inférieure à $d$. En reprenant les arguments pour $j' : U \to V$ et en utilisant la récurrence, on conclut que $j'_*\mathcal{O}_U$ est un $\mathcal{O}_V$-module cohérent. Choisissons un élément non nul $f \in A$ qui s'annule sur $Z$. Puisque $D(f) \cap V \subset U$, on trouve $n$ tel que la multiplication par $f^n$ sur $U$ se prolonge en une application $f^n : j'_*\mathcal{O}_U \to \mathcal{O}_V$ sur $V$ (par exemple d'après Cohomologie des schémas, lemme \ref{coherent-lemma-homs-over-open}). Cette application est injective ; on a donc une application injective
$$
j_*\mathcal{O}_U = j''_* j'_* \mathcal{O}_U \to j''_*\mathcal{O}_V
$$
sur $X$, où $j'' : V \to X$ est le morphisme d'inclusion. Il suffit donc de montrer que $j''_*\mathcal{O}_V$ est cohérent. Autrement dit, on peut supposer que $X$ est le spectre d'un anneau local noethérien intègre et que $Z$ est constitué du point fermé.

\medskip\noindent
Supposons $X = \Spec(A)$, avec $(A, \mathfrak m)$ local et $Z = \{\mathfrak m\}$. Soit $A^\wedge$ le complété de $A$. Posons $X^\wedge = \Spec(A^\wedge)$, $Z^\wedge = \{\mathfrak m^\wedge\}$, $U^\wedge = X^\wedge \setminus Z^\wedge$ et $\mathcal{F}^\wedge = \mathcal{O}_{U^\wedge}$. L'anneau $A^\wedge$ est universellement caténaire et de Nagata (Algèbre, remarque \ref{algebra-remark-Noetherian-complete-local-ring-universally-catenary} et lemme \ref{algebra-lemma-Noetherian-complete-local-Nagata}). De plus, la condition (3) du lemme \ref{lemma-finiteness-pushforward} pour $X^\wedge, Z^\wedge, U^\wedge, \mathcal{F}^\wedge$ est vérifiée par hypothèse ! On voit donc que $(U^\wedge \to X^\wedge)_*\mathcal{O}_{U^\wedge}$ est cohérent. Puisque le morphisme $c : X^\wedge \to X$ est plat, l'image réciproque de $j_*\mathcal{O}_U$ est $(U^\wedge \to X^\wedge)_*\mathcal{O}_{U^\wedge}$ (Cohomologie des schémas, lemme \ref{coherent-lemma-flat-base-change-cohomology}). Enfin, puisque $c$ est fidèlement plat, on conclut que $j_*\mathcal{O}_U$ est cohérent d'après Descente, lemme \ref{descent-lemma-finite-type-descends}.
\end{proof}

\begin{remark}
\label{remark-closure}
Soit $j : U \to X$ une immersion ouverte de schémas localement noethériens. Soit $x \in U$. Soit $i_x : W_x \to U$ le sous-schéma fermé intègre de point générique $x$, et soit $\overline{\{x\}}$ l'adhérence dans $X$. On a alors un diagramme commutatif
$$
\xymatrix{
W_x \ar[d]_{i_x} \ar[r]_{j'} & \overline{\{x\}} \ar[d]^i \\
U \ar[r]^j & X
}
$$
On a $j_*i_{x, *}\mathcal{O}_{W_x} = i_*j'_*\mathcal{O}_{W_x}$. Puisque la flèche verticale de gauche est une immersion fermée, $j_*i_{x, *}\mathcal{O}_{W_x}$ est cohérent si et seulement si $j'_*\mathcal{O}_{W_x}$ est cohérent.
\end{remark}

\begin{remark}
\label{remark-no-finiteness-pushforward}
Soit $X$ un schéma localement noethérien. Soit $j : U \to X$ l'inclusion d'un sous-schéma ouvert de complémentaire $Z$. Soit $\mathcal{F}$ un $\mathcal{O}_U$-module cohérent. S'il existe $x \in \text{Ass}(\mathcal{F})$ et $z \in Z \cap \overline{\{x\}}$ tels que $\dim(\mathcal{O}_{\overline{\{x\}}, z}) \leq 1$, alors $j_*\mathcal{F}$ n'est pas cohérent. Pour le montrer, on peut effectuer un changement de base plat vers le spectre de $\mathcal{O}_{X, z}$. Posons $X' = \overline{\{x\}}$. L'hypothèse implique $\mathcal{O}_{X' \cap U} \subset \mathcal{F}$. Il suffit donc de voir que $j_*\mathcal{O}_{X' \cap U}$ n'est pas cohérent. Cela est clair, car $X' = \{x, z\}$ ; ainsi $j_*\mathcal{O}_{X' \cap U}$ correspond à $\kappa(x)$ en tant que $\mathcal{O}_{X, z}$-module, qui ne peut être de type fini puisque $x$ n'est pas un point fermé.

\medskip\noindent
En fait, la réciproque du lemme \ref{lemma-sharp-finiteness-pushforward} est vraie : étant donnée une immersion ouverte $j : U \to X$ de schémas noethériens intègres, s'il existe $z \in X \setminus U$ et un idéal premier associé $\mathfrak p$ du complété $\mathcal{O}_{X, z}^\wedge$ tel que $\dim(\mathcal{O}_{X, z}^\wedge/\mathfrak p) = 1$, alors $j_*\mathcal{O}_U$ n'est pas cohérent. En effet, on peut passer à l'anneau local, agrandir $U$ au spectre épointé, passer au complété, puis appliquer l'argument ci-dessus pour obtenir la non-finitude.
\end{remark}

\begin{proposition}[Koll\'ar]
\label{proposition-kollar}
\begin{reference}
Voir \cite{k-coherent}, et voir \cite[IV, Proposition 7.2.2]{EGA} pour un cas particulier.
\end{reference}
\begin{slogan}
Analogue faible du théorème de Hartogs : sur un schéma noethérien, la restriction d'un faisceau cohérent à un ouvert dont le complémentaire est de codimension 2 dans le support du faisceau est cohérente.
\end{slogan}
Soit $j : U \to X$ une immersion ouverte de schémas localement noethériens de complémentaire $Z$. Soit $\mathcal{F}$ un $\mathcal{O}_U$-module cohérent. Les conditions suivantes sont équivalentes :
\begin{enumerate}
\item $j_*\mathcal{F}$ est cohérent,
\item pour $x \in \text{Ass}(\mathcal{F})$, $z \in Z \cap \overline{\{x\}}$ et tout idéal premier associé $\mathfrak p$ du complété $\mathcal{O}_{\overline{\{x\}}, z}^\wedge$, on a $\dim(\mathcal{O}_{\overline{\{x\}}, z}^\wedge/\mathfrak p) \geq 2$.
\end{enumerate}
\end{proposition}

\begin{proof}
Si (2) est vérifié, on obtient (1) en combinant le lemme \ref{lemma-check-finiteness-pushforward-on-associated-points}, la remarque \ref{remark-closure} et le lemme \ref{lemma-sharp-finiteness-pushforward}. Si (2) n'est pas vérifié, alors $j_*i_{x, *}\mathcal{O}_{W_x}$ n'est pas de type fini pour un certain $x \in \text{Ass}(\mathcal{F})$, d'après la discussion de la remarque \ref{remark-no-finiteness-pushforward} (et la remarque \ref{remark-closure}). Ainsi $j_*\mathcal{F}$ n'est pas cohérent d'après le lemme \ref{lemma-check-finiteness-pushforward-on-associated-points}.
\end{proof}

\begin{lemma}
\label{lemma-kollar-finiteness-H1-local}
Soient $A$ un anneau noethérien et $I \subset A$ un idéal. Posons $Z = V(I)$. Soit $M$ un $A$-module de type fini. Les conditions suivantes sont équivalentes :
\begin{enumerate}
\item $H^1_Z(M)$ est un $A$-module de type fini,
\item pour tous $\mathfrak p \in \text{Ass}(M)$, $\mathfrak p \not \in Z$ et tout $\mathfrak q \in V(\mathfrak p + I)$, le complété de $(A/\mathfrak p)_\mathfrak q$ n'a pas d'idéal premier associé de dimension $1$.
\end{enumerate}
\end{lemma}

\begin{proof}
Cela résulte immédiatement de la proposition \ref{proposition-kollar}, par l'intermédiaire du lemme \ref{lemma-finiteness-pushforwards-and-H1-local}.
\end{proof}

\noindent
La formulation du lemme suivant a l'avantage que les conditions (1) et (2) se transmettent aux schémas de type fini sur $X$. De plus, c'est cette forme de finitude que nous généraliserons aux images directes supérieures dans la section \ref{section-finiteness-pushforward-II}.

\begin{lemma}
\label{lemma-finiteness-pushforward-general}
Soit $X$ un schéma localement noethérien. Soit $j : U \to X$ l'inclusion d'un sous-schéma ouvert de complémentaire $Z$. Soit $\mathcal{F}$ un $\mathcal{O}_U$-module cohérent. Supposons que
\begin{enumerate}
\item $X$ soit universellement caténaire,
\item pour tout $z \in Z$, les fibres formelles de $\mathcal{O}_{X, z}$ vérifient $(S_1)$.
\end{enumerate}
Dans cette situation, les conditions suivantes sont équivalentes :
\begin{enumerate}
\item[(a)] pour $x \in \text{Ass}(\mathcal{F})$ et $z \in Z \cap \overline{\{x\}}$, on a $\dim(\mathcal{O}_{\overline{\{x\}}, z}) \geq 2$,
\item[(b)] $j_*\mathcal{F}$ est cohérent.
\end{enumerate}
\end{lemma}

\begin{proof}
Soit $x \in \text{Ass}(\mathcal{F})$. D'après la proposition \ref{proposition-kollar}, il suffit de vérifier que $A = \mathcal{O}_{\overline{\{x\}}, z}$ satisfait à la condition de la proposition concernant les idéaux premiers associés de son complété si et seulement si $\dim(A) \geq 2$. Observons que $A$ est universellement caténaire (ce qui est clair) et que ses fibres formelles vérifient $(S_1)$, d'après Compléments d'algèbre, lemme \ref{more-algebra-lemma-formal-fibres-normal} et proposition \ref{more-algebra-proposition-finite-type-over-P-ring}. Soit $\mathfrak p' \subset A^\wedge$ un idéal premier associé. Puisque $A \to A^\wedge$ est plat, Algèbre, lemme \ref{algebra-lemma-bourbaki}, montre que $\mathfrak p'$ est au-dessus de $(0) \subset A$. La fibre formelle $A^\wedge \otimes_A F$ vérifie $(S_1)$, où $F$ est le corps des fractions de $A$. On en déduit que $\mathfrak p'$ est un idéal premier minimal ; voir Algèbre, lemme \ref{algebra-lemma-criterion-no-embedded-primes}. Puisque $A$ est universellement caténaire, il est formellement caténaire d'après Compléments d'algèbre, proposition \ref{more-algebra-proposition-ratliff}. Ainsi $\dim(A^\wedge/\mathfrak p') = \dim(A)$, ce qui prouve l'équivalence.
\end{proof}






\section{Profondeur et dimension}
\label{section-dept-dimension}

\noindent
Quelques lemmes auxiliaires.

\begin{lemma}
\label{lemma-ideal-depth-function}
Soit $A$ un anneau noethérien. Soit $I \subset A$ un idéal. Soit $M$ un $A$-module de type fini. Soit $\mathfrak p \in V(I)$ un idéal premier. Supposons $e = \text{depth}_{IA_\mathfrak p}(M_\mathfrak p) < \infty$. Alors il existe un ouvert non vide $U \subset V(\mathfrak p)$ tel que $\text{depth}_{IA_\mathfrak q}(M_\mathfrak q) \geq e$ pour tout $\mathfrak q \in U$.
\end{lemma}

\begin{proof}
Par définition de la profondeur, on a $IM_\mathfrak p \not = M_\mathfrak p$ et il existe une suite $M_\mathfrak p$-régulière $f_1, \ldots, f_e \in IA_\mathfrak p$. Après avoir remplacé $A$ par un localisé principal, on peut supposer que $f_1, \ldots, f_e \in I$ forment une suite $M$-régulière ; voir Algèbre, lemme \ref{algebra-lemma-regular-sequence-in-neighbourhood}. Considérons le module $M' = M/IM$. Puisque $\mathfrak p \in \text{Supp}(M')$ et que le support d'un module de type fini est fermé, on trouve $V(\mathfrak p) \subset \text{Supp}(M')$. Ainsi, pour $\mathfrak q \in V(\mathfrak p)$, on a $IM_\mathfrak q \not = M_\mathfrak q$. Par conséquent, l'exactitude de la localisation et la définition de la profondeur donnent $\text{depth}_{IA_\mathfrak q}(M_\mathfrak q) \geq e$ pour tout $\mathfrak q \in V(I)$.
\end{proof}

\begin{lemma}
\label{lemma-depth-function}
Soit $A$ un anneau noethérien. Soit $M$ un $A$-module de type fini. Soit $\mathfrak p$ un idéal premier. Supposons $e = \text{depth}_{A_\mathfrak p}(M_\mathfrak p) < \infty$. Alors il existe un ouvert non vide $U \subset V(\mathfrak p)$ tel que $\text{depth}_{A_\mathfrak q}(M_\mathfrak q) \geq e$ pour tout $\mathfrak q \in U$ et tel que, pour tous les $\mathfrak q \in U$ sauf un nombre fini, on ait $\text{depth}_{A_\mathfrak q}(M_\mathfrak q) > e$.
\end{lemma}

\begin{proof}
Par définition de la profondeur, on a $\mathfrak p M_\mathfrak p \not = M_\mathfrak p$ et il existe une suite $M_\mathfrak p$-régulière $f_1, \ldots, f_e \in \mathfrak p A_\mathfrak p$. Après avoir remplacé $A$ par un localisé principal, on peut supposer que $f_1, \ldots, f_e \in \mathfrak p$ forment une suite $M$-régulière ; voir Algèbre, lemme \ref{algebra-lemma-regular-sequence-in-neighbourhood}. Considérons le module $M' = M/(f_1, \ldots, f_e)M$. Puisque $\mathfrak p \in \text{Supp}(M')$ et que le support d'un module de type fini est fermé, on trouve $V(\mathfrak p) \subset \text{Supp}(M')$. Ainsi, pour $\mathfrak q \in V(\mathfrak p)$, on a $\mathfrak q M_\mathfrak q \not = M_\mathfrak q$. Par conséquent, l'exactitude de la localisation et la définition de la profondeur donnent $\text{depth}_{A_\mathfrak q}(M_\mathfrak q) \geq e$ pour tout $\mathfrak q \in V(I)$. De plus, dès que $\mathfrak q$ n'est pas un idéal premier associé au module $M'$, la profondeur augmente. La dernière assertion résulte donc d'Algèbre, lemme \ref{algebra-lemma-finite-ass}.
\end{proof}

\begin{lemma}
\label{lemma-finite-nr-points-next-S}
Soit $X$ un schéma noethérien muni d'un complexe dualisant $\omega_X^\bullet$. Soit $\mathcal{F}$ un $\mathcal{O}_X$-module cohérent. Soit $k \geq 0$ un entier. Supposons que $\mathcal{F}$ vérifie $(S_k)$. Alors les points $x \in X$ tels que
$$
\text{depth}(\mathcal{F}_x) = k
\quad\text{et}\quad
\dim(\text{Supp}(\mathcal{F}_x)) > k
$$
sont en nombre fini.
\end{lemma}

\begin{proof}
Démontrons ce lemme par récurrence sur $k$. Le cas initial $k = 0$ affirme que $\mathcal{F}$ possède un nombre fini de points associés immergés, ce qui résulte de Diviseurs, lemme \ref{divisors-lemma-finite-ass}.

\medskip\noindent
Supposons $k > 0$ et le résultat établi pour toute valeur plus petite de $k$. On peut recouvrir $X$ par un nombre fini d'ouverts affines ; on peut donc supposer $X = \Spec(A)$ affine. Alors $\mathcal{F}$ est le $\mathcal{O}_X$-module cohérent associé à un $A$-module de type fini $M$ qui vérifie $(S_k)$. Nous utiliserons Algèbre, lemmes \ref{algebra-lemma-one-equation-module} et \ref{algebra-lemma-depth-drops-by-one}, sans le mentionner à nouveau.

\medskip\noindent
Soit $f \in A$ un non-diviseur de zéro sur $M$. Alors $M/fM$ vérifie $(S_{k - 1})$. Par récurrence, les idéaux premiers $\mathfrak p \in V(f)$ tels que $\text{depth}((M/fM)_\mathfrak p) = k - 1$ et $\dim(\text{Supp}((M/fM)_\mathfrak p)) > k - 1$ sont en nombre fini. Ce sont exactement les idéaux premiers $\mathfrak p \in V(f)$ tels que $\text{depth}(M_\mathfrak p) = k$ et $\dim(\text{Supp}(M_\mathfrak p)) > k$. On peut donc remplacer $A$ par $A_f$ et $M$ par $M_f$ pour démontrer l'assertion de finitude.

\medskip\noindent
Puisque $M$ vérifie $(S_k)$ et que $k > 0$, le module $M$ n'a pas d'idéal premier associé immergé (Algèbre, lemme \ref{algebra-lemma-criterion-no-embedded-primes}). Ainsi $\text{Ass}(M)$ est l'ensemble des points génériques du support de $M$. Complexes dualisants, lemme \ref{dualizing-lemma-CM-open}, montre donc que l'ensemble $U = \{\mathfrak q \mid M_\mathfrak q\text{ est de Cohen-Macaulay}\}$ est un ouvert contenant $\text{Ass}(M)$. Par évitement des idéaux premiers (Algèbre, lemme \ref{algebra-lemma-silly}), on peut choisir $f \in A$ tel que $f \not \in \mathfrak p$ pour $\mathfrak p \in \text{Ass}(M)$ et tel que $D(f) \subset U$. Alors $f$ est un non-diviseur de zéro sur $M$ (Algèbre, lemme \ref{algebra-lemma-ass-zero-divisors}). Après avoir remplacé $A$ par $A_f$ et $M$ par $M_f$ (voir ci-dessus), on obtient que $M$ est de Cohen-Macaulay. Ainsi, pour tout $\mathfrak q \subset A$, on a $\dim(M_\mathfrak q) = \text{depth}(M_\mathfrak q)$ ; l'ensemble décrit dans le lemme est donc vide, et a fortiori fini.
\end{proof}

\begin{lemma}
\label{lemma-sitting-in-degrees}
Soit $(A, \mathfrak m)$ un anneau local noethérien muni d'un complexe dualisant normalisé $\omega_A^\bullet$. Soit $M$ un $A$-module de type fini. Posons $E^i = \text{Ext}_A^{-i}(M, \omega_A^\bullet)$. Alors
\begin{enumerate}
\item {\emergencystretch=2em $E^i$ est un $A$-module de type fini, non nul seulement pour $0 \leq i \leq \dim(\text{Supp}(M))$,\par}
\item $\dim(\text{Supp}(E^i)) \leq i$,
\item $\text{depth}(M)$ est le plus petit entier $\delta \geq 0$ tel que $E^\delta \not = 0$,
\item $\mathfrak p \in \text{Supp}(E^0 \oplus \ldots \oplus E^i)
\Leftrightarrow
\text{depth}_{A_\mathfrak p}(M_\mathfrak p) + \dim(A/\mathfrak p) \leq i$,
\item l'annulateur de $E^i$ est égal à celui de $H^i_\mathfrak m(M)$.
\end{enumerate}
\end{lemma}

\begin{proof}
{\emergencystretch=3em
Les points (1), (2) et (3) reproduisent les énoncés de Complexes dualisants, lemme \ref{dualizing-lemma-sitting-in-degrees}. Pour un idéal premier $\mathfrak p$ de $A$, le complexe $(\omega_A^\bullet)_\mathfrak p[-\dim(A/\mathfrak p)]$ est un complexe dualisant normalisé pour $A_\mathfrak p$. Voir Complexes dualisants, lemme \ref{dualizing-lemma-dimension-function}. Ainsi
\par}
$$
E^i_\mathfrak p =
\text{Ext}^{-i}_A(M, \omega_A^\bullet)_\mathfrak p =
\text{Ext}^{-i + \dim(A/\mathfrak p)}_{A_\mathfrak p}
(M_\mathfrak p, (\omega_A^\bullet)_\mathfrak p[-\dim(A/\mathfrak p)])
$$
est nul pour $i - \dim(A/\mathfrak p) < \text{depth}_{A_\mathfrak p}(M_\mathfrak p)$ et non nul pour $i = \dim(A/\mathfrak p) + \text{depth}_{A_\mathfrak p}(M_\mathfrak p)$, d'après (3) appliqué sur $A_\mathfrak p$. Cela prouve (4). Si $E$ est une enveloppe injective du corps résiduel de $A$, alors on a
$$
\Hom_A(H^i_\mathfrak m(M), E) =
\text{Ext}^{-i}_A(M, \omega_A^\bullet)^\wedge =
(E^i)^\wedge =
E^i \otimes_A A^\wedge
$$
d'après le théorème de dualité locale (sous la forme de Complexes dualisants, lemme \ref{dualizing-lemma-special-case-local-duality}). Puisque $A \to A^\wedge$ est fidèlement plat, la dualité de Matlis (Complexes dualisants, proposition \ref{dualizing-proposition-matlis}) prouve (5).
\end{proof}






\section{Annulateurs de la cohomologie locale, I}
\label{section-annihilators}

\noindent
Cette section présente un résultat dû à Faltings ; voir \cite{Faltings-annulators}.

\begin{proposition}
\label{proposition-annihilator}
\begin{reference}
\cite{Faltings-annulators}.
\end{reference}
Soit $A$ un anneau noethérien possédant un complexe dualisant. Soient $T \subset T' \subset \Spec(A)$ des parties stables par spécialisation. Soit $s \geq 0$ un entier. Soit $M$ un $A$-module de type fini. Les conditions suivantes sont équivalentes :
\begin{enumerate}
\item il existe un idéal $J \subset A$ tel que $V(J) \subset T'$ et tel que $J$ annule $H^i_T(M)$ pour $i \leq s$,
\item pour tous $\mathfrak p \not \in T'$, $\mathfrak q \in T$ tels que $\mathfrak p \subset \mathfrak q$, on a
$$
\text{depth}_{A_\mathfrak p}(M_\mathfrak p) +
\dim((A/\mathfrak p)_\mathfrak q) > s
$$
\end{enumerate}
\end{proposition}

\begin{proof}
Soit $\omega_A^\bullet$ un complexe dualisant. Soit $\delta$ sa fonction de dimension ; voir Complexes dualisants, section \ref{dualizing-section-dimension-function}. Les $A$-modules de type fini
$$
E^i = \Ext_A^i(M, \omega_A^\bullet)
$$
joueront un rôle important. Pour $\mathfrak p \subset A$, nous écrirons $H^i_\mathfrak p$ pour désigner la cohomologie locale d'un $A_\mathfrak p$-module relativement à $\mathfrak pA_\mathfrak p$. On voit alors que le complété $\mathfrak pA_\mathfrak p$-adique de
$$
(E^i)_\mathfrak p =
\Ext^{\delta(\mathfrak p) + i}_{A_\mathfrak p}(M_\mathfrak p,
(\omega_A^\bullet)_\mathfrak p[-\delta(\mathfrak p)])
$$
est le dual de Matlis de
$$
H^{-\delta(\mathfrak p) - i}_{\mathfrak p}(M_\mathfrak p)
$$
d'après Complexes dualisants, lemme \ref{dualizing-lemma-special-case-local-duality}. On en déduit en particulier le fait suivant : un idéal $J \subset A$ annule $(E^i)_\mathfrak p$ si et seulement si $J$ annule $H^{-\delta(\mathfrak p) - i}_{\mathfrak p}(M_\mathfrak p)$.

\medskip\noindent
Posons $T_n = \{\mathfrak p \in T \mid \delta(\mathfrak p) \leq n\}$. Puisque $\delta$ est une fonction bornée, on voit que $T_a = \emptyset$ pour $a \ll 0$ et $T_b = T$ pour $b \gg 0$.

\medskip\noindent
Supposons (2). Démontrons l'existence de $J$ comme dans (1). Nous utiliserons une double récurrence. Pour $i \leq s$, considérons l'hypothèse de récurrence $IH_i$ : $H^a_T(M)$ est annulé par un idéal $J \subset A$ tel que $V(J) \subset T'$ pour $0 \leq a \leq i$. Le cas $IH_0$ est trivial, car $H^0_T(M)$ est un sous-module de $M$, donc de type fini, et est par conséquent annulé par un idéal $J$ tel que $V(J) \subset T$.

\medskip\noindent
Étape de récurrence. Supposons $IH_{i - 1}$ vérifiée pour un certain $0 < i \leq s$. Choisissons $J'$ tel que $V(J') \subset T'$, annulant $H^a_T(M)$ pour $0 \leq a \leq i - 1$ (ce que permet l'hypothèse de récurrence). Montrons par récurrence descendante sur $n$ qu'il existe un idéal $J$ tel que $V(J) \subset T'$ et tel que les idéaux premiers associés à $J H^i_T(M)$ appartiennent à $T_n$. Pour $n \ll 0$, cela implique $JH^i_T(M) = 0$ (Algèbre, lemme \ref{algebra-lemma-ass-zero}), et donc $IH_i$ sera vérifiée. Le cas initial $n \gg 0$ est trivial, car $T = T_n$ dans ce cas, et tous les idéaux premiers associés à $H^i_T(M)$ appartiennent à $T$.

\medskip\noindent
Supposons donc donné $J$ ayant la propriété pour $n$. Soit $\mathfrak q \in T_n$. Soit $T_\mathfrak q \subset \Spec(A_\mathfrak q)$ l'image réciproque de $T$. On a $H^j_T(M)_\mathfrak q = H^j_{T_\mathfrak q}(M_\mathfrak q)$ d'après le lemme \ref{lemma-torsion-change-rings}. Considérons la suite spectrale
$$
H_\mathfrak q^p(H^q_{T_\mathfrak q}(M_\mathfrak q))
\Rightarrow
H^{p + q}_\mathfrak q(M_\mathfrak q)
$$
du lemme \ref{lemma-local-cohomology-ss}. Nous trouverons ci-dessous un idéal $J'' \subset A$ tel que $V(J'') \subset T'$ et tel que $H^i_\mathfrak q(M_\mathfrak q)$ soit annulé par $J''$ pour tout $\mathfrak q \in T_n \setminus T_{n - 1}$. Affirmation : $J (J')^i J''$ convient pour $n - 1$. En effet, soit $\mathfrak q \in T_n \setminus T_{n - 1}$. La suite spectrale ci-dessus définit une filtration
$$
E_\infty^{0, i} = E_{i + 2}^{0, i} \subset \ldots \subset E_3^{0, i} \subset
E_2^{0, i} = H^0_\mathfrak q(H^i_{T_\mathfrak q}(M_\mathfrak q))
$$
Le module $E_\infty^{0, i}$ est annulé par $J''$. Les sous-quotients $E_j^{0, i}/E_{j + 1}^{0, i}$ pour $i + 1 \geq j \geq 2$ sont annulés par $J'$, car le but de $d_j^{0, i}$ est un sous-quotient de
$$
H^j_\mathfrak q(H^{i - j + 1}_{T_\mathfrak q}(M_\mathfrak q)) =
H^j_\mathfrak q(H^{i - j + 1}_T(M)_\mathfrak q)
$$
et $H^{i - j + 1}_T(M)_\mathfrak q$ est annulé par $J'$, par le choix de $J'$. Enfin, par le choix de $J$, on a $J H^i_T(M)_\mathfrak q \subset H^0_\mathfrak q(H^i_T(M)_\mathfrak q)$, puisque les points non fermés de $\Spec(A_\mathfrak q)$ ont des valeurs de $\delta$ plus grandes. Ainsi $\mathfrak q$ ne peut être un idéal premier associé à $J(J')^iJ'' H^i_T(M)$, comme voulu.

\medskip\noindent
D'après les remarques initiales, $J''$ doit annuler
$$
(E^{-\delta(\mathfrak q) - i})_\mathfrak q =
(E^{-n - i})_\mathfrak q
$$
pour tout $\mathfrak q \in T_n \setminus T_{n - 1}$. Mais si $J''$ convient pour un $\mathfrak q$, il convient pour tous les $\mathfrak q$ d'un voisinage ouvert de $\mathfrak q$, puisque les modules $E^{-n - i}$ sont de type fini. Toute partie de $\Spec(A)$ étant noethérienne pour la topologie induite (Topologie, lemme \ref{topology-lemma-Noetherian}), il suffit donc de prouver l'existence de $J''$ pour un seul $\mathfrak q$.

\medskip\noindent
Les modules Ext étant de type fini, l'existence de $J''$ équivaut à
$$
\text{Supp}(E^{-n - i}) \cap \Spec(A_\mathfrak q) \subset T'.
$$
Cela revient à montrer que le localisé de $E^{-n - i}$ en tout $\mathfrak p \subset \mathfrak q$, $\mathfrak p \not \in T'$ est nul. Par dualité locale sur $A_\mathfrak p$, il faut donc montrer que
$$
H^{i + n - \delta(\mathfrak p)}_\mathfrak p(M_\mathfrak p) =
H^{i - \dim((A/\mathfrak p)_\mathfrak q)}_\mathfrak p(M_\mathfrak p)
$$
est nul (on utilise ici le fait que $\delta$ est une fonction de dimension). Cette annulation résulte de l'hypothèse du lemme, de $i \leq s$ et de Complexes dualisants, lemme \ref{dualizing-lemma-depth}.

\medskip\noindent
Pour démontrer l'implication réciproque, supposons que (2) ne soit pas vérifié et reprenons les arguments précédents en sens inverse. Choisissons d'abord $\mathfrak q \in T$, $\mathfrak p \subset \mathfrak q$, avec $\mathfrak p \not \in T'$, tels que
$$
i = \text{depth}_{A_\mathfrak p}(M_\mathfrak p) +
\dim((A/\mathfrak p)_\mathfrak q) \leq s
$$
soit minimal. Alors $H^{i - \dim((A/\mathfrak p)_\mathfrak q)}_\mathfrak p(M_\mathfrak p)$ est non nul d'après la non-annulation de Complexes dualisants, lemme \ref{dualizing-lemma-depth}. Posons $n = \delta(\mathfrak q)$. Il n'existe alors aucun idéal $J \subset A$ tel que $V(J) \subset T'$ et $J(E^{-n - i})_\mathfrak q = 0$. Ainsi $H^i_\mathfrak q(M_\mathfrak q)$ n'est annulé par aucun idéal $J \subset A$ tel que $V(J) \subset T'$. Par minimalité de $i$, la suite spectrale ci-dessus montre que le module $H^i_T(M)_\mathfrak q$ n'est annulé par aucun idéal $J \subset A$ tel que $V(J) \subset T'$. Ainsi $H^i_T(M)$ n'est annulé par aucun idéal $J \subset A$ tel que $V(J) \subset T'$. Cela achève la démonstration de la proposition.
\end{proof}

\begin{lemma}
\label{lemma-kill-local-cohomology-at-prime}
Soit $I$ un idéal d'un anneau noethérien $A$. Soient $M$ un $A$-module de type fini, $\mathfrak p \subset A$ un idéal premier et $s \geq 0$ un entier. Supposons que
\begin{enumerate}
\item $A$ possède un complexe dualisant,
\item $\mathfrak p \not \in V(I)$,
\item pour tous les idéaux premiers $\mathfrak p' \subset \mathfrak p$ et $\mathfrak q \in V(I)$ tels que $\mathfrak p' \subset \mathfrak q$, on ait
$$
\text{depth}_{A_{\mathfrak p'}}(M_{\mathfrak p'}) +
\dim((A/\mathfrak p')_\mathfrak q) > s
$$
\end{enumerate}
Alors il existe $f \in A$, $f \not \in \mathfrak p$, qui annule $H^i_{V(I)}(M)$ pour $i \leq s$.
\end{lemma}

\begin{proof}
Considérons les ensembles
$$
T = V(I)
\quad\text{et}\quad
T' = \bigcup\nolimits_{f \in A, f \not \in \mathfrak p} V(f)
$$
Ce sont des parties de $\Spec(A)$ stables par spécialisation. Observons que $T \subset T'$ et $\mathfrak p \not \in T'$. L'hypothèse (3) signifie que l'hypothèse (2) de la proposition \ref{proposition-annihilator} est vérifiée. On peut donc trouver $J \subset A$ tel que $V(J) \subset T'$ et tel que $J H^i_{V(I)}(M) = 0$ pour $i \leq s$. Choisissons $f \in A$, $f \not \in \mathfrak p$, tel que $V(J) \subset V(f)$. Une puissance de $f$ annule $H^i_{V(I)}(M)$ pour $i \leq s$.
\end{proof}





\section{Finitude de la cohomologie locale, II}
\label{section-finiteness-II}

\noindent
Nous poursuivons l'étude de la finitude de la cohomologie locale commencée à la section \ref{section-finiteness}. Le théorème des annulateurs de Faltings permet de démontrer aisément le résultat fondamental suivant.

\begin{proposition}
\label{proposition-finiteness}
\begin{reference}
\cite{Faltings-annulators}.
\end{reference}
Soit $A$ un anneau noethérien possédant un complexe dualisant. Soit $T \subset \Spec(A)$ une partie stable par spécialisation. Soit $s \geq 0$ un entier. Soit $M$ un $A$-module de type fini. Les conditions suivantes sont équivalentes :
\begin{enumerate}
\item $H^i_T(M)$ est un $A$-module de type fini pour $i \leq s$,
\item pour tous $\mathfrak p \not \in T$, $\mathfrak q \in T$ tels que $\mathfrak p \subset \mathfrak q$, on a
$$
\text{depth}_{A_\mathfrak p}(M_\mathfrak p) +
\dim((A/\mathfrak p)_\mathfrak q) > s
$$
\end{enumerate}
\end{proposition}

\begin{proof}
Conséquence formelle de la proposition \ref{proposition-annihilator} et du lemme \ref{lemma-check-finiteness-local-cohomology-by-annihilator}.
\end{proof}

\noindent
Outre quelques lemmes destinés à un usage ultérieur, le reste de cette section examine dans quelle mesure on peut affaiblir la condition de la proposition \ref{proposition-finiteness} selon laquelle $A$ possède un complexe dualisant. La réponse est, en gros, qu'il faut supposer que les fibres formelles de $A$ vérifient $(S_n)$ pour $n$ assez grand.

\medskip\noindent
Soient $A$ un anneau noethérien et $I \subset A$ un idéal. Posons $X = \Spec(A)$ et $Z = V(I) \subset X$. Soit $M$ un $A$-module de type fini. Définissons
\begin{equation}
\label{equation-cutoff}
s_{A, I}(M) =
\min \{
\text{depth}_{A_\mathfrak p}(M_\mathfrak p) + \dim((A/\mathfrak p)_\mathfrak q)
\mid
\mathfrak p \in X \setminus Z, \mathfrak q \in Z,
\mathfrak p \subset \mathfrak q
\}
\end{equation}
Notre convention sur la profondeur est que la profondeur de $0$ vaut $\infty$ ; il suffit donc de considérer les idéaux premiers $\mathfrak p$ du support de $M$. Nous verrons que $s_{A, I}(M)$ est un invariant important de la situation.

\begin{lemma}
\label{lemma-cutoff}
Soit $A \to B$ un homomorphisme fini d'anneaux noethériens. Soit $I \subset A$ un idéal et posons $J = IB$. Soit $M$ un $B$-module de type fini. Si $A$ est universellement caténaire, alors $s_{B, J}(M) = s_{A, I}(M)$.
\end{lemma}

\begin{proof}
Soient $\mathfrak p \subset \mathfrak q \subset A$ des idéaux premiers tels que $I \subset \mathfrak q$ et $I \not \subset \mathfrak p$. Puisque $A \to B$ est fini, les idéaux premiers $\mathfrak p_i$ au-dessus de $\mathfrak p$ sont en nombre fini. D'après Algèbre, lemme \ref{algebra-lemma-depth-goes-down-finite}, on a
$$
\text{depth}(M_\mathfrak p) = \min \text{depth}(M_{\mathfrak p_i})
$$
Soient $\mathfrak p_i \subset \mathfrak q_{ij}$ des idéaux premiers au-dessus de $\mathfrak q$. Par le théorème de montée pour $A \to B$ (Algèbre, lemme \ref{algebra-lemma-integral-going-up}), il existe au moins un $\mathfrak q_{ij}$ pour chaque $i$. Alors on a
$$
\dim((B/\mathfrak p_i)_{\mathfrak q_{ij}}) =
\dim((A/\mathfrak p)_\mathfrak q)
$$
d'après la formule des dimensions ; voir Algèbre, lemme \ref{algebra-lemma-dimension-formula}. Cela implique que le minimum des quantités utilisées pour définir $s_{B, J}(M)$ pour les couples $(\mathfrak p_i, \mathfrak q_{ij})$ est égal à la quantité correspondant au couple $(\mathfrak p, \mathfrak q)$. Cela prouve le lemme.
\end{proof}

\begin{lemma}
\label{lemma-change-completion}
Soit $A$ un anneau noethérien possédant un complexe dualisant. Soit $I \subset A$ un idéal. Soit $M$ un $A$-module de type fini. Soient $A', M'$ les complétés $I$-adiques de $A, M$. Soient $\mathfrak p' \subset \mathfrak q'$ des idéaux premiers de $A'$, avec $\mathfrak q' \in V(IA')$, au-dessus de $\mathfrak p \subset \mathfrak q$ dans $A$. Alors
$$
\text{depth}_{A_{\mathfrak p'}}(M'_{\mathfrak p'})
\geq
\text{depth}_{A_\mathfrak p}(M_\mathfrak p)
$$
et
$$
\text{depth}_{A_{\mathfrak p'}}(M'_{\mathfrak p'}) +
\dim((A'/\mathfrak p')_{\mathfrak q'}) =
\text{depth}_{A_\mathfrak p}(M_\mathfrak p) +
\dim((A/\mathfrak p)_\mathfrak q)
$$
\end{lemma}

\begin{proof}
On a
$$
\text{depth}(M'_{\mathfrak p'}) =
\text{depth}(M_\mathfrak p) +
\text{depth}(A'_{\mathfrak p'}/\mathfrak p A'_{\mathfrak p'})
\geq \text{depth}(M_\mathfrak p)
$$
par platitude de $A \to A'$ ; voir Algèbre, lemme \ref{algebra-lemma-apply-grothendieck-module}. Puisque les fibres de $A \to A'$ sont de Cohen-Macaulay (Complexes dualisants, lemme \ref{dualizing-lemma-dualizing-gorenstein-formal-fibres}, et Compléments d'algèbre, section \ref{more-algebra-section-properties-formal-fibres}), on voit que $\text{depth}(A'_{\mathfrak p'}/\mathfrak p A'_{\mathfrak p'}) =
\dim(A'_{\mathfrak p'}/\mathfrak p A'_{\mathfrak p'})$. On obtient donc
\begin{align*}
\text{depth}(M'_{\mathfrak p'}) +
\dim((A'/\mathfrak p')_{\mathfrak q'})
& =
\text{depth}(M_\mathfrak p) +
\dim(A'_{\mathfrak p'}/\mathfrak p A'_{\mathfrak p'}) +
\dim((A'/\mathfrak p')_{\mathfrak q'}) \\
& =
\text{depth}(M_\mathfrak p) +
\dim((A'/\mathfrak p A')_{\mathfrak q'}) \\
& =
\text{depth}(M_\mathfrak p) +
\dim((A/\mathfrak p)_\mathfrak q)
\end{align*}
La deuxième égalité vient du fait que $A'$ est caténaire et la troisième de Compléments d'algèbre, lemme \ref{more-algebra-lemma-completion-dimension}, puisque $(A/\mathfrak p)_\mathfrak q$ et $(A'/\mathfrak p A')_{\mathfrak q'}$ ont les mêmes complétés $I$-adiques.
\end{proof}





\begin{lemma}
\label{lemma-cutoff-completion}
Soit $A$ un anneau local noethérien universellement caténaire. Soit $I \subset A$ un idéal. Soit $M$ un $A$-module de type fini. Alors
$$
s_{A, I}(M) \geq s_{A^\wedge, I^\wedge}(M^\wedge)
$$
Si les fibres formelles de $A$ vérifient $(S_n)$, alors $\min(n + 1, s_{A, I}(M)) \leq s_{A^\wedge, I^\wedge}(M^\wedge)$.
\end{lemma}

\begin{proof}
Écrivons $X = \Spec(A)$, $X^\wedge = \Spec(A^\wedge)$, $Z = V(I) \subset X$ et $Z^\wedge = V(I^\wedge)$. Soient $\mathfrak p' \subset \mathfrak q' \subset A^\wedge$ des idéaux premiers tels que $\mathfrak p' \not \in Z^\wedge$ et $\mathfrak q' \in Z^\wedge$. Soient $\mathfrak p \subset \mathfrak q$ les idéaux premiers correspondants de $A$. Alors $\mathfrak p \not \in Z$ et $\mathfrak q \in Z$. On a le diagramme
$$
\xymatrix{
\mathfrak p' \ar[r] & \mathfrak q' \ar[r] & A^\wedge \\
\mathfrak p \ar[r] \ar@{-}[u] &
\mathfrak q \ar[r] \ar@{-}[u] & A \ar[u]
}
$$
Écrivons
\begin{align*}
a & = \dim(A/\mathfrak p) = \dim(A^\wedge/\mathfrak pA^\wedge),\\
b & = \dim(A/\mathfrak q) = \dim(A^\wedge/\mathfrak qA^\wedge),\\
a' & = \dim(A^\wedge/\mathfrak p'),\\
b' & = \dim(A^\wedge/\mathfrak q')
\end{align*}
Les égalités résultent de Compléments d'algèbre, lemme \ref{more-algebra-lemma-completion-dimension}. Écrivons aussi
\begin{align*}
p & = \dim(A^\wedge_{\mathfrak p'}/\mathfrak p A^\wedge_{\mathfrak p'}) =
\dim((A^\wedge/\mathfrak p A^\wedge)_{\mathfrak p'}) \\
q & = \dim(A^\wedge_{\mathfrak q'}/\mathfrak p A^\wedge_{\mathfrak q'}) =
\dim((A^\wedge/\mathfrak q A^\wedge)_{\mathfrak q'})
\end{align*}
Puisque $A$ est universellement caténaire, $A^\wedge/\mathfrak pA^\wedge = (A/\mathfrak p)^\wedge$ est équidimensionnel de dimension $a$ (Compléments d'algèbre, proposition \ref{more-algebra-proposition-ratliff}). Ainsi $a = a' + p$. De même, $b = b' + q$. D'après Algèbre, lemme \ref{algebra-lemma-apply-grothendieck-module}, appliqué à l'homomorphisme local plat $A_\mathfrak p \to A^\wedge_{\mathfrak p'}$, on a
$$
\text{depth}(M^\wedge_{\mathfrak p'})
=
\text{depth}(M_\mathfrak p) +
\text{depth}(A^\wedge_{\mathfrak p'} / \mathfrak p A^\wedge_{\mathfrak p'})
$$
La quantité dont on prend le minimum pour $s_{A, I}(M)$ est
$$
s(\mathfrak p, \mathfrak q) =
\text{depth}(M_\mathfrak p) + \dim((A/\mathfrak p)_\mathfrak q) =
\text{depth}(M_\mathfrak p) + a - b
$$
(la dernière égalité vient du fait que $A$ est caténaire). La quantité dont on prend le minimum pour $s_{A^\wedge, I^\wedge}(M^\wedge)$ est
$$
s(\mathfrak p', \mathfrak q') =
\text{depth}(M^\wedge_{\mathfrak p'}) +
\dim((A^\wedge/\mathfrak p')_{\mathfrak q'}) =
\text{depth}(M^\wedge_{\mathfrak p'}) + a' - b'
$$
(la dernière égalité vient du fait que $A^\wedge$ est caténaire). Nous disposons maintenant des notations nécessaires pour commencer la démonstration.

\medskip\noindent
Soient $\mathfrak p \subset \mathfrak q \subset A$ des idéaux premiers tels que $\mathfrak p \not \in Z$, $\mathfrak q \in Z$ et $s_{A, I}(M) = s(\mathfrak p, \mathfrak q)$. On peut alors choisir $\mathfrak q'$ minimal au-dessus de $\mathfrak q A^\wedge$ et $\mathfrak p' \subset \mathfrak q'$ minimal au-dessus de $\mathfrak p A^\wedge$ (par le théorème de descente pour $A \to A^\wedge$). On dispose alors de quatre idéaux premiers comme ci-dessus, avec $p = 0$ et $q = 0$. De plus, on a $\text{depth}(A^\wedge_{\mathfrak p'} / \mathfrak p A^\wedge_{\mathfrak p'})=0$, là encore parce que $p = 0$. Cela signifie que $s(\mathfrak p', \mathfrak q') = s(\mathfrak p, \mathfrak q)$. On obtient ainsi la première inégalité.

\medskip\noindent
Supposons que les fibres formelles de $A$ vérifient $(S_n)$. Alors $\text{depth}(A^\wedge_{\mathfrak p'} / \mathfrak p A^\wedge_{\mathfrak p'})
\geq \min(n, p)$. D'où
$$
s(\mathfrak p', \mathfrak q') \geq
s(\mathfrak p, \mathfrak q) + q + \min(n, p) - p \geq
s_{A, I}(M) + q + \min(n, p) - p
$$
Le seul cas susceptible de poser problème est donc $p > n$. Dans ce cas,
\begin{align*}
s(\mathfrak p', \mathfrak q')
& =
\text{depth}(M^\wedge_{\mathfrak p'}) +
\dim((A^\wedge/\mathfrak p')_{\mathfrak q'}) \\
& =
\text{depth}(M_\mathfrak p) +
\text{depth}(A^\wedge_{\mathfrak p'} / \mathfrak p A^\wedge_{\mathfrak p'}) +
\dim((A^\wedge/\mathfrak p')_{\mathfrak q'}) \\
& \geq
0 + n + 1
\end{align*}
car $(A^\wedge/\mathfrak p')_{\mathfrak q'}$ possède au moins deux idéaux premiers. Cela prouve la seconde inégalité.
\end{proof}

\noindent
La méthode de démonstration du lemme suivant s'applique plus généralement, mais les résultats plus forts ainsi obtenus seront englobés par le théorème \ref{theorem-finiteness} ci-dessous.

\begin{lemma}
\label{lemma-local-annihilator}
\begin{reference}
C'est un cas particulier de \cite[Satz 1]{Faltings-annulators}.
\end{reference}
Soit $A$ un anneau local noethérien de Gorenstein. Soit $I \subset A$ un idéal et posons $Z = V(I) \subset \Spec(A)$. Soit $M$ un $A$-module de type fini. Soit $s = s_{A, I}(M)$ comme dans (\ref{equation-cutoff}). Alors $H^i_Z(M)$ est de type fini pour $i < s$, mais $H^s_Z(M)$ n'est pas de type fini.
\end{lemma}

\begin{proof}
Puisqu'un anneau local de Gorenstein possède un complexe dualisant, c'est un cas particulier de la proposition \ref{proposition-finiteness}. Il serait utile de disposer d'une démonstration courte de ce cas particulier, qui interviendra ci-dessous dans la démonstration d'un théorème général de finitude.
\end{proof}

\noindent
Observons que les hypothèses du théorème suivant sont satisfaites par les anneaux noethériens excellents (par définition), par les anneaux noethériens possédant un complexe dualisant (Complexes dualisants, lemme \ref{dualizing-lemma-universally-catenary}, et Complexes dualisants, lemme \ref{dualizing-lemma-dualizing-gorenstein-formal-fibres}), ainsi que par les quotients d'anneaux noethériens réguliers.

\begin{theorem}
\label{theorem-finiteness}
\begin{reference}
C'est un cas particulier de \cite[Satz 2]{Faltings-finiteness}.
\end{reference}
Soient $A$ un anneau noethérien et $I \subset A$ un idéal. Posons $Z = V(I) \subset \Spec(A)$. Soit $M$ un $A$-module de type fini. Posons $s = s_{A, I}(M)$ comme dans (\ref{equation-cutoff}). Supposons que
\begin{enumerate}
\item $A$ soit universellement caténaire,
\item les fibres formelles des anneaux locaux de $A$ soient de Cohen-Macaulay.
\end{enumerate}
Alors $H^i_Z(M)$ est de type fini pour $0 \leq i < s$ et $H^s_Z(M)$ n'est pas de type fini.
\end{theorem}

\begin{proof}
D'après le lemme \ref{lemma-check-finiteness-local-cohomology-locally}, on peut supposer que $A$ est un anneau local.

\medskip\noindent
Si $A$ est un anneau local noethérien complet, on peut écrire $A$ comme quotient d'un anneau local régulier complet $B$ par le théorème de structure de Cohen (Algèbre, théorème \ref{algebra-theorem-cohen-structure-theorem}). À l'aide du lemme \ref{lemma-cutoff} et de Complexes dualisants, lemme \ref{dualizing-lemma-local-cohomology-and-restriction}, on se ramène au cas d'un anneau local régulier, qui résulte du lemme \ref{lemma-local-annihilator} puisqu'un anneau local régulier est de Gorenstein (Complexes dualisants, lemme \ref{dualizing-lemma-regular-gorenstein}).

\medskip\noindent
Soit $A$ un anneau local noethérien. Soit $\mathfrak m$ son idéal maximal. On peut supposer $I \subset \mathfrak m$, sinon le lemme est trivial. Soit $A^\wedge$ le complété de $A$, soit $Z^\wedge = V(IA^\wedge)$, et soit $M^\wedge = M \otimes_A A^\wedge$ le complété de $M$ (Algèbre, lemme \ref{algebra-lemma-completion-tensor}). Alors $H^i_Z(M) \otimes_A A^\wedge = H^i_{Z^\wedge}(M^\wedge)$ d'après Complexes dualisants, lemme \ref{dualizing-lemma-torsion-change-rings}, et la platitude de $A \to A^\wedge$ (Algèbre, lemme \ref{algebra-lemma-completion-flat}). Il suffit donc de montrer que $H^i_{Z^\wedge}(M^\wedge)$ est de type fini pour $i < s$ et n'est pas de type fini pour $i = s$ ; voir Algèbre, lemme \ref{algebra-lemma-descend-properties-modules}. Puisque le résultat est connu pour $A^\wedge$, il suffit de montrer que $s_{A, I}(M) = s_{A^\wedge, I^\wedge}(M^\wedge)$. Cela résulte du lemme \ref{lemma-cutoff-completion}.
\end{proof}

\begin{remark}
\label{remark-astute-reader}
Le lecteur attentif aura remarqué qu'une condition légèrement plus faible sur les fibres formelles des anneaux locaux de $A$ suffit. Plus précisément, dans la situation du théorème \ref{theorem-finiteness}, supposons $A$ universellement caténaire, sans faire d'hypothèse sur les fibres formelles. Supposons donné $n$ et cherchons à montrer que les $H^i_Z(M)$ sont de type fini pour $i \leq n$. La même démonstration montre qu'il suffit d'avoir $s_{A, I}(M) > n$ et que les fibres formelles des anneaux locaux de $A$ vérifient $(S_n)$. D'autre part, si l'on veut montrer que $H^s_Z(M)$ n'est pas de type fini, où $s = s_{A, I}(M)$, nos arguments le prouvent lorsque les fibres formelles vérifient $(S_{s - 1})$.
\end{remark}







\section{Finitude des images directes, II}
\label{section-finiteness-pushforward-II}

\noindent
Cette section poursuit la section \ref{section-finiteness-pushforward}. Nous y recueillons les fruits du travail effectué à la section \ref{section-finiteness-II}.

\begin{lemma}
\label{lemma-finiteness-Rjstar}
Soit $X$ un schéma localement noethérien. Soit $j : U \to X$ l'inclusion d'un sous-schéma ouvert de complémentaire $Z$. Soit $\mathcal{F}$ un $\mathcal{O}_U$-module cohérent. Soit $n \geq 0$ un entier. Supposons que
\begin{enumerate}
\item $X$ soit universellement caténaire,
\item pour tout $z \in Z$, les fibres formelles de $\mathcal{O}_{X, z}$ vérifient $(S_n)$.
\end{enumerate}
Dans cette situation, les conditions suivantes sont équivalentes :
\begin{enumerate}
\item[(a)] pour $x \in \text{Supp}(\mathcal{F})$ et $z \in Z \cap \overline{\{x\}}$, on a $\text{depth}_{\mathcal{O}_{X, x}}(\mathcal{F}_x) +
\dim(\mathcal{O}_{\overline{\{x\}}, z}) > n$,
\item[(b)] $R^pj_*\mathcal{F}$ est cohérent pour $0 \leq p < n$.
\end{enumerate}
\end{lemma}

\begin{proof}
L'énoncé est local sur $X$ ; on peut donc supposer $X$ affine. Écrivons $X = \Spec(A)$ et $Z = V(I)$. Soit $M$ un $A$-module de type fini dont le $\mathcal{O}_X$-module cohérent associé se restreint à $\mathcal{F}$ sur $U$ ; voir le lemme \ref{lemma-finiteness-pushforwards-and-H1-local}. Ce lemme dit aussi que $R^pj_*\mathcal{F}$ est cohérent si et seulement si $H^{p + 1}_Z(M)$ est un $A$-module de type fini. Observons que le minimum des expressions $\text{depth}_{\mathcal{O}_{X, x}}(\mathcal{F}_x) +
\dim(\mathcal{O}_{\overline{\{x\}}, z})$ est le nombre $s_{A, I}(M)$ de (\ref{equation-cutoff}). Cela étant dit, le lemme résulte du théorème \ref{theorem-finiteness}, avec les précisions apportées par la remarque \ref{remark-astute-reader}.
\end{proof}

\begin{lemma}
\label{lemma-finiteness-for-finite-locally-free}
Soit $X$ un schéma localement noethérien. Soit $j : U \to X$ l'inclusion d'un sous-schéma ouvert de complémentaire $Z$. Soit $n \geq 0$ un entier. Si $R^pj_*\mathcal{O}_U$ est cohérent pour $0 \leq p < n$, il en est de même de $R^pj_*\mathcal{F}$, $0 \leq p < n$, pour tout $\mathcal{O}_U$-module localement libre de type fini $\mathcal{F}$.
\end{lemma}

\begin{proof}
La question est locale sur $X$ ; on peut donc supposer $X$ affine. Écrivons $X = \Spec(A)$ et $Z = V(I)$. Par l'intermédiaire du lemme \ref{lemma-finiteness-pushforwards-and-H1-local}, notre lemme résulte du lemme \ref{lemma-local-finiteness-for-finite-locally-free}.
\end{proof}

\begin{lemma}
\label{lemma-annihilate-Hp}
\begin{reference}
\cite[Lemma 1.9]{Bhatt-local}
\end{reference}
Soient $A$ un anneau et $J \subset I \subset A$ des idéaux de type fini. Soit $p \geq 0$ un entier. Posons $U = \Spec(A) \setminus V(I)$. Si $H^p(U, \mathcal{O}_U)$ est annulé par $J^n$ pour un certain $n$, alors $H^p(U, \mathcal{F})$ est annulé par $J^m$ pour un certain $m = m(\mathcal{F})$, pour tout $\mathcal{O}_U$-module localement libre de type fini $\mathcal{F}$.
\end{lemma}

\begin{proof}
Considérons l'annulateur $\mathfrak a$ de $H^p(U, \mathcal{F})$. Soit $u \in U$. Il existe un voisinage ouvert $u \in U' \subset U$ et un isomorphisme $\varphi : \mathcal{O}_{U'}^{\oplus r} \to \mathcal{F}|_{U'}$. Choisissons $f \in A$ tel que $u \in D(f) \subset U'$. Il existe des applications
$$
a : \mathcal{O}_U^{\oplus r} \longrightarrow \mathcal{F}
\quad\text{et}\quad
b : \mathcal{F} \longrightarrow \mathcal{O}_U^{\oplus r}
$$
dont les restrictions à $D(f)$ sont égales à $f^N \varphi$ et $f^N \varphi^{-1}$ pour un certain $N$. De plus, on peut supposer que $a \circ b$ et $b \circ a$ sont égales à la multiplication par $f^{2N}$. Cela résulte de Propriétés, lemme \ref{properties-lemma-section-maps-backwards}, puisque $U$ est quasi-compact ($I$ est de type fini), séparé, et que $\mathcal{F}$ et $\mathcal{O}_U^{\oplus r}$ sont de présentation finie. On voit donc que $H^p(U, \mathcal{F})$ est annulé par $f^{2N}J^n$, c'est-à-dire que $f^{2N}J^n \subset \mathfrak a$.

\medskip\noindent
Puisque $U$ est quasi-compact, on peut trouver un nombre fini de $f_1, \ldots, f_t$ et de $N_1, \ldots, N_t$ tels que $U = \bigcup D(f_i)$ et $f_i^{2N_i}J^n \subset \mathfrak a$. Alors $V(I) = V(f_1, \ldots, f_t)$ et, puisque $I$ est de type fini, on en déduit $I^M \subset (f_1, \ldots, f_t)$ pour un certain $M$. Finalement, on voit que $J^m \subset \mathfrak a$ pour $m \gg 0$ ; par exemple, $m = M (2N_1 + \ldots + 2N_t) n$ convient.
\end{proof}











\section{Annulateurs de la cohomologie locale, II}
\label{section-annihilators-II}

\noindent
Nous étendons la discussion des annulateurs de la cohomologie locale de la section \ref{section-annihilators} aux complexes bornés inférieurement dont les modules de cohomologie sont de type fini.

\begin{definition}
\label{definition-depth-complex}
Soit $I$ un idéal d'un anneau noethérien $A$. Soit $K \in D^+_{\textit{Coh}}(A)$. On définit la {\it $I$-profondeur} de $K$, notée $\text{depth}_I(K)$, comme le plus grand $m \in \mathbf{Z} \cup \{\infty\}$ tel que $H^i_I(K) = 0$ pour tout $i < m$. Si $A$ est local d'idéal maximal $\mathfrak m$, on appelle simplement $\text{depth}_\mathfrak m(K)$ la {\it profondeur} de $K$.
\end{definition}

\noindent
Cette définition est compatible avec Algèbre, définition \ref{algebra-definition-depth}, d'après Complexes dualisants, lemme \ref{dualizing-lemma-depth}.

\begin{proposition}
\label{proposition-annihilator-complex}
Soit $A$ un anneau noethérien possédant un complexe dualisant. Soient $T \subset T' \subset \Spec(A)$ des parties stables par spécialisation. Soit $s \in \mathbf{Z}$. Soit $K$ un objet de $D_{\textit{Coh}}^+(A)$. Les conditions suivantes sont équivalentes :
\begin{enumerate}
\item il existe un idéal $J \subset A$ tel que $V(J) \subset T'$ et tel que $J$ annule $H^i_T(K)$ pour $i \leq s$,
\item pour tous $\mathfrak p \not \in T'$, $\mathfrak q \in T$ tels que $\mathfrak p \subset \mathfrak q$, on a
$$
\text{depth}_{A_\mathfrak p}(K_\mathfrak p) +
\dim((A/\mathfrak p)_\mathfrak q) > s
$$
\end{enumerate}
\end{proposition}

\begin{proof}
Ce lemme est la généralisation naturelle de la proposition \ref{proposition-annihilator}, dont le lecteur est invité à lire d'abord la démonstration. Soit $\omega_A^\bullet$ un complexe dualisant. Soit $\delta$ sa fonction de dimension ; voir Complexes dualisants, section \ref{dualizing-section-dimension-function}. Les $A$-modules de type fini
$$
E^i = \Ext_A^i(K, \omega_A^\bullet)
$$
joueront un rôle important. Pour $\mathfrak p \subset A$, nous écrirons $H^i_\mathfrak p$ pour désigner la cohomologie locale d'un objet de $D(A_\mathfrak p)$ relativement à $\mathfrak pA_\mathfrak p$. On voit alors que le complété $\mathfrak pA_\mathfrak p$-adique de
$$
(E^i)_\mathfrak p =
\Ext^{\delta(\mathfrak p) + i}_{A_\mathfrak p}(K_\mathfrak p,
(\omega_A^\bullet)_\mathfrak p[-\delta(\mathfrak p)])
$$
est le dual de Matlis de
$$
H^{-\delta(\mathfrak p) - i}_{\mathfrak p}(K_\mathfrak p)
$$
d'après Complexes dualisants, lemme \ref{dualizing-lemma-special-case-local-duality}. On en déduit en particulier le fait suivant : un idéal $J \subset A$ annule $(E^i)_\mathfrak p$ si et seulement si $J$ annule $H^{-\delta(\mathfrak p) - i}_{\mathfrak p}(K_\mathfrak p)$.

\medskip\noindent
Posons $T_n = \{\mathfrak p \in T \mid \delta(\mathfrak p) \leq n\}$. Puisque $\delta$ est une fonction bornée, on voit que $T_a = \emptyset$ pour $a \ll 0$ et $T_b = T$ pour $b \gg 0$.

\medskip\noindent
Supposons (2). Démontrons l'existence de $J$ comme dans (1). Nous utiliserons une double récurrence. Pour $i \leq s$, considérons l'hypothèse de récurrence $IH_i$ : $H^a_T(K)$ est annulé par un idéal $J \subset A$ tel que $V(J) \subset T'$ pour $a \leq i$. Le cas $IH_i$ est trivial pour $i$ assez petit, puisque $K$ est borné inférieurement.

\medskip\noindent
Étape de récurrence. Supposons $IH_{i - 1}$ vérifiée pour un certain $i \leq s$. Choisissons $J'$ tel que $V(J') \subset T'$, annulant $H^a_T(K)$ pour $a \leq i - 1$ (ce que permet l'hypothèse de récurrence). Montrons par récurrence descendante sur $n$ qu'il existe un idéal $J$ tel que $V(J) \subset T'$ et tel que les idéaux premiers associés à $J H^i_T(K)$ appartiennent à $T_n$. Pour $n \ll 0$, cela implique $JH^i_T(K) = 0$ (Algèbre, lemme \ref{algebra-lemma-ass-zero}), et donc $IH_i$ sera vérifiée. Le cas initial $n \gg 0$ est trivial, car $T = T_n$ dans ce cas, et tous les idéaux premiers associés à $H^i_T(K)$ appartiennent à $T$.

\medskip\noindent
Supposons donc donné $J$ ayant la propriété pour $n$. Soit $\mathfrak q \in T_n$. Soit $T_\mathfrak q \subset \Spec(A_\mathfrak q)$ l'image réciproque de $T$. On a $H^j_T(K)_\mathfrak q = H^j_{T_\mathfrak q}(K_\mathfrak q)$ d'après le lemme \ref{lemma-torsion-change-rings}. Considérons la suite spectrale
$$
H_\mathfrak q^p(H^q_{T_\mathfrak q}(K_\mathfrak q))
\Rightarrow
H^{p + q}_\mathfrak q(K_\mathfrak q)
$$
du lemme \ref{lemma-local-cohomology-ss}. Nous trouverons ci-dessous un idéal $J'' \subset A$ tel que $V(J'') \subset T'$ et tel que $H^i_\mathfrak q(K_\mathfrak q)$ soit annulé par $J''$ pour tout $\mathfrak q \in T_n \setminus T_{n - 1}$. Affirmation : $J (J')^i J''$ convient pour $n - 1$. En effet, soit $\mathfrak q \in T_n \setminus T_{n - 1}$. La suite spectrale ci-dessus définit une filtration
$$
E_\infty^{0, i} = E_{i + 2}^{0, i} \subset \ldots \subset E_3^{0, i} \subset
E_2^{0, i} = H^0_\mathfrak q(H^i_{T_\mathfrak q}(K_\mathfrak q))
$$
Le module $E_\infty^{0, i}$ est annulé par $J''$. Les sous-quotients $E_j^{0, i}/E_{j + 1}^{0, i}$ pour $i + 1 \geq j \geq 2$ sont annulés par $J'$, car le but de $d_j^{0, i}$ est un sous-quotient de
$$
H^j_\mathfrak q(H^{i - j + 1}_{T_\mathfrak q}(K_\mathfrak q)) =
H^j_\mathfrak q(H^{i - j + 1}_T(K)_\mathfrak q)
$$
et $H^{i - j + 1}_T(K)_\mathfrak q$ est annulé par $J'$, par le choix de $J'$. Enfin, par le choix de $J$, on a $J H^i_T(K)_\mathfrak q \subset H^0_\mathfrak q(H^i_T(K)_\mathfrak q)$, puisque les points non fermés de $\Spec(A_\mathfrak q)$ ont des valeurs de $\delta$ plus grandes. Ainsi $\mathfrak q$ ne peut être un idéal premier associé à $J(J')^iJ'' H^i_T(K)$, comme voulu.

\medskip\noindent
D'après les remarques initiales, $J''$ doit annuler
$$
(E^{-\delta(\mathfrak q) - i})_\mathfrak q =
(E^{-n - i})_\mathfrak q
$$
pour tout $\mathfrak q \in T_n \setminus T_{n - 1}$. Mais si $J''$ convient pour un $\mathfrak q$, il convient pour tous les $\mathfrak q$ d'un voisinage ouvert de $\mathfrak q$, puisque les modules $E^{-n - i}$ sont de type fini. Toute partie de $\Spec(A)$ étant noethérienne pour la topologie induite (Topologie, lemme \ref{topology-lemma-Noetherian}), il suffit donc de prouver l'existence de $J''$ pour un seul $\mathfrak q$.

\medskip\noindent
Les modules Ext étant de type fini, l'existence de $J''$ équivaut à
$$
\text{Supp}(E^{-n - i}) \cap \Spec(A_\mathfrak q) \subset T'.
$$
Cela revient à montrer que le localisé de $E^{-n - i}$ en tout $\mathfrak p \subset \mathfrak q$, $\mathfrak p \not \in T'$ est nul. Par dualité locale sur $A_\mathfrak p$, il faut donc montrer que
$$
H^{i + n - \delta(\mathfrak p)}_\mathfrak p(K_\mathfrak p) =
H^{i - \dim((A/\mathfrak p)_\mathfrak q)}_\mathfrak p(K_\mathfrak p)
$$
est nul (on utilise ici le fait que $\delta$ est une fonction de dimension). Cette annulation résulte de l'hypothèse du lemme, de $i \leq s$ et de notre définition de la profondeur, définition \ref{definition-depth-complex}.

\medskip\noindent
Pour démontrer l'implication réciproque, supposons que (2) ne soit pas vérifié et reprenons les arguments précédents en sens inverse. Choisissons d'abord $\mathfrak q \in T$, $\mathfrak p \subset \mathfrak q$, avec $\mathfrak p \not \in T'$, tels que
$$
i = \text{depth}_{A_\mathfrak p}(K_\mathfrak p) +
\dim((A/\mathfrak p)_\mathfrak q) \leq s
$$
soit minimal. Alors $H^{i - \dim((A/\mathfrak p)_\mathfrak q)}_\mathfrak p(K_\mathfrak p)$ est non nul d'après notre définition de la profondeur, définition \ref{definition-depth-complex}. Posons $n = \delta(\mathfrak q)$. Il n'existe alors aucun idéal $J \subset A$ tel que $V(J) \subset T'$ et $J(E^{-n - i})_\mathfrak q = 0$. Ainsi $H^i_\mathfrak q(K_\mathfrak q)$ n'est annulé par aucun idéal $J \subset A$ tel que $V(J) \subset T'$. Par minimalité de $i$, la suite spectrale ci-dessus montre que le module $H^i_T(K)_\mathfrak q$ n'est annulé par aucun idéal $J \subset A$ tel que $V(J) \subset T'$. Ainsi $H^i_T(K)$ n'est annulé par aucun idéal $J \subset A$ tel que $V(J) \subset T'$. Cela achève la démonstration de la proposition.
\end{proof}







\section{Finitude de la cohomologie locale, III}
\label{section-finiteness-III}

\noindent
Nous étendons la discussion de la finitude de la cohomologie locale des sections \ref{section-finiteness} et \ref{section-finiteness-II} aux complexes bornés inférieurement dont les modules de cohomologie sont de type fini.

\begin{lemma}
\label{lemma-check-finiteness-local-cohomology-by-annihilator-complex}
Soit $A$ un anneau noethérien. Soit $T \subset \Spec(A)$ une partie stable par spécialisation. Soit $K$ un objet de $D_{\textit{Coh}}^+(A)$. Soit $n \in \mathbf{Z}$. Les conditions suivantes sont équivalentes :
\begin{enumerate}
\item $H^i_T(K)$ est de type fini pour $i \leq n$,
\item il existe un idéal $J \subset A$ tel que $V(J) \subset T$ et tel que $J$ annule $H^i_T(K)$ pour $i \leq n$.
\end{enumerate}
Si $T = V(I) = Z$ pour un idéal $I \subset A$, ces conditions sont également équivalentes à
\begin{enumerate}
\item[(3)] il existe $e \geq 0$ tel que $I^e$ annule $H^i_Z(K)$ pour $i \leq n$.
\end{enumerate}
\end{lemma}

\begin{proof}
Ce lemme est la généralisation naturelle du lemme \ref{lemma-check-finiteness-local-cohomology-by-annihilator}, dont le lecteur est invité à lire d'abord la démonstration. Supposons (1) vrai. Rappelons que $H^i_J(K) = H^i_{V(J)}(K)$ ; voir Complexes dualisants, lemme \ref{dualizing-lemma-local-cohomology-noetherian}. Ainsi $H^i_T(K) = \colim H^i_J(K)$, où la limite inductive porte sur les idéaux $J \subset A$ tels que $V(J) \subset T$ ; voir le lemme \ref{lemma-adjoint-ext}. Puisque $H^i_T(K)$ est de type fini pour $i \leq n$, on peut trouver $J \subset A$ comme dans (2) tel que $H^i_J(K) \to H^i_T(K)$ soit surjective pour $i \leq n$. Les éléments de la liste finie de générateurs sont donc de torsion par rapport aux puissances de $J$, et (2) est vérifié en remplaçant $J$ par une puissance convenable.

\medskip\noindent
Soit $a \in \mathbf{Z}$ un entier tel que $H^i(K) = 0$ pour $i < a$. Démontrons (2) $\Rightarrow$ (1) par récurrence descendante sur $a$. Si $a > n$, on a $H^i_T(K) = 0$ pour $i \leq n$ ; ainsi (1) et (2) sont tous deux vrais, et il n'y a rien à démontrer.

\medskip\noindent
Supposons donné $J$ comme dans (2). Observons que $N = H^a_T(K) = H^0_T(H^a(K))$ est de type fini, comme sous-module du $A$-module de type fini $H^a(K)$. Si $n = a$, la démonstration est terminée ; supposons donc désormais $a < n$. Par construction de $R\Gamma_T$, on a $H^i_T(N) = 0$ pour $i > 0$ et $H^0_T(N) = N$ ; voir la remarque \ref{remark-upshot}. Choisissons un triangle distingué
$$
N[-a] \to K \to K' \to N[-a + 1]
$$
On voit alors que $H^a_T(K') = 0$ et $H^i_T(K) = H^i_T(K')$ pour $i > a$. On peut donc remplacer $K$ par $K'$. Ainsi on peut supposer $H^a_T(K) = 0$. Cela signifie qu'aucun des idéaux premiers associés à $H^a(K)$, qui sont en nombre fini, n'appartient à $T$. Par évitement des idéaux premiers (Algèbre, lemme \ref{algebra-lemma-silly}), on peut trouver $f \in J$ n'appartenant à aucun des idéaux premiers associés à $H^a(K)$. Choisissons un triangle distingué
$$
L \to K \xrightarrow{f} K \to L[1]
$$
Par construction, on a $H^i(L) = 0$ pour $i \leq a$. D'autre part, on a une suite exacte longue de cohomologie
$$
0 \to H^{a + 1}_T(L) \to H^{a + 1}_T(K) \xrightarrow{f}
H^{a + 1}_T(K) \to H^{a + 2}_T(L) \to H^{a + 2}_T(K) \xrightarrow{f} \ldots
$$
qui se décompose en l'identification $H^{a + 1}_T(L) = H^{a + 1}_T(K)$ et en suites exactes courtes
$$
0 \to H^{i - 1}_T(K) \to H^i_T(L) \to H^i_T(K) \to 0
$$
pour $i \leq n$, puisque $f \in J$. On en déduit que $J^2$ annule $H^i_T(L)$ pour $i \leq n$. Par l'hypothèse de récurrence appliquée à $L$, on voit que $H^i_T(L)$ est de type fini pour $i \leq n$. En utilisant une nouvelle fois la suite exacte courte, on voit que $H^i_T(K)$ est de type fini pour $i \leq n$, comme voulu.

\medskip\noindent
Nous omettons la démonstration de l'équivalence de (2) et (3) dans le cas $T = V(I)$.
\end{proof}

\begin{proposition}
\label{proposition-finiteness-complex}
Soit $A$ un anneau noethérien possédant un complexe dualisant. Soit $T \subset \Spec(A)$ une partie stable par spécialisation. Soit $s \in \mathbf{Z}$. Soit $K \in D_{\textit{Coh}}^+(A)$. Les conditions suivantes sont équivalentes :
\begin{enumerate}
\item $H^i_T(K)$ est un $A$-module de type fini pour $i \leq s$,
\item pour tous $\mathfrak p \not \in T$, $\mathfrak q \in T$ tels que $\mathfrak p \subset \mathfrak q$, on a
$$
\text{depth}_{A_\mathfrak p}(K_\mathfrak p) +
\dim((A/\mathfrak p)_\mathfrak q) > s
$$
\end{enumerate}
\end{proposition}

\begin{proof}
Conséquence formelle de la proposition \ref{proposition-annihilator-complex} et du lemme \ref{lemma-check-finiteness-local-cohomology-by-annihilator-complex}.
\end{proof}






\section{Amélioration des modules cohérents}
\label{section-improve}

\noindent
Des constructions analogues figurent dans \cite{EGA} et, plus récemment, dans \cite{Kollar-local-global-hulls} et \cite{Kollar-variants}.

\begin{lemma}
\label{lemma-get-depth-1-along-Z}
Soit $X$ un schéma noethérien. Soit $T \subset X$ une partie stable par spécialisation. Soit $\mathcal{F}$ un $\mathcal{O}_X$-module cohérent. Alors il existe une unique application $\mathcal{F} \to \mathcal{F}'$ de $\mathcal{O}_X$-modules cohérents telle que
\begin{enumerate}
\item $\mathcal{F} \to \mathcal{F}'$ soit surjective,
\item $\mathcal{F}_x \to \mathcal{F}'_x$ soit un isomorphisme pour $x \not \in T$,
\item $\text{depth}_{\mathcal{O}_{X, x}}(\mathcal{F}'_x) \geq 1$ pour $x \in T$.
\end{enumerate}
Si $f : Y \to X$ est un morphisme plat, avec $Y$ noethérien, alors $f^*\mathcal{F} \to f^*\mathcal{F}'$ est le quotient correspondant pour $f^{-1}(T) \subset Y$ et $f^*\mathcal{F}$.
\end{lemma}

\begin{proof}
La condition (3) signifie simplement que $\text{Ass}(\mathcal{F}') \cap T = \emptyset$. Ainsi $\mathcal{F} \to \mathcal{F}'$ est le quotient de $\mathcal{F}$ par le sous-faisceau des sections dont le support est contenu dans $T$. Cela prouve l'unicité. L'assertion concernant les images réciproques résulte de Diviseurs, lemme \ref{divisors-lemma-bourbaki}, et de l'unicité.

\medskip\noindent
Existence de $\mathcal{F} \to \mathcal{F}'$. Par unicité, il suffit de prouver l'existence et l'unicité localement sur $X$ ; nous omettons ce petit détail. On peut donc supposer $X = \Spec(A)$ affine et $\mathcal{F}$ égal au module cohérent associé au $A$-module de type fini $M$. Posons $M' = M / H^0_T(M)$, avec $H^0_T(M)$ comme dans la section \ref{section-supports}. Alors $M_\mathfrak p = M'_\mathfrak p$ pour $\mathfrak p \not \in T$, ce qui prouve (1). D'autre part, on a $H^0_T(M) = \colim H^0_Z(M)$, où $Z$ parcourt les parties fermées de $X$ contenues dans $T$. Ainsi, d'après Complexes dualisants, lemme \ref{dualizing-lemma-divide-by-torsion}, on a $H^0_T(M') = 0$, c'est-à-dire qu'aucun idéal premier associé à $M'$ n'appartient à $T$. Par conséquent, $\text{depth}(M'_\mathfrak p) \geq 1$ pour $\mathfrak p \in T$.
\end{proof}

\begin{lemma}
\label{lemma-get-depth-2-along-Z}
Soit $j : U \to X$ une immersion ouverte de schémas noethériens. Soit $\mathcal{F}$ un $\mathcal{O}_X$-module cohérent. Supposons $\mathcal{F}' = j_*(\mathcal{F}|_U)$ cohérent. Alors $\mathcal{F} \to \mathcal{F}'$ est l'unique application de $\mathcal{O}_X$-modules cohérents telle que
\begin{enumerate}
\item $\mathcal{F}|_U \to \mathcal{F}'|_U$ soit un isomorphisme,
\item $\text{depth}_{\mathcal{O}_{X, x}}(\mathcal{F}'_x) \geq 2$ pour $x \in X$, $x \not \in U$.
\end{enumerate}
Si $f : Y \to X$ est un morphisme plat, avec $Y$ noethérien, alors $f^*\mathcal{F} \to f^*\mathcal{F}'$ est l'application correspondante pour $f^{-1}(U) \subset Y$.
\end{lemma}

\begin{proof}
On a $\text{depth}_{\mathcal{O}_{X, x}}(\mathcal{F}'_x) \geq 2$ d'après Diviseurs, lemme \ref{divisors-lemma-depth-pushforward}, point (3). L'unicité de $\mathcal{F} \to \mathcal{F}'$ résulte de Diviseurs, lemme \ref{divisors-lemma-depth-2-hartog}. La compatibilité aux images réciproques plates résulte du changement de base plat ; voir Cohomologie des schémas, lemme \ref{coherent-lemma-flat-base-change-cohomology}.
\end{proof}

\begin{lemma}
\label{lemma-make-S2-along-Z}
Soit $X$ un schéma noethérien. Soit $Z \subset X$ un sous-schéma fermé. Soit $\mathcal{F}$ un $\mathcal{O}_X$-module cohérent. Supposons $X$ universellement caténaire et les fibres formelles de ses anneaux locaux vérifiant $(S_1)$. Alors il existe une unique application $\mathcal{F} \to \mathcal{F}''$ de $\mathcal{O}_X$-modules cohérents telle que
\begin{enumerate}
\item $\mathcal{F}_x \to \mathcal{F}''_x$ soit un isomorphisme pour $x \in X \setminus Z$,
\item $\mathcal{F}_x \to \mathcal{F}''_x$ soit surjective et $\text{depth}_{\mathcal{O}_{X, x}}(\mathcal{F}''_x) = 1$ pour $x \in Z$ tel qu'il existe une spécialisation immédiate $x' \leadsto x$ avec $x' \not \in Z$ et $x' \in \text{Ass}(\mathcal{F})$,
\item $\text{depth}_{\mathcal{O}_{X, x}}(\mathcal{F}''_x) \geq 2$ pour les autres $x \in Z$.
\end{enumerate}
Si $f : Y \to X$ est un morphisme de Cohen-Macaulay, avec $Y$ noethérien, alors $f^*\mathcal{F} \to f^*\mathcal{F}''$ vérifie les mêmes propriétés relativement à $f^{-1}(Z) \subset Y$.
\end{lemma}

\begin{proof}
Soit $\mathcal{F} \to \mathcal{F}'$ l'application construite au lemme \ref{lemma-get-depth-1-along-Z} pour la partie $Z$ de $X$. Rappelons que $\mathcal{F}'$ est le quotient de $\mathcal{F}$ par le sous-faisceau des sections à support dans $Z$.

\medskip\noindent
Démontrons d'abord l'unicité. Soit $\mathcal{F} \to \mathcal{F}''$ comme dans le lemme. On obtient une factorisation $\mathcal{F} \to \mathcal{F}' \to \mathcal{F}''$, puisque $\text{Ass}(\mathcal{F}'') \cap Z = \emptyset$ d'après (2) et (3). Soit $U \subset X$ le plus grand sous-schéma ouvert tel que $\mathcal{F}'|_U \to \mathcal{F}''|_U$ soit un isomorphisme. On voit que $U$ contient tous les points de (2). Alors Diviseurs, lemme \ref{divisors-lemma-depth-2-hartog}, permet de conclure que $\mathcal{F}'' = j_*(\mathcal{F}'|_U)$. On obtient ainsi l'unicité (petit détail : si l'on dispose de deux tels $\mathcal{F}''$, on prend l'intersection des ouverts $U$ obtenus pour chacun).

\medskip\noindent
Démonstration de l'existence. Rappelons que $\text{Ass}(\mathcal{F}') = \{x_1, \ldots, x_n\}$ est fini et que $x_i \not \in Z$. Soit $Y_i$ l'adhérence de $\{x_i\}$. Soient $Z_{i, j}$ les composantes irréductibles de $Z \cap Y_i$. Observons que $\text{Supp}(\mathcal{F}') \cap Z = \bigcup Z_{i, j}$. Soit $z_{i, j} \in Z_{i, j}$ le point générique. Posons
$$
d_{i, j} = \dim(\mathcal{O}_{\overline{\{x_i\}}, z_{i, j}})
$$
Si $d_{i, j} = 1$, alors $z_{i, j}$ est l'un des points de (2). Il n'est donc pas nécessaire de modifier $\mathcal{F}'$ en ces points. De plus, en supposant toujours $d_{i, j} = 1$, le lemme \ref{lemma-depth-function} permet de trouver un voisinage ouvert $z_{i, j} \in V_{i, j} \subset X$ tel que $\text{depth}_{\mathcal{O}_{X, z}}(\mathcal{F}'_z) \geq 2$ pour $z \in Z_{i, j} \cap V_{i, j}$, $z \not = z_{i, j}$. Posons
$$
Z' = X \setminus
\left(
X \setminus Z \cup \bigcup\nolimits_{d_{i, j} = 1} V_{i, j})
\right)
$$
Notons $j' : X \setminus Z' \to X$. Le choix de $Z'$ assure que les hypothèses du lemme \ref{lemma-finiteness-pushforward-general} sont satisfaites. On conclut en posant $\mathcal{F}'' = j'_*(\mathcal{F}'|_{X \setminus Z'})$ et en appliquant le lemme \ref{lemma-get-depth-2-along-Z}.

\medskip\noindent
La dernière assertion résulte de la formule de variation de la profondeur par un homomorphisme local plat ; voir Algèbre, lemme \ref{algebra-lemma-apply-grothendieck-module}, et de l'hypothèse sur les fibres de $f$ contenue dans le fait que $f$ est de Cohen-Macaulay. Détails omis.
\end{proof}

\begin{lemma}
\label{lemma-make-S2-along-T-simple}
Soit $X$ un schéma noethérien possédant localement un complexe dualisant. Soit $T' \subset X$ une partie stable par spécialisation. Soit $\mathcal{F}$ un $\mathcal{O}_X$-module cohérent. Supposons que, si $x \leadsto x'$ est une spécialisation immédiate de points de $X$ avec $x' \in T'$ et $x \not \in T'$, alors $\text{depth}(\mathcal{F}_x) \geq 1$. Il existe alors une unique application $\mathcal{F} \to \mathcal{F}''$ de $\mathcal{O}_X$-modules cohérents telle que
\begin{enumerate}
\item $\mathcal{F}_x \to \mathcal{F}''_x$ soit un isomorphisme pour $x \not \in T'$,
\item $\text{depth}_{\mathcal{O}_{X, x}}(\mathcal{F}''_x) \geq 2$ pour $x \in T'$.
\end{enumerate}
Si $f : Y \to X$ est un morphisme de Cohen-Macaulay, avec $Y$ noethérien, alors $f^*\mathcal{F} \to f^*\mathcal{F}''$ vérifie les mêmes propriétés relativement à $f^{-1}(T') \subset Y$.
\end{lemma}

\begin{proof}
Démonstration de l'unicité. Soit $\mathcal{F} \to \mathcal{F}''$ comme dans le lemme. Soit $U \subset X$ le plus grand sous-schéma ouvert tel que $\mathcal{F}|_U \to \mathcal{F}''|_U$ soit un isomorphisme. D'après (1), $U$ contient l'ensemble $X \setminus T'$. D'après (2) et Diviseurs, lemme \ref{divisors-lemma-depth-2-hartog}, on conclut que $\mathcal{F}'' = j_*(\mathcal{F}|_U)$, où $j : U \to X$ est le morphisme d'inclusion. On obtient ainsi l'unicité (petit détail : si l'on dispose de deux tels $\mathcal{F}''$, on prend l'intersection des ouverts $U$ obtenus pour chacun).

\medskip\noindent
Démonstration de l'existence. Soit $\mathcal{F} \to \mathcal{F}'$ le quotient de $\mathcal{F}$ construit au lemme \ref{lemma-get-depth-1-along-Z} à l'aide de $T'$. Rappelons que $\mathcal{F}'$ est le quotient de $\mathcal{F}$ par le sous-faisceau des sections à support dans $T'$. Définissons
$$
\mathcal{F}'' = \colim j_*(\mathcal{F}'|_V)
$$
où $j : V \to X$ parcourt les sous-schémas ouverts tels que $X \setminus V \subset T'$. Observons que la limite inductive est filtrante, puisque $T'$ est stable par spécialisation. Chacune des applications $\mathcal{F}' \to j_*(\mathcal{F}'|_V)$ est injective, puisque $\text{Ass}(\mathcal{F}')$ est disjoint de $T'$. Ainsi $\mathcal{F}' \to \mathcal{F}''$ est injective.

\medskip\noindent
Supposons $X = \Spec(A)$ affine et $\mathcal{F}$ correspondant au $A$-module de type fini $M$. Alors $\mathcal{F}'$ correspond à $M' = M / H^0_{T'}(M)$ ; voir la démonstration du lemme \ref{lemma-get-depth-1-along-Z}. En appliquant les lemmes \ref{lemma-local-cohomology} et \ref{lemma-adjoint-ext}, on voit que $\mathcal{F}''$ correspond à un $A$-module $M''$ qui s'insère dans la suite exacte courte
$$
0 \to M' \to M'' \to H^1_{T'}(M') \to 0
$$
D'après la proposition \ref{proposition-finiteness} et la condition de l'énoncé sur les spécialisations immédiates, $M''$ est un $A$-module de type fini. On voit ainsi que $\mathcal{F}''$ est cohérent.

\medskip\noindent
La cohérence de $\mathcal{F}''$ et l'injectivité des applications de transition dans la limite inductive filtrante ci-dessus impliquent que $\mathcal{F}'' = j_*(\mathcal{F}'|_V)$ pour tout $V$ assez petit comme ci-dessus. Les conditions (1) et (2) résultent alors du lemme \ref{lemma-get-depth-2-along-Z}.

\medskip\noindent
La dernière assertion résulte de la formule de variation de la profondeur par un homomorphisme local plat ; voir Algèbre, lemme \ref{algebra-lemma-apply-grothendieck-module}, et de l'hypothèse sur les fibres de $f$ contenue dans le fait que $f$ est de Cohen-Macaulay. Détails omis.
\end{proof}

\begin{lemma}
\label{lemma-make-S2-along-T}
Soit $X$ un schéma noethérien possédant localement un complexe dualisant. Soient $T' \subset T \subset X$ des parties stables par spécialisation telles que, si $x \leadsto x'$ est une spécialisation immédiate de points de $X$ et $x' \in T'$, alors $x \in T$. Soit $\mathcal{F}$ un $\mathcal{O}_X$-module cohérent. Il existe alors une unique application $\mathcal{F} \to \mathcal{F}''$ de $\mathcal{O}_X$-modules cohérents telle que
\begin{enumerate}
\item $\mathcal{F}_x \to \mathcal{F}''_x$ soit un isomorphisme pour $x \not \in T$,
\item $\mathcal{F}_x \to \mathcal{F}''_x$ soit surjective et $\text{depth}_{\mathcal{O}_{X, x}}(\mathcal{F}''_x) \geq 1$ pour $x \in T$, $x \not \in T'$,
\item $\text{depth}_{\mathcal{O}_{X, x}}(\mathcal{F}''_x) \geq 2$ pour $x \in T'$.
\end{enumerate}
Si $f : Y \to X$ est un morphisme de Cohen-Macaulay, avec $Y$ noethérien, alors $f^*\mathcal{F} \to f^*\mathcal{F}''$ vérifie les mêmes propriétés relativement à $f^{-1}(T') \subset f^{-1}(T) \subset Y$.
\end{lemma}

\begin{proof}
Soit d'abord $\mathcal{F} \to \mathcal{F}'$ le quotient de $\mathcal{F}$ construit au lemme \ref{lemma-get-depth-1-along-Z} à l'aide de $T$. Soit ensuite $\mathcal{F}' \to \mathcal{F}''$ l'unique application de modules cohérents construite au lemme \ref{lemma-make-S2-along-T-simple} à l'aide de $T'$. Alors $\mathcal{F} \to \mathcal{F}''$ possède les propriétés voulues.
\end{proof}










\section{Annulation de Hartshorne-Lichtenbaum}
\label{section-Hartshorne-Lichtenbaum-vanishing}

\noindent
Ce résultat d'annulation est l'analogue local du théorème de Lichtenbaum présenté dans Dualité pour les schémas, section \ref{duality-section-lichtenbaum}. On le trouvera, ainsi que bien d'autres résultats, dans \cite{CD}.

\begin{lemma}
\label{lemma-cd-top-vanishing}
Soit $A$ un anneau noethérien de dimension $d$. Soient $I \subset I' \subset A$ des idéaux. Si $I'$ est contenu dans le radical de Jacobson de $A$ et si $\text{cd}(A, I') < d$, alors $\text{cd}(A, I) < d$.
\end{lemma}

\begin{proof}
D'après le lemme \ref{lemma-cd-dimension}, on sait que $\text{cd}(A, I) \leq d$. Nous utiliserons le lemme \ref{lemma-isomorphism} pour montrer que
$$
H^d_{V(I')}(A) \to H^d_{V(I)}(A)
$$
est surjective, ce qui achèvera la démonstration. Choisissons $\mathfrak p \in V(I) \setminus V(I')$. L'hypothèse sur $I'$ montre que $\mathfrak p$ n'est pas un idéal maximal de $A$. Ainsi $\dim(A_\mathfrak p) < d$. Alors $H^d_{\mathfrak pA_\mathfrak p}(A_\mathfrak p) = 0$ d'après le lemme \ref{lemma-cd-dimension}.
\end{proof}

\begin{lemma}
\label{lemma-cd-top-vanishing-some-module}
Soit $A$ un anneau noethérien de dimension $d$. Soit $I \subset A$ un idéal. Si $H^d_{V(I)}(M) = 0$ pour un $A$-module de type fini dont le support contient toutes les composantes irréductibles de dimension $d$, alors $\text{cd}(A, I) < d$.
\end{lemma}

\begin{proof}
D'après le lemme \ref{lemma-cd-dimension}, on sait que $\text{cd}(A, I) \leq d$. Ainsi, pour tout $A$-module de type fini $N$, on a $H^i_{V(I)}(N) = 0$ pour $i > d$. Disons que la propriété $\mathcal{P}$ est vérifiée pour le $A$-module de type fini $N$ si $H^d_{V(I)}(N) = 0$. L'une des hypothèses est que $\mathcal{P}(M)$ est vérifiée. Observons que $\mathcal{P}(N_1 \oplus N_2)
\Leftrightarrow (\mathcal{P}(N_1) \wedge \mathcal{P}(N_2))$. Si $N \to N'$ est surjective, alors $\mathcal{P}(N) \Rightarrow \mathcal{P}(N')$, puisque $H^{d + 1}_{V(I)}$ est nul (voir ci-dessus). Soient $\mathfrak p_1, \ldots, \mathfrak p_n$ les idéaux premiers minimaux de $A$ tels que $\dim(A/\mathfrak p_i) = d$. Observons que $\mathcal{P}(N)$ est vérifiée si le support de $N$ est disjoint de $\{\mathfrak p_1, \ldots, \mathfrak p_n\}$, pour des raisons de dimension ; voir le lemme \ref{lemma-cd-dimension}. Pour chaque $i$, posons $M_i = M/\mathfrak p_i M$. C'est un $A$-module de type fini annulé par $\mathfrak p_i$, dont le support est égal à $V(\mathfrak p_i)$ (on utilise ici l'hypothèse sur le support de $M$). Enfin, si $J \subset A$ est un idéal, alors $\mathcal{P}(JM_i)$ est vérifiée, car $JM_i$ est un quotient d'une somme directe de copies de $M$. Il résulte donc de Cohomologie des schémas, lemme \ref{coherent-lemma-property-higher-rank-cohomological}, que $\mathcal{P}$ est vérifiée pour tout $A$-module de type fini.
\end{proof}

\begin{lemma}
\label{lemma-top-coh-divisible}
Soit $A$ un anneau local noethérien de dimension $d$. Soit $f \in A$ un élément n'appartenant à aucun idéal premier minimal de dimension $d$. Alors $f : H^d_{V(I)}(M) \to H^d_{V(I)}(M)$ est surjective pour tout $A$-module de type fini $M$ et tout idéal $I \subset A$.
\end{lemma}

\begin{proof}
Le support de $M/fM$ est de dimension $< d$, par l'hypothèse sur $f$. Ainsi $H^d_{V(I)}(M/fM) = 0$ d'après le lemme \ref{lemma-cd-dimension}. Par conséquent, $H^d_{V(I)}(fM) \to H^d_{V(I)}(M)$ est surjective. Puisque le lemme \ref{lemma-cd-dimension} donne $\text{cd}(A, I) \leq d$, on voit aussi que la surjection $M \to fM$, $x \mapsto fx$ induit une surjection $H^d_{V(I)}(M) \to H^d_{V(I)}(fM)$.
\end{proof}

\begin{lemma}
\label{lemma-cd-bound-dualizing}
Soit $A$ un anneau local noethérien muni d'un complexe dualisant normalisé $\omega_A^\bullet$. Soit $I \subset A$ un idéal. Si $H^0_{V(I)}(\omega_A^\bullet) = 0$, alors $\text{cd}(A, I) < \dim(A)$.
\end{lemma}

\begin{proof}
Posons $d = \dim(A)$. Soient $\mathfrak p_1, \ldots, \mathfrak p_n \subset A$ les idéaux premiers minimaux de dimension $d$. Rappelons que le $A$-module de type fini $H^{-i}(\omega_A^\bullet)$ est non nul seulement pour $i \in \{0, \ldots, d\}$ et que le support de $H^{-i}(\omega_A^\bullet)$ est de dimension $\leq i$ ; voir le lemme \ref{lemma-sitting-in-degrees}. Posons $\omega_A = H^{-d}(\omega_A^\bullet)$. Par évitement des idéaux premiers (Algèbre, lemme \ref{algebra-lemma-silly}), on peut trouver $f \in A$, $f \not \in \mathfrak p_i$, qui annule $H^{-i}(\omega_A^\bullet)$ pour $i < d$. Considérons le triangle distingué
$$
\omega_A[d] \to \omega_A^\bullet \to
\tau_{\geq -d + 1}\omega_A^\bullet \to \omega_A[d + 1]
$$
Voir Catégories dérivées, remarque \ref{derived-remark-truncation-distinguished-triangle}. D'après Catégories dérivées, lemme \ref{derived-lemma-trick-vanishing-composition}, $f^d$ induit l'endomorphisme nul de $\tau_{\geq -d + 1}\omega_A^\bullet$. Les axiomes des catégories triangulées fournissent une application
$$
\omega_A^\bullet \to \omega_A[d]
$$
dont le composé avec $\omega_A[d] \to \omega_A^\bullet$ est la multiplication par $f^d$ sur $\omega_A[d]$. On en déduit que $f^d$ annule $H^d_{V(I)}(\omega_A)$. Le lemme \ref{lemma-top-coh-divisible} donne alors $H^d_{V(I)}(\omega_A) = 0$. On conclut par le lemme \ref{lemma-cd-top-vanishing-some-module} et le fait que $(\omega_A)_{\mathfrak p_i}$ est non nul (voir par exemple Complexes dualisants, lemme \ref{dualizing-lemma-nonvanishing-generically-local}).
\end{proof}

\begin{lemma}
\label{lemma-inverse-system-symbolic-powers}
Soit $(A, \mathfrak m)$ un anneau local noethérien complet intègre. Soit $\mathfrak p \subset A$ un idéal premier de dimension $1$. Pour tout $n \geq 1$, il existe $m \geq n$ tel que $\mathfrak p^{(m)} \subset \mathfrak p^n$.
\end{lemma}

\begin{proof}
Rappelons que la puissance symbolique $\mathfrak p^{(m)}$ est définie comme le noyau de $A \to A_\mathfrak p/\mathfrak p^mA_\mathfrak p$. Puisque la localisation est exacte, on en déduit que, dans la suite exacte courte
$$
0 \to \mathfrak a_n \to A/\mathfrak p^n \to A/\mathfrak p^{(n)} \to 0
$$
le support de $\mathfrak a_n$ est contenu dans $\{\mathfrak m\}$. En particulier, le système projectif $(\mathfrak a_n)$ vérifie la condition de Mittag-Leffler, puisque chaque $\mathfrak a_n$ est un $A$-module artinien. Le lemme équivaut donc à la condition $\lim \mathfrak a_n = 0$. Soit $f \in \lim \mathfrak a_n$. Alors $f$ est un élément de $A = \lim A/\mathfrak p^n$ (on utilise ici la complétude de $A$) dont l'image dans le complété $A_\mathfrak p^\wedge$ de $A_\mathfrak p$ est nulle. Puisque $A_\mathfrak p \to A_\mathfrak p^\wedge$ est fidèlement plat, on voit que $f$ a une image nulle dans $A_\mathfrak p$. Puisque $A$ est intègre, $f$ est nul, comme voulu.
\end{proof}

\begin{proposition}
\label{proposition-Hartshorne-Lichtenbaum-vanishing}
\begin{reference}
\cite[Theorem 3.1]{CD}
\end{reference}
Soit $A$ un anneau local noethérien de complété $A^\wedge$. Soit $I \subset A$ un idéal tel que
$$
\dim V(IA^\wedge + \mathfrak p) \geq 1
$$
pour tout idéal premier minimal $\mathfrak p \subset A^\wedge$ de dimension $\dim(A)$. Alors $\text{cd}(A, I) < \dim(A)$.
\end{proposition}

\begin{proof}
Puisque $A \to A^\wedge$ est fidèlement plat, on a $H^d_{V(I)}(A) \otimes_A A^\wedge = H^d_{V(IA^\wedge)}(A^\wedge)$ d'après Complexes dualisants, lemme \ref{dualizing-lemma-torsion-change-rings}. On peut donc supposer $A$ complet.

\medskip\noindent
Supposons $A$ complet. Soient $\mathfrak p_1, \ldots, \mathfrak p_n \subset A$ les idéaux premiers minimaux de dimension $d$. Considérons l'anneau local complet $A_i = A/\mathfrak p_i$. On a $H^d_{V(I)}(A_i) = H^d_{V(IA_i)}(A_i)$ d'après Complexes dualisants, lemme \ref{dualizing-lemma-local-cohomology-and-restriction}. D'après le lemme \ref{lemma-cd-top-vanishing-some-module}, il suffit de prouver le lemme pour $(A_i, IA_i)$. On peut donc supposer que $A$ est un anneau local complet intègre.

\medskip\noindent
Supposons que $A$ soit un anneau local complet intègre. On peut choisir un idéal premier $\mathfrak p \supset I$ tel que $\dim(A/\mathfrak p) = 1$. D'après le lemme \ref{lemma-cd-top-vanishing}, il suffit de prouver le lemme pour $\mathfrak p$.

\medskip\noindent
D'après le lemme \ref{lemma-cd-bound-dualizing}, il suffit de montrer que $H^0_{V(\mathfrak p)}(\omega_A^\bullet) = 0$. Rappelons que
$$
H^0_{V(\mathfrak p)}(\omega_A^\bullet) =
\colim \text{Ext}^0_A(A/\mathfrak p^n, \omega_A^\bullet)
$$
D'après le lemme \ref{lemma-inverse-system-symbolic-powers}, cette limite inductive est égale à
$$
\colim \text{Ext}^0_A(A/\mathfrak p^{(n)}, \omega_A^\bullet)
$$
Puisque $\text{depth}(A/\mathfrak p^{(n)}) = 1$, ces groupes Ext sont nuls d'après le lemme \ref{lemma-sitting-in-degrees}, comme voulu.
\end{proof}

\begin{lemma}
\label{lemma-affine-complement}
Soit $(A, \mathfrak m)$ un anneau local noethérien. Soit $I \subset A$ un idéal. Supposons $A$ excellent, normal, et $\dim V(I) \geq 1$. Alors $\text{cd}(A, I) < \dim(A)$. En particulier, si $\dim(A) = 2$, alors $\Spec(A) \setminus V(I)$ est affine.
\end{lemma}

\begin{proof}
D'après Compléments d'algèbre, lemme \ref{more-algebra-lemma-completion-normal-local-ring}, le complété $A^\wedge$ est normal, donc intègre. L'hypothèse de la proposition \ref{proposition-Hartshorne-Lichtenbaum-vanishing} est donc vérifiée, et l'on conclut. L'assertion d'affinité résulte du lemme \ref{lemma-cd-is-one}.
\end{proof}
















\section{Action du Frobenius}
\label{section-frobenius}

\noindent
Soit $p$ un nombre premier. Soit $A$ un anneau tel que $p = 0$ dans $A$. L'{\it endomorphisme de Frobenius} de $A$ est l'application
$$
F : A \longrightarrow A,
\quad
a \longmapsto a^p
$$
Dans cette section, nous démontrons des lemmes sur les modules munis d'une action du Frobenius.

\begin{lemma}
\label{lemma-annihilator-frobenius-module}
Soit $p$ un nombre premier. Soit $(A, \mathfrak m, \kappa)$ un anneau local noethérien tel que $p = 0$ dans $A$. Soit $M$ un $A$-module de type fini tel que $M \otimes_{A, F} A \cong M$. Alors $M$ est libre de type fini.
\end{lemma}

\begin{proof}
Choisissons une présentation $A^{\oplus m} \to A^{\oplus n} \to M$ qui induit un isomorphisme $\kappa^{\oplus n} \to M/\mathfrak m M$. Soit $T = (a_{ij})$ la matrice de l'application $A^{\oplus m} \to A^{\oplus n}$. Observons que $a_{ij} \in \mathfrak m$. En effectuant le changement de base par $F$ et en utilisant son exactitude à droite, on obtient une présentation $A^{\oplus m} \to A^{\oplus n} \to M$ dont la matrice est $T = (a_{ij}^p)$. On a donc une présentation avec $a_{ij} \in \mathfrak m^p$. En répétant cette construction, on trouve, pour chaque $e \geq 1$, une présentation avec $a_{ij} \in \mathfrak m^e$. Cela implique que les idéaux de Fitting (Compléments d'algèbre, définition \ref{more-algebra-definition-fitting-ideal}) $\text{Fit}_k(M)$ pour $k < n$ sont contenus dans $\bigcap_{e \geq 1} \mathfrak m^e$. Ce dernier étant nul par le théorème d'intersection de Krull (Algèbre, lemme \ref{algebra-lemma-intersect-powers-ideal-module-zero}), on conclut que $M$ est libre de rang $n$, d'après Compléments d'algèbre, lemme \ref{more-algebra-lemma-fitting-ideal-finite-locally-free}.
\end{proof}

\noindent
Dans cette section, nous dirons que des éléments $f_1, \ldots, f_r$ d'un anneau $A$ sont {\it indépendants} si $\sum a_if_i = 0$ implique $a_i \in (f_1, \ldots, f_r)$. Autrement dit, en posant $I = (f_1, \ldots, f_r)$, le module $I/I^2$ est libre sur $A/I$, de base $f_1, \ldots, f_r$.

\begin{lemma}
\label{lemma-1}
\begin{reference}
Voir \cite{Lech-inequalities} et \cite[Lemma 1 page 299]{MatCA}.
\end{reference}
Soit $A$ un anneau. Si $f_1, \ldots, f_{r - 1}, f_rg_r$ sont indépendants, alors $f_1, \ldots, f_r$ sont indépendants.
\end{lemma}

\begin{proof}
{\emergencystretch=5em
Supposons $\sum a_if_i = 0$. Alors $\sum a_ig_rf_i = 0$. Ainsi $a_r \in (f_1, \ldots, f_{r - 1}, f_rg_r)$. Écrivons $a_r = \sum_{i < r} b_i f_i + b f_rg_r$. Alors $0 = \sum_{i < r} (a_i + b_if_r)f_i + bf_r^2g_r$. Par conséquent, $a_i + b_i f_r \in (f_1, \ldots, f_{r - 1}, f_rg_r)$, ce qui implique $a_i \in (f_1, \ldots, f_r)$, comme voulu.
\par}
\end{proof}

\begin{lemma}
\label{lemma-2}
\begin{reference}
Voir \cite{Lech-inequalities} et \cite[Lemma 2 page 300]{MatCA}.
\end{reference}
Soit $A$ un anneau. Si $f_1, \ldots, f_{r - 1}, f_rg_r$ sont indépendants et si le $A$-module $A/(f_1, \ldots, f_{r - 1}, f_rg_r)$ est de longueur finie, alors
\begin{align*}
& \text{length}_A(A/(f_1, \ldots, f_{r - 1}, f_rg_r)) \\
& =
\text{length}_A(A/(f_1, \ldots, f_{r - 1}, f_r)) +
\text{length}_A(A/(f_1, \ldots, f_{r - 1}, g_r))
\end{align*}
\end{lemma}

\begin{proof}
Montrons qu'il existe une suite exacte
$$
0 \to
A/(f_1, \ldots, f_{r - 1}, g_r) \xrightarrow{f_r}
A/(f_1, \ldots, f_{r - 1}, f_rg_r) \to
A/(f_1, \ldots, f_{r - 1}, f_r) \to 0
$$
En effet, si $a f_r \in (f_1, \ldots, f_{r - 1}, f_rg_r)$, alors $\sum_{i < r} a_i f_i + (a + bg_r)f_r = 0$ pour certains $b, a_i \in A$. Ainsi $\sum_{i < r} a_i g_r f_i + (a + bg_r)g_rf_r = 0$, ce qui implique $a + bg_r \in (f_1, \ldots, f_{r - 1}, f_rg_r)$ et signifie que $a$ a une image nulle dans $A/(f_1, \ldots, f_{r - 1}, g_r)$. Cela prouve l'affirmation. Pour conclure, utiliser l'additivité des longueurs (Algèbre, lemme \ref{algebra-lemma-length-additive}).
\end{proof}

\begin{lemma}
\label{lemma-3}
\begin{reference}
Voir \cite{Lech-inequalities} et \cite[Lemma 3 page 300]{MatCA}.
\end{reference}
Soit $(A, \mathfrak m)$ un anneau local. Si $\mathfrak m = (x_1, \ldots, x_r)$ et si $x_1^{e_1}, \ldots, x_r^{e_r}$ sont indépendants pour certains $e_i > 0$, alors $\text{length}_A(A/(x_1^{e_1}, \ldots, x_r^{e_r})) = e_1\ldots e_r$.
\end{lemma}

\begin{proof}
Utiliser les lemmes \ref{lemma-1} et \ref{lemma-2} et une récurrence.
\end{proof}

\begin{lemma}
\label{lemma-flat-extension-independent}
Soit $\varphi : A \to B$ un homomorphisme plat d'anneaux. Si $f_1, \ldots, f_r \in A$ sont indépendants, alors $\varphi(f_1), \ldots, \varphi(f_r) \in B$ sont indépendants.
\end{lemma}

\begin{proof}
Posons $I = (f_1, \ldots, f_r)$ et $J = \varphi(I)B$. Par platitude, on a $I/I^2 \otimes_A B = J/J^2$. La liberté de $I/I^2$ sur $A/I$ implique donc celle de $J/J^2$ sur $B/J$.
\end{proof}

\begin{lemma}[Kunz]
\label{lemma-frobenius-flat-regular}
\begin{reference}
\cite{Kunz-flat}
\end{reference}
Soit $p$ un nombre premier. Soit $A$ un anneau noethérien tel que $p = 0$. Les conditions suivantes sont équivalentes :
\begin{enumerate}
\item $A$ est régulier,
\item $F : A \to A$, $a \mapsto a^p$ est plat.
\end{enumerate}
\end{lemma}

\begin{proof}
Observons que $\Spec(F) : \Spec(A) \to \Spec(A)$ est l'application identique. La régularité se définit à l'aide des anneaux locaux et la platitude se vérifie sur les anneaux locaux ; voir Algèbre, lemme \ref{algebra-lemma-flat-localization}. On peut donc supposer, et l'on suppose, que $A$ est un anneau local noethérien d'idéal maximal $\mathfrak m$.

\medskip\noindent
Supposons $A$ régulier. Soit $x_1, \ldots, x_d$ un système de paramètres de $A$. En appliquant $F$, on obtient $F(x_1), \ldots, F(x_d) = x_1^p, \ldots, x_d^p$, qui est un système de paramètres de $A$. Ainsi $F$ est plat ; voir Algèbre, lemmes \ref{algebra-lemma-CM-over-regular-flat} et \ref{algebra-lemma-regular-ring-CM}.

\medskip\noindent
Réciproquement, supposons $F$ plat. Écrivons $\mathfrak m = (x_1, \ldots, x_r)$ avec $r$ minimal. Alors $x_1, \ldots, x_r$ sont indépendants au sens défini ci-dessus. Puisque $F$ est plat, $x_1^p, \ldots, x_r^p$ sont indépendants ; voir le lemme \ref{lemma-flat-extension-independent}. Ainsi $\text{length}_A(A/(x_1^p, \ldots, x_r^p)) = p^r$ d'après le lemme \ref{lemma-3}. Posons $\chi(n) = \text{length}_A(A/\mathfrak m^n)$ et rappelons qu'il s'agit d'un polynôme numérique de degré $\dim(A)$ ; voir Algèbre, proposition \ref{algebra-proposition-dimension}. Choisissons $n \gg 0$. Observons que
$$
\mathfrak m^{pn + pr} \subset F(\mathfrak m^n)A \subset \mathfrak m^{pn}
$$
comme on le voit en considérant les monômes en $x_1, \ldots, x_r$. On a
$$
A/F(\mathfrak m^n)A = A/\mathfrak m^n \otimes_{A, F} A
$$
Par platitude de $F$, ce module est de longueur $\chi(n) \text{length}_A(A/F(\mathfrak m)A)$ (Algèbre, lemme \ref{algebra-lemma-pullback-module}), qui vaut $p^r\chi(n)$ d'après ce qui précède. On en déduit
$$
\chi(pn + pr) \geq p^r\chi(n) \geq \chi(pn)
$$
La comparaison des termes dominants donne $r = \dim(A)$, c'est-à-dire que $A$ est régulier.
\end{proof}




\section{Structure de certains modules}
\label{section-structure}

\noindent
Quelques résultats sur la structure de certains types de modules sur les anneaux locaux réguliers. On trouvera de tels résultats et bien davantage dans \cite{Huneke-Sharp}, \cite{Lyubeznik}, \cite{Lyubeznik2}.

\begin{lemma}
\label{lemma-structure-torsion-D-module-regular}
\begin{reference}
Cas particulier de \cite[Theorem 2.4]{Lyubeznik}
\end{reference}
{\emergencystretch=3em
Soit $k$ un corps de caractéristique $0$. Soit $d \geq 1$. Soit $A = k[[x_1, \ldots, x_d]]$, d'idéal maximal $\mathfrak m$. Soit $M$ un module de torsion par rapport aux puissances de $\mathfrak m$, sur l'anneau $A$, muni d'opérateurs additifs $D_1, \ldots, D_d$ vérifiant la règle de Leibniz
\par}
$$
D_i(fz) = \partial_i(f) z + f D_i(z)
$$
pour $f \in A$ et $z \in M$. Ici $\partial_i$ désigne la dérivation par rapport à $x_i$. Alors $M$ est isomorphe à une somme directe de copies de l'enveloppe injective $E$ de $k$.
\end{lemma}

\begin{proof}
Choisissons un ensemble $J$ et un isomorphisme $M[\mathfrak m] \to \bigoplus_{j \in J} k$. Puisque $\bigoplus_{j \in J} E$ est injectif (Complexes dualisants, lemme \ref{dualizing-lemma-sum-injective-modules}), on peut prolonger cet isomorphisme en un homomorphisme de $A$-modules $\varphi : M \to \bigoplus_{j \in J} E$. Montrons que $\varphi$ est un isomorphisme, c'est-à-dire une bijection.

\medskip\noindent
Injectivité. Soit $z \in M$ non nul. Puisque $M$ est de torsion par rapport aux puissances de $\mathfrak m$, on peut choisir un élément $f \in A$ tel que $fz \in M[\mathfrak m]$ et $fz \not = 0$. Alors $\varphi(fz) = f\varphi(z)$ est non nul, donc $\varphi(z)$ est non nul.

\medskip\noindent
Surjectivité. Soit $z \in M$. Alors $x_1^n z = 0$ pour un certain $n \geq 0$. Montrons que $z \in x_1M$ par récurrence sur $n$. Si $n = 0$, alors $z = 0$, et le résultat est vrai. Si $n > 0$, l'application de $D_1$ donne $0 = n x_1^{n - 1} z + x_1^nD_1(z)$. Ainsi $x_1^{n - 1}(nz + x_1D_1(z)) = 0$. Par récurrence, on obtient $nz + x_1D_1(z) \in x_1M$. Puisque $n$ est inversible, on en déduit $z \in x_1M$. Ainsi $M$ est $x_1$-divisible. Si $\varphi$ n'est pas surjective, on peut choisir $e \in \bigoplus_{j \in J} E$ n'appartenant pas à $M$. En raisonnant comme ci-dessus, on peut supposer $\mathfrak m e \subset M$, en particulier $x_1 e \in M$. Il existe un élément $z_1 \in M$ tel que $x_1 z_1 = x_1 e$. D'où $x_1(z_1 - e) = 0$. En remplaçant $e$ par $e - z_1$, on peut supposer $e$ annulé par $x_1$. Il suffit donc de montrer que
$$
\varphi[x_1] :
M[x_1]
\longrightarrow
\left(\bigoplus\nolimits_{j \in J} E\right)[x_1] =
\bigoplus\nolimits_{j \in J} E[x_1]
$$
est surjective. Si $d = 1$, c'est vrai par construction de $\varphi$. Si $d > 1$, on observe que $E[x_1]$ est l'enveloppe injective du corps résiduel de $k[[x_2, \ldots, x_d]]$ ; voir Complexes dualisants, lemme \ref{dualizing-lemma-quotient}. Observons que $M[x_1]$, en tant que module sur $k[[x_2, \ldots, x_d]]$, est de torsion par rapport aux puissances de $\mathfrak m/(x_1)$ et est muni d'opérateurs $D_2, \ldots, D_d$ vérifiant la règle de Leibniz ci-dessus. Par récurrence sur $d$, on en déduit que $\varphi[x_1]$ est surjective, comme voulu.
\end{proof}

\begin{lemma}
\label{lemma-structure-torsion-Frobenius-regular}
\begin{reference}
Résulte de \cite[Corollary 3.6]{Huneke-Sharp} moyennant quelques arguments supplémentaires. Résulte également directement de \cite[Theorem 1.4]{Lyubeznik2}.
\end{reference}
Soit $p$ un nombre premier. Soit $(A, \mathfrak m, k)$ un anneau local régulier tel que $p = 0$. Notons $F : A \to A$, $a \mapsto a^p$ l'endomorphisme de Frobenius. Soit $M$ un module de torsion par rapport aux puissances de $\mathfrak m$ tel que $M \otimes_{A, F} A \cong M$. Alors $M$ est isomorphe à une somme directe de copies de l'enveloppe injective $E$ de $k$.
\end{lemma}

\begin{proof}
Choisissons un ensemble $J$ et un homomorphisme de $A$-modules $\varphi : M \to \bigoplus_{j \in J} E$ qui envoie $M[\mathfrak m]$ isomorphiquement sur $(\bigoplus_{j \in J} E)[\mathfrak m] = \bigoplus_{j \in J} k$. Montrons que $\varphi$ est un isomorphisme, c'est-à-dire une bijection.

\medskip\noindent
Injectivité. Soit $z \in M$ non nul. Puisque $M$ est de torsion par rapport aux puissances de $\mathfrak m$, on peut choisir un élément $f \in A$ tel que $fz \in M[\mathfrak m]$ et $fz \not = 0$. Alors $\varphi(fz) = f\varphi(z)$ est non nul, donc $\varphi(z)$ est non nul.

\medskip\noindent
Surjectivité. Rappelons que $F$ est plat ; voir le lemme \ref{lemma-frobenius-flat-regular}. Soit $x_1, \ldots, x_d$ un système minimal de générateurs de $\mathfrak m$. Notons
$$
M_n = M[x_1^{p^n}, \ldots, x_d^{p^n}]
$$
le sous-module de $M$ constitué des éléments annulés par $x_1^{p^n}, \ldots, x_d^{p^n}$. Ainsi $M_0 = M[\mathfrak m]$ est un espace vectoriel sur $k$. De plus, $M = \bigcup M_n$, puisque $M$ est de torsion par rapport aux puissances de $\mathfrak m$. Puisque $F^n$ est plat et que $F^n(x_i) = x_i^{p^n}$, on a
\begingroup\small
$$
M_n \cong (M \otimes_{A, F^n} A)[x_1^{p^n}, \ldots, x_d^{p^n}] =
M[x_1, \ldots, x_d] \otimes_{A, F^n} A =
M_0 \otimes_k A/(x_1^{p^n}, \ldots, x_d^{p^n})
$$
\endgroup
Ainsi $M_n$ est libre sur $A/(x_1^{p^n}, \ldots, x_d^{p^n})$. Un calcul montre que tout élément de $A/(x_1^{p^n}, \ldots, x_d^{p^n})$ annulé par $x_1^{p^n - 1}$ est divisible par $x_1$ ; par exemple, on peut utiliser le fait que $A/(x_1^{p^n}, \ldots, x_d^{p^n}) \cong
k[x_1, \ldots, x_d]/(x_1^{p^n}, \ldots, x_d^{p^n})$ d'après Algèbre, lemme \ref{algebra-lemma-regular-complete-containing-coefficient-field}. Il en est donc de même pour tout élément de $M_n$. Puisque tout élément de $M$ appartient à $M_n$ pour tout $n \gg 0$ et que tout élément de $M$ est annulé par une puissance de $x_1$, on conclut que $M$ est $x_1$-divisible.

\medskip\noindent
Posons $x = x_1$. Nous avons vu ci-dessus que $M$ est $x$-divisible. Si $\varphi$ n'est pas surjective, on peut choisir $e \in \bigoplus_{j \in J} E$ n'appartenant pas à $M$. En raisonnant comme ci-dessus, on peut supposer $\mathfrak m e \subset M$, en particulier $x e \in M$. Il existe un élément $z_1 \in M$ tel que $x z_1 = x e$. D'où $x(z_1 - e) = 0$. En remplaçant $e$ par $e - z_1$, on peut supposer $e$ annulé par $x$. Il suffit donc de montrer que
$$
\varphi[x] :
M[x]
\longrightarrow
\left(\bigoplus\nolimits_{j \in J} E\right)[x] =
\bigoplus\nolimits_{j \in J} E[x]
$$
est surjective. Si $d = 1$, c'est vrai par construction de $\varphi$. Si $d > 1$, on observe que $E[x]$ est l'enveloppe injective du corps résiduel de l'anneau régulier $A/xA$ ; voir Complexes dualisants, lemme \ref{dualizing-lemma-quotient}. Observons que $M[x]$, en tant que module sur $A/xA$, est de torsion par rapport aux puissances de $\mathfrak m/(x)$ et que l'on a
\begin{align*}
M[x] \otimes_{A/xA, F} A/xA
& = M[x] \otimes_{A, F} A \otimes_A A/xA \\
& = (M \otimes_{A, F} A)[x^p] \otimes_A A/xA \\
& \cong M[x^p] \otimes_A A/xA
\end{align*}
Raisonner à l'aide de la platitude de $F$, comme précédemment. Montrons que $M[x^p] \otimes_A A/xA \to M[x]$, $z \otimes 1 \mapsto x^{p - 1}z$ est un isomorphisme. Il suffit de le prouver pour chacun des modules $M_n$, $n > 0$ définis ci-dessus ; cela résulte alors du même résultat pour $A/(x_1^{p^n}, \ldots, x_d^{p^n})$ et $x = x_1$. Par récurrence sur $\dim(A)$, on en déduit que $\varphi[x]$ est surjective, comme voulu.
\end{proof}





\section{Structures supplémentaires sur la cohomologie locale}
\label{section-additional}

\noindent
Voici un exemple de résultat.

\begin{lemma}
\label{lemma-derivation}
Soit $A$ un anneau. Soit $I \subset A$ un idéal de type fini. Posons $Z = V(I)$. Pour toute dérivation $\theta : A \to A$, il existe un opérateur additif canonique $D$ sur les modules de cohomologie locale $H^i_Z(A)$, vérifiant la règle de Leibniz relativement à $\theta$.
\end{lemma}

\begin{proof}
{\emergencystretch=3em
Soient $f_1, \ldots, f_r$ des éléments engendrant $I$. Rappelons que $R\Gamma_Z(A)$ est calculé par le complexe
\par}
$$
A \to \prod\nolimits_{i_0} A_{f_{i_0}} \to
\prod\nolimits_{i_0 < i_1} A_{f_{i_0}f_{i_1}}
\to \ldots \to A_{f_1\ldots f_r}
$$
Voir Complexes dualisants, lemme \ref{dualizing-lemma-local-cohomology-adjoint}. Puisque $\theta$ se prolonge de manière unique en un opérateur additif sur tout localisé de $A$, vérifiant la règle de Leibniz relativement à $\theta$, le lemme est clair.
\end{proof}

\begin{lemma}
\label{lemma-frobenius}
Soit $p$ un nombre premier. Soit $A$ un anneau tel que $p = 0$. Notons $F : A \to A$, $a \mapsto a^p$ l'endomorphisme de Frobenius. Soit $I \subset A$ un idéal de type fini. Posons $Z = V(I)$. Il existe un isomorphisme $R\Gamma_Z(A) \otimes_{A, F}^\mathbf{L} A \cong R\Gamma_Z(A)$.
\end{lemma}

\begin{proof}
Cela résulte de Complexes dualisants, lemme \ref{dualizing-lemma-torsion-change-rings}, et du fait que $Z = V(f_1^p, \ldots, f_r^p)$ si $I = (f_1, \ldots, f_r)$.
\end{proof}

\begin{lemma}
\label{lemma-etale-derivation}
Soit $A$ un anneau. Soit $V \to \Spec(A)$ quasi-compact, quasi-séparé et \'etale. Pour toute dérivation $\theta : A \to A$, il existe un opérateur additif canonique $D$ sur $H^i(V, \mathcal{O}_V)$ vérifiant la règle de Leibniz relativement à $\theta$.
\end{lemma}

\begin{proof}
Si $V$ est séparé, on peut raisonner à l'aide d'un recouvrement ouvert affine $V = \bigcup_{j = 1, \ldots m} V_j$. En effet, puisque $V$ est séparé, on peut écrire $V_{j_0 \ldots j_p} = \Spec(B_{j_0 \ldots j_p})$. Voir Schémas, lemme \ref{schemes-lemma-characterize-separated}. On trouve alors que le $A$-module $H^i(V, \mathcal{O}_V)$ est le groupe de cohomologie en degré $i$ du complexe de {\v C}ech
$$
\prod B_{j_0} \to
\prod B_{j_0j_1} \to
\prod B_{j_0j_1j_2} \to \ldots
$$
Voir Cohomologie des schémas, lemme \ref{coherent-lemma-cech-cohomology-quasi-coherent}. Chaque $B = B_{j_0 \ldots j_p}$ est une algèbre \'etale sur $A$. Ainsi $\Omega_B = \Omega_A \otimes_A B$, et l'on en déduit que $\theta$ se prolonge de manière unique en une dérivation $\theta_B : B \to B$. Ces applications définissent un endomorphisme du complexe de {\v C}ech et les opérateurs voulus sur les groupes de cohomologie.

\medskip\noindent
Dans le cas général, on utilise un hyperrecouvrement de $V$ par des ouverts affines, exactement comme dans la première partie de la démonstration de Cohomologie des schémas, lemme \ref{coherent-lemma-hypercoverings}. Nous omettons les détails.
\end{proof}

\begin{remark}
\label{remark-higher-order-operators}
On peut renforcer les lemmes \ref{lemma-derivation} et \ref{lemma-etale-derivation} pour y inclure les opérateurs différentiels d'ordre supérieur, à l'aide d'Algèbre, lemme \ref{algebra-lemma-formally-etale-principal-parts} et remarque \ref{algebra-remark-formally-etale-differential-operators}. Si le besoin s'en fait sentir, nous énoncerons et démontrerons ici un lemme précis.
\end{remark}

\begin{lemma}
\label{lemma-etale-frobenius}
Soit $p$ un nombre premier. Soit $A$ un anneau tel que $p = 0$. Notons $F : A \to A$, $a \mapsto a^p$ l'endomorphisme de Frobenius. Si $V \to \Spec(A)$ est quasi-compact, quasi-séparé et \'etale, alors il existe un isomorphisme $R\Gamma(V, \mathcal{O}_V) \otimes_{A, F}^\mathbf{L} A \cong
R\Gamma(V, \mathcal{O}_V)$.
\end{lemma}

\begin{proof}
Observons que le morphisme de Frobenius relatif
$$
V \longrightarrow V \times_{\Spec(A), \Spec(F)} \Spec(A)
$$
de $V$ sur $A$ est un isomorphisme ; voir Morphismes \'etales, lemme \ref{etale-lemma-relative-frobenius-etale}. Le lemme résulte donc de la cohomologie et du changement de base ; voir Catégories dérivées des schémas, lemme \ref{perfect-lemma-compare-base-change}. Observons que, puisque $V$ est \'etale sur $A$, il est plat sur $A$.
\end{proof}





\section{Un peu d'uniformité, I}
\label{section-uniform-vanishing}

\noindent
Le principal objectif de cette section est d'énoncer et de démontrer le lemme \ref{lemma-characterize-vanishing-tor-ext-above-e}.

\begin{lemma}
\label{lemma-map-tor-1-zero}
Soit $R$ un anneau. Soit $M \to M'$ un homomorphisme de $R$-modules, avec $M$ de présentation finie, tel que $\text{Tor}_1^R(M, N) \to \text{Tor}_1^R(M', N)$ soit nul pour tous les $R$-modules $N$. Alors $M \to M'$ se factorise par un $R$-module libre.
\end{lemma}

\begin{proof}
On peut choisir un morphisme de suites exactes courtes
$$
\xymatrix{
0 \ar[r] &
K \ar[r] \ar[d] &
R^{\oplus r} \ar[r] \ar[d] &
M \ar[r] \ar[d] &
0 \\
0 \ar[r] &
K' \ar[r] &
\bigoplus_{i \in I} R \ar[r] &
M' \ar[r] &
0
}
$$
dont la flèche verticale de droite est l'application donnée. Ce morphisme se factorise par la suite exacte courte
\begin{equation}
\label{equation-pushout}
0 \to K' \to E \to M \to 0
\end{equation}
qui est la somme amalgamée de la première suite exacte courte par $K \to K'$. Une chasse au diagramme montre que l'hypothèse du lemme implique que l'homomorphisme de bord $\text{Tor}_1^R(M, N) \to K' \otimes_R N$ induit par (\ref{equation-pushout}) est nul, c'est-à-dire que la suite (\ref{equation-pushout}) est universellement exacte. Algèbre, lemme \ref{algebra-lemma-universally-exact-split}, implique alors que (\ref{equation-pushout}) est scindée (c'est ici que l'on utilise le fait que $M$ est de présentation finie). Ainsi l'application $M \to M'$ se factorise par $\bigoplus_{i \in I} R$, ce qui conclut.
\end{proof}

\begin{lemma}
\label{lemma-characterize-vanishing-tor-ext-above-e}
Soit $R$ un anneau. Soit $\alpha : M \to M'$ un homomorphisme de $R$-modules. Soient $P_\bullet \to M$ et $P'_\bullet \to M'$ des résolutions par des $R$-modules projectifs. Soit $e \geq 0$ un entier. Considérons les conditions suivantes :
\begin{enumerate}
\item Il existe un morphisme de complexes $a_\bullet : P_\bullet \to P'_\bullet$ induisant $\alpha$ en cohomologie, avec $a_i = 0$ pour $i > e$.
\item Il existe un morphisme de complexes $a_\bullet : P_\bullet \to P'_\bullet$ induisant $\alpha$ en cohomologie, avec $a_{e + 1} = 0$.
\item L'application $\Ext^i_R(M', N) \to \Ext^i_R(M, N)$ est nulle pour tous les $R$-modules $N$ et tout $i > e$.
\item L'application $\Ext^{e + 1}_R(M', N) \to \Ext^{e + 1}_R(M, N)$ est nulle pour tous les $R$-modules $N$.
\item Soit $N = \Im(P'_{e + 1} \to P'_e)$, et notons $\xi \in \Ext^{e + 1}_R(M', N)$ l'élément canonique (voir la démonstration). Alors $\xi$ a une image nulle dans $\Ext^{e + 1}_R(M, N)$.
\item L'application $\text{Tor}_i^R(M, N) \to \text{Tor}_i^R(M', N)$ est nulle pour tous les $R$-modules $N$ et tout $i > e$.
\item L'application $\text{Tor}_{e + 1}^R(M, N) \to \text{Tor}_{e + 1}^R(M', N)$ est nulle pour tous les $R$-modules $N$.
\end{enumerate}
On a toujours les implications
$$
(1) \Leftrightarrow (2) \Leftrightarrow (3) \Leftrightarrow (4)
\Leftrightarrow (5) \Rightarrow (6) \Leftrightarrow (7)
$$
Si $M$ est $(-e - 1)$-pseudo-cohérent (par exemple si $R$ est noethérien et $M$ est un $R$-module de type fini), alors toutes les conditions sont équivalentes.
\end{lemma}

\begin{proof}
Il est clair que (2) implique (1). Si $a_\bullet$ est comme dans (1), on peut considérer le morphisme de complexes $a'_\bullet : P_\bullet \to P'_\bullet$ défini par $a'_i = a_i$ pour $i \leq e + 1$ et $a'_i = 0$ pour $i \geq e + 1$, afin d'obtenir un morphisme de complexes comme dans (2). Ainsi (1) équivaut à (2).

\medskip\noindent
La construction des foncteurs $\Ext$ et $\text{Tor}$ à l'aide de résolutions (Algèbre, sections \ref{algebra-section-ext} et \ref{algebra-section-tor}) montre que (1) et (2) impliquent toutes les autres conditions.

\medskip\noindent
Il est clair que (3) implique (4), qui implique (5). Soit $N$ comme dans (5). L'application canonique $\tilde \xi : P'_{e + 1} \to N$ précomposée avec $P'_{e + 2} \to P'_{e + 1}$ est nulle. On peut donc considérer la classe $\xi$ de $\tilde \xi$ dans
$$
\Ext^{e + 1}_R(M', N) =
\frac{\Ker(\Hom(P'_{e + 1}, N \to \Hom(P'_{e + 2}, N)}{
\Im(\Hom(P'_e, N \to \Hom(P'_{e + 1}, N)}
$$
Choisissons un morphisme de complexes $a_\bullet : P_\bullet \to P'_\bullet$ relevant $\alpha$ ; voir Catégories dérivées, lemme \ref{derived-lemma-morphisms-lift-projective}. Si $\xi$ a une image nulle dans $\Ext^{e + 1}_R(M', N)$, on trouve alors une application $\varphi : P_e \to N$ telle que $\tilde \xi  \circ a_{e + 1} = \varphi \circ d$. On obtient ainsi un morphisme de complexes
$$
\xymatrix{
\ldots \ar[r] &
P_{e + 1} \ar[r] \ar[d]^0 &
P_e \ar[r] \ar[d]^{a_e - \varphi} &
P_{e - 1} \ar[r] \ar[d]^{a_{e - 1}} &
\ldots \\
\ldots \ar[r] &
P'_{e + 1} \ar[r] &
P'_e \ar[r] &
P'_{e - 1} \ar[r] &
\ldots
}
$$
comme dans (2). Ainsi (1) -- (5) sont équivalentes.

\medskip\noindent
L'équivalence de (6) et (7) résulte du décalage de dimension ; nous omettons les détails.

\medskip\noindent
Supposons $M$ $(-e - 1)$-pseudo-cohérent. (L'assertion entre parenthèses du lemme résulte de Compléments d'algèbre, lemme \ref{more-algebra-lemma-Noetherian-pseudo-coherent}.) Montrons que (7) implique (4), ce qui achèvera la démonstration. Nous raisonnons par récurrence sur $e$. Le cas initial est $e = 0$. Alors $M$ est de présentation finie d'après Compléments d'algèbre, lemme \ref{more-algebra-lemma-n-pseudo-module}, et le lemme \ref{lemma-map-tor-1-zero} montre que $M \to M'$ se factorise par un module libre. Bien entendu, si $M \to M'$ se factorise par un module libre, alors $\Ext^i_R(M', N) \to \Ext^i_R(M, N)$ est nulle pour tout $i > 0$, comme voulu. Supposons $e > 0$. On peut choisir un morphisme de suites exactes courtes
$$
\xymatrix{
0 \ar[r] &
K \ar[r] \ar[d] &
R^{\oplus r} \ar[r] \ar[d] &
M \ar[r] \ar[d] &
0 \\
0 \ar[r] &
K' \ar[r] &
\bigoplus_{i \in I} R \ar[r] &
M' \ar[r] &
0
}
$$
{\emergencystretch=4em
dont la flèche verticale de droite est l'application donnée. On obtient $\text{Tor}_{i + 1}^R(M, N) = \text{Tor}^R_i(K, N)$ et $\Ext^{i + 1}_R(M, N) = \Ext^i_R(K, N)$ pour $i \geq 1$ et pour tous les $R$-modules $N$, et de même pour $M', K'$. Ainsi $\text{Tor}_e^R(K, N) \to \text{Tor}_e^R(K', N)$ est nulle pour tous les $R$-modules $N$. D'après Compléments d'algèbre, lemme \ref{more-algebra-lemma-cone-pseudo-coherent}, $K$ est $(-e)$-pseudo-cohérent. Par récurrence, on conclut que $\Ext^e(K', N) \to \Ext^e(K, N)$ est nulle pour tous les $R$-modules $N$, ce qui donne le résultat voulu.
\par}
\end{proof}

\begin{lemma}
\label{lemma-cd-sequence-Koszul}
Soit $I$ un idéal d'un anneau noethérien $A$. Pour tout $n \geq 1$, il existe $m > n$ tel que l'application $A/I^m \to A/I^n$ vérifie les conditions équivalentes du lemme \ref{lemma-characterize-vanishing-tor-ext-above-e} avec $e = \text{cd}(A, I)$.
\end{lemma}

\begin{proof}
{\emergencystretch=4em
Soit $\xi \in \Ext^{e + 1}_A(A/I^n, N)$ l'élément construit au lemme \ref{lemma-characterize-vanishing-tor-ext-above-e}, point (5). Puisque $e = \text{cd}(A, I)$, on a $0 = H^{e + 1}_Z(N) = H^{e + 1}_I(N) = \colim \Ext^{e + 1}(A/I^m, N)$ d'après Complexes dualisants, lemmes \ref{dualizing-lemma-local-cohomology-noetherian} et \ref{dualizing-lemma-local-cohomology-ext}. On peut donc choisir $m \geq n$ tel que $\xi$ ait une image nulle dans $\Ext^{e + 1}_A(A/I^m, N)$, comme voulu.
\par}
\end{proof}






\section{Un peu d'uniformité, II}
\label{section-uniform-vanishing-bis}

\noindent
Soit $I$ un idéal d'un anneau noethérien $A$. Soit $M$ un $A$-module de type fini. Soit $i > 0$. D'après Compléments d'algèbre, lemme \ref{more-algebra-lemma-tor-strictly-pro-zero}, il existe $c = c(A, I, M, i)$ tel que $\text{Tor}^A_i(M, A/I^n) \to \text{Tor}^A_i(M, A/I^{n - c})$ soit nulle pour tout $n \geq c$. Dans cette section, nous présentons des résultats montrant que l'on peut parfois choisir une constante $c$ convenant simultanément à tous les $A$-modules $M$ (et à une plage d'indices $i$). Ce sujet est lié au théorème d'Artin-Rees uniforme, étudié dans \cite{Huneke-uniform} et \cite{AHS}.

\medskip\noindent
Dans la remarque \ref{remark-strict-pro-isomorphism}, nous l'appliquerons pour montrer que divers systèmes projectifs liés à la complétion dérivée sont (ou ne sont pas) strictement pro-isomorphes.

\medskip\noindent
Le lemme suivant peut être sensiblement renforcé.

\begin{lemma}
\label{lemma-maps-zero-fixed-torsion}
Soit $I$ un idéal d'un anneau noethérien $A$. Pour tous $m \geq 0$ et $i > 0$, il existe $c = c(A, I, m, i) \geq 0$ tel que, pour tout $A$-module $M$ annulé par $I^m$, l'application
$$
\text{Tor}^A_i(M, A/I^n) \to \text{Tor}^A_i(M, A/I^{n - c})
$$
soit nulle pour tout $n \geq c$.
\end{lemma}

\begin{proof}
Raisonnons par récurrence sur $i$. Cas initial $i = 1$. La suite exacte courte $0 \to I^n \to A \to A/I^n \to 0$ définit une injection $\text{Tor}_1^A(M, A/I^n) \subset I^n \otimes_A M$ ; voir Algèbre, remarque \ref{algebra-remark-Tor-ring-mod-ideal}. Puisque $M$ est annulé par $I^m$, l'application $I^n \otimes_A M \to I^{n - m} \otimes_A M$ est nulle pour $n \geq m$. Le résultat est donc vrai avec $c = m$.

\medskip\noindent
Étape de récurrence. Soit $i > 1$ et supposons que $c$ convienne pour $i - 1$. D'après Compléments d'algèbre, lemme \ref{more-algebra-lemma-tor-strictly-pro-zero}, appliqué à $M = A/I^m$, on peut choisir $c' \geq 0$ tel que $\text{Tor}_i(A/I^m, A/I^n) \to \text{Tor}_i(A/I^m, A/I^{n - c'})$ soit nulle pour $n \geq c'$. Soit $M$ annulé par $I^m$. Choisissons une suite exacte courte
$$
0 \to S \to \bigoplus\nolimits_{i \in I} A/I^m \to M \to 0
$$
La suite exacte longue de Tor correspondante fournit une suite exacte
$$
\text{Tor}_i^A(\bigoplus\nolimits_{i \in I} A/I^m, A/I^n) \to
\text{Tor}_i^A(M, A/I^n) \to
\text{Tor}_{i - 1}^A(S, A/I^n)
$$
{\emergencystretch=5em
pour tous les entiers $n \geq 0$. Si $n \geq c + c'$, alors l'application $\text{Tor}_{i - 1}^A(S, A/I^n) \to \text{Tor}_{i - 1}^A(S, A/I^{n - c})$ est nulle et l'application $\text{Tor}_i^A(A/I^m, A/I^{n - c}) \to 
\text{Tor}_i^A(A/I^m, A/I^{n - c - c'})$ est nulle. En combinant cela avec les suites exactes courtes, on voit que le résultat est vrai pour $i$ avec la constante $c + c'$.
\par}
\end{proof}

\begin{lemma}
\label{lemma-annihilates-affine}
Soit $I = (a_1, \ldots, a_t)$ un idéal d'un anneau noethérien $A$. Posons $a = a_1$ et notons $B = A[\frac{I}{a}]$ l'algèbre d'éclatement affine. Il existe $c > 0$ tel que $\text{Tor}_i^A(B, M)$ soit annulé par $I^c$ pour tous les $A$-modules $M$ et tout $i \geq t$.
\end{lemma}

\begin{proof}
{\emergencystretch=4em
Rappelons que $B$ est le quotient de $A[x_2, \ldots, x_t]/(a_1x_2 - a_2, \ldots, a_1x_t - a_t)$ par sa $a_1$-torsion ; voir Algèbre, lemme \ref{algebra-lemma-affine-blowup-quotient-description}. Soit
\par}
$$
B_\bullet = \text{complexe de Koszul sur }a_1x_2 - a_2, \ldots, a_1x_t - a_t
\text{ sur }A[x_2, \ldots, x_t]
$$
considéré comme complexe de chaînes concentré en degrés $(t - 1), \ldots, 0$. Le complexe $B_\bullet[1/a_1]$ est isomorphe au complexe de Koszul sur $x_2 - a_2/a_1, \ldots, x_t - a_t/a_1$, qui est une suite régulière dans $A[1/a_1][x_2, \ldots, x_t]$. Puisque les suites régulières sont Koszul-régulières, on en déduit que l'augmentation
$$
\epsilon : B_\bullet \longrightarrow B
$$
est un quasi-isomorphisme après inversion de $a_1$. Puisque les modules d'homologie du cône $C_\bullet$ de $\epsilon$ sont des $A[x_2, \ldots, x_n]$-modules de type fini et que $C_\bullet$ est borné, il existe $c \geq 0$ tel que $a_1^c$ les annule tous. D'après Catégories dérivées, lemme \ref{derived-lemma-trick-vanishing-composition}, cela implique, après remplacement éventuel de $c$ par un entier plus grand, que $a_1^c$ est nul sur $C_\bullet$ dans $D(A)$. La démonstration s'achève en considérant le triangle distingué
$$
B_\bullet \otimes_A^\mathbf{L} M \to
B \otimes_A^\mathbf{L} M \to
C_\bullet \otimes_A^\mathbf{L} M
$$
En effet, le premier terme est représenté par $B_\bullet \otimes_A M$, concentré en degrés homologiques $(t - 1), \ldots, 0$, puisque les termes du complexe de Koszul $B_\bullet$ sont des $A$-modules libres (donc plats). Ainsi $\text{Tor}_i^A(B, M) = H_i(C_\bullet \otimes_A^\mathbf{L} M)$ pour $i > t - 1$, et ce module est annulé par $a_1^c$. Puisque $a_1^cB = I^cB$ et que le module Tor est un module sur $B$, on conclut.
\end{proof}

\noindent
Pour le reste de cette section, fixons un anneau noethérien $A$ et un idéal $I \subset A$. Notons
$$
p : X \to \Spec(A)
$$
l'éclatement de $\Spec(A)$ en l'idéal $I$. Autrement dit, $X$ est le $\text{Proj}$ de l'algèbre de Rees $\bigoplus_{n \geq 0} I^n$. D'après Cohomologie des schémas, lemmes \ref{coherent-lemma-coherent-on-proj} et \ref{coherent-lemma-recover-tail-graded-module}, on peut choisir un entier $q(A, I) \geq 0$ tel que, pour tout $q \geq q(A, I)$, on ait $H^i(X, \mathcal{O}_X(q)) = 0$ pour $i > 0$ et $H^0(X, \mathcal{O}_X(q)) = I^q$.

\begin{lemma}
\label{lemma-compute-tor-Iq}
Dans la situation ci-dessus, pour $q \geq q(A, I)$ et tout $A$-module $M$, on a
$$
R\Gamma(X, Lp^*\widetilde{M}(q)) \cong M \otimes_A^\mathbf{L} I^q
$$
dans $D(A)$.
\end{lemma}

\begin{proof}
Choisissons une résolution libre $F_\bullet \to M$. Alors $\widetilde{F}_\bullet$ est une résolution plate de $\widetilde{M}$. Ainsi $Lp^*\widetilde{M}$ est donné par le complexe $p^*\widetilde{F}_\bullet$. Donc $Lp^*\widetilde{M}(q)$ est donné par le complexe $p^*\widetilde{F}_\bullet(q)$. Puisque les $p^*\widetilde{F}_i(q)$ sont acycliques à droite pour $\Gamma(X, -)$, par le choix de $q \geq q(A, I)$, et que l'on a $\Gamma(X, p^*\widetilde{F}_i(q)) = I^qF_i$, par le choix de $q \geq q(A, I)$, on obtient que $R\Gamma(X, Lp^*\widetilde{M}(q))$ est donné par le complexe de termes $I^qF_i$, d'après Catégories dérivées des schémas, lemme \ref{perfect-lemma-acyclicity-lemma-global}. Le résultat s'ensuit, car le complexe $I^qF_\bullet$ calcule $M \otimes_A^\mathbf{L} I^q$ par définition.
\end{proof}

\begin{lemma}
\label{lemma-annihilates}
Dans la situation ci-dessus, soit $t$ une borne supérieure du nombre de générateurs de $I$. Il existe un entier $c = c(A, I) \geq 0$ tel que, pour tout $A$-module $M$, les faisceaux de cohomologie $H^j(Lp^*\widetilde{M})$ soient annulés par $I^c$ pour $j \leq -t$.
\end{lemma}

\begin{proof}
Écrivons $I = (a_1, \ldots, a_t)$. La question est locale sur les ouverts affines de $X$. Pour $1 \leq i \leq t$, soit $B_i = A[\frac{I}{a_i}]$ l'algèbre d'éclatement affine. Alors $X$ possède un recouvrement ouvert affine par les spectres des anneaux $B_i$ ; voir Diviseurs, lemme \ref{divisors-lemma-blowing-up-affine}. D'après la description de l'image réciproque dérivée donnée dans Catégories dérivées des schémas, lemme \ref{perfect-lemma-quasi-coherence-pullback}, il suffit de montrer que, pour chaque $i$, il existe $c \geq 0$ tel que
$$
\text{Tor}_j^A(B_i, M)
$$
soit annulé par $I^c$ pour $j \geq t$. C'est le lemme \ref{lemma-annihilates-affine}.
\end{proof}

\begin{lemma}
\label{lemma-annihilates-tors}
Dans la situation ci-dessus, soit $t$ une borne supérieure du nombre de générateurs de $I$. Il existe un entier $c = c(A, I) \geq 0$ tel que, pour tout $A$-module $M$, les modules Tor $\text{Tor}_i^A(M, A/I^q)$ soient annulés par $I^c$ pour $i > t$ et tout $q \geq 0$.
\end{lemma}

\begin{proof}
Soit $q(A, I)$ comme ci-dessus. Pour $q \geq q(A, I)$, on a
$$
R\Gamma(X, Lp^*\widetilde{M}(q)) = M \otimes_A^\mathbf{L} I^q
$$
d'après le lemme \ref{lemma-compute-tor-Iq}. On dispose d'une suite spectrale bornée et convergente
$$
H^a(X, H^b(Lp^*\widetilde{M}(q))) \Rightarrow
\text{Tor}_{-a - b}^A(M, I^q)
$$
d'après Catégories dérivées des schémas, lemme \ref{perfect-lemma-spectral-sequence}. Soit $d$ un entier comme dans Cohomologie des schémas, lemme \ref{coherent-lemma-vanishing-nr-affines-quasi-separated} (on peut en fait prendre $d = t$ ; voir Cohomologie des schémas, lemme \ref{coherent-lemma-vanishing-nr-affines}). On voit alors que $H^{-i}(X, Lp^*\widetilde{M}(q)) = \text{Tor}_i^A(M, I^q)$ possède une filtration finie d'au plus $d$ étapes, dont les gradués sont des sous-quotients des modules
$$
H^a(X, H^{- i - a}(Lp^*\widetilde{M})(q)),\quad
a = 0, 1, \ldots, d - 1
$$
Si $i \geq t$, tous ces modules sont annulés par $I^c$, où $c = c(A, I)$ est comme dans le lemme \ref{lemma-annihilates}, car les faisceaux de cohomologie $H^{- i - a}(Lp^*\widetilde{M})$ sont tous annulés par $I^c$ d'après ce lemme. Ainsi $\text{Tor}_i^A(M, I^q)$ est annulé par $I^{dc}$ pour $q \geq q(A, I)$ et $i \geq t$. La suite exacte courte $0 \to I^q \to A \to A/I^q \to 0$ montre alors que $\text{Tor}_i(M, A/I^q)$ est annulé par $I^{dc}$ pour $q \geq q(A, I)$ et $i > t$. On conclut que $I^m$, avec $m = \max(dc, q(A, I) - 1)$, annule $\text{Tor}_i^A(M, A/I^q)$ pour tout $q \geq 0$ et tout $i > t$, comme voulu.
\end{proof}

\begin{lemma}
\label{lemma-tor-maps-vanish}
Soit $I$ un idéal d'un anneau noethérien $A$. Soit $t \geq 0$ une borne supérieure du nombre de générateurs de $I$. Il existe $N, c \geq 0$ tels que les applications
$$
\text{Tor}_{t + 1}^A(M, A/I^n) \to \text{Tor}_{t + 1}^A(M, A/I^{n - c})
$$
soient nulles pour tout $A$-module $M$ et tout $n \geq N$.
\end{lemma}

\begin{proof}
Soit $c_1$ la constante obtenue au lemme \ref{lemma-annihilates-tors}. Rappelons que cette constante $c_1$ convient simultanément pour $\text{Tor}_i$ pour tous les $i > t$.

\medskip\noindent
Écrivons $I = (a_1, \ldots, a_t)$. Pour un $A$-module $M$, posons
$$
\ell(M) =
\#\{i \mid 1 \leq i \leq t,\ a_i^{c_1}\text{ est nul sur }M\}
$$
C'est un élément de $\{0, 1, \ldots, t\}$. Démontrons par récurrence descendante sur $0 \leq s \leq t$ l'assertion suivante $H_s$ : il existe $N, c \geq 0$ tels que, pour tout module $M$ avec $\ell(M) \geq s$, les applications
$$
\text{Tor}_{t + 1 + i}^A(M, A/I^n) \to \text{Tor}_{t + 1 + i}^A(M, A/I^{n - c})
$$
soient nulles pour $i = 0, \ldots, s$ et tout $n \geq N$.

\medskip\noindent
Cas initial : $s = t$. Si $\ell(M) = t$, alors $M$ est annulé par $(a_1^{c_1}, \ldots, a_t^{c_1}\}$ et donc par $I^{t(c_1 - 1) + 1}$. Le lemme \ref{lemma-maps-zero-fixed-torsion} montre que $H_t$ est vérifiée en prenant pour $c = N$ le maximum des entiers $c(A, I, t(c_1 - 1) + 1, t + 1), \ldots, c(A, I, t(c_1 - 1) + 1, 2t + 1)$ obtenus dans ce lemme.

\medskip\noindent
Étape de récurrence. Soit $0 \leq s < t$ et supposons donnés $N, c$ comme dans $H_{s + 1}$. Considérons un module $M$ tel que $\ell(M) = s$. On peut alors choisir $i$ tel que $a_i^{c_1}$ soit non nul sur $M$. Il s'ensuit que $\ell(M[a_i^c]) \geq s + 1$ et $\ell(M/a_i^{c_1}M) \geq s + 1$, et l'hypothèse de récurrence s'applique à ces modules. Considérons la suite exacte
$$
0 \to M[a_i^{c_1}] \to M \xrightarrow{a_i^{c_1}} M \to M/a_i^{c_1}M \to 0
$$
Notons $E \subset M$ l'image de la flèche centrale. Considérons le diagramme de modules Tor correspondant
$$
\xymatrix{
&
&
\text{Tor}_{i + 1}(M/a_i^{c_1}M, A/I^q) \ar[d] \\
\text{Tor}_i(M[a_i^{c_1}], A/I^q) \ar[r] &
\text{Tor}_i(M, A/I^q) \ar[r] \ar[rd]^0 &
\text{Tor}_i(E, A/I^q) \ar[d] \\
&
&
\text{Tor}_i(M, A/I^q)
}
$$
à lignes et colonnes exactes (pour tout $q$). La flèche dirigée vers le sud-est est nulle par le choix de $c_1$. On en déduit que le module $\text{Tor}_i(M, A/I^q)$ s'insère dans une suite exacte entre un quotient de $\text{Tor}_i(M[a_i^{c_1}], A/I^q)$ et un sous-module de $\text{Tor}_{i + 1}(M/a_i^{c_1}M, A/I^q)$ dans le diagramme. On conclut donc que $H_s$ est vérifiée en remplaçant $N$ par $N + c$ et $c$ par $2c$. Nous omettons certains détails.
\end{proof}

\begin{proposition}
\label{proposition-uniform-artin-rees}
Soit $I$ un idéal d'un anneau noethérien $A$. Soit $t \geq 0$ une borne supérieure du nombre de générateurs de $I$. Il existe $N, c \geq 0$ tels que, pour $n \geq N$, les applications
$$
A/I^n \to A/I^{n - c}
$$
vérifient les conditions équivalentes du lemme \ref{lemma-characterize-vanishing-tor-ext-above-e} avec $e = t$.
\end{proposition}

\begin{proof}
Conséquence immédiate des lemmes \ref{lemma-tor-maps-vanish} et \ref{lemma-characterize-vanishing-tor-ext-above-e}.
\end{proof}

\begin{remark}
\label{remark-better-bound}
L'article \cite{AHS} montre, parmi bien d'autres résultats, que si $A$ est local, la proposition \ref{proposition-uniform-artin-rees} reste vraie en remplaçant $e = t$ par $e = \dim(A)$. Le lemme \ref{lemma-cd-sequence-Koszul} conduit naturellement à se demander si la proposition \ref{proposition-uniform-artin-rees} reste vraie en remplaçant $e = t$ par $e = \text{cd}(A, I)$. Nous ne le savons pas.
\end{remark}

\begin{remark}
\label{remark-strict-pro-isomorphism}
{\emergencystretch=3em
Soit $I$ un idéal d'un anneau noethérien $A$. Écrivons $I = (f_1, \ldots, f_r)$. Notons $K_n^\bullet$ le complexe de Koszul sur $f_1^n, \ldots, f_r^n$, comme dans Compléments d'algèbre, situation \ref{more-algebra-situation-koszul}, et notons $K_n \in D(A)$ l'objet correspondant. Soit $M^\bullet$ un complexe borné de $A$-modules de type fini et notons $M \in D(A)$ l'objet correspondant. Considérons les systèmes projectifs suivants dans $D(A)$ :
\par}
\begin{enumerate}
\item $M^\bullet/I^nM^\bullet$, c'est-à-dire le complexe de termes $M^i/I^nM^i$,
\item $M \otimes_A^\mathbf{L} A/I^n$,
\item $M \otimes_A^\mathbf{L} K_n$,
\item $M \otimes_P^\mathbf{L} P/J^n$ (voir ci-dessous).
\end{enumerate}
Tous ces systèmes projectifs sont isomorphes en tant que pro-objets : l'isomorphisme entre (2) et (3) résulte de Compléments d'algèbre, lemme \ref{more-algebra-lemma-sequence-Koszul-complexes}. L'isomorphisme entre (1) et (2) est donné dans Compléments d'algèbre, lemme \ref{more-algebra-lemma-derived-completion-plain-completion}. Pour le dernier, voir ci-dessous.

\medskip\noindent
On peut toutefois se demander si ces isomorphismes de systèmes projectifs sont « stricts » ; cette terminologie et cette question se rattachent à la discussion de \cite[pages 61, 62]{quillenhomology}. Plus précisément, étant donnée une catégorie $\mathcal{C}$, on peut définir une « pro-catégorie stricte » dont les objets sont les systèmes projectifs $(X_n)$ et dont les morphismes $(X_n) \to (Y_n)$ sont donnés par des familles $(c, \varphi_n)$ constituées de $c \geq 0$ et de morphismes $\varphi_n : X_n \to Y_{n - c}$ pour tout $n \geq c$, satisfaisant à une condition évidente de compatibilité, à une certaine équivalence près (donnée essentiellement par l'augmentation de $c$). On demande alors si les systèmes projectifs ci-dessus sont isomorphes dans cette pro-catégorie stricte.

\medskip\noindent
Cela ne peut manifestement pas être le cas pour (1) et (3), même lorsque $M = A[0]$. En effet, le système $H^0(K_n) = A/(f_1^n, \ldots, f_r^n)$ n'est en général pas strictement pro-isomorphe, dans la catégorie des modules, au système $A/I^n$. Par exemple, si l'on prend $A = \mathbf{Z}[x_1, \ldots, x_r]$ et $f_i = x_i$, alors $H^0(K_n)$ n'est pas annulé par $I^{r(n - 1)}$.\footnote{Bien entendu, on peut demander si ces systèmes projectifs sont isomorphes dans une catégorie dont les objets sont les systèmes projectifs et dont les applications sont données par des familles $(r, c, \varphi_n)$ constituées de $r \geq 1$, de $c \geq 0$ et d'applications $\varphi_n : X_{rn} \to Y_{n - c}$ pour $n \geq c$.}

\medskip\noindent
{\emergencystretch=3em
Il se trouve que les résultats précédents montrent que l'application naturelle de (2) vers (1), étudiée dans Compléments d'algèbre, lemme \ref{more-algebra-lemma-derived-completion-plain-completion}, est un pro-isomorphisme strict. Esquissons la démonstration. Des arguments usuels faisant intervenir les troncations bêtes permettent d'abord de se ramener au cas où $M^\bullet$ est donné par un seul $A$-module de type fini $M$ placé en degré $0$. Choisissons $N, c \geq 0$ comme dans la proposition \ref{proposition-uniform-artin-rees}. La proposition implique que, pour $n \geq N$, on obtient des factorisations
\par}
$$
M \otimes_A^\mathbf{L} A/I^n
\to
\tau_{\geq -t}(M \otimes_A^\mathbf{L} A/I^n)
\to
M \otimes_A^\mathbf{L} A/I^{n - c}
$$
des applications de transition du système (2). D'autre part, d'après Compléments d'algèbre, lemme \ref{more-algebra-lemma-tor-strictly-pro-zero}, on peut trouver une autre constante $c' = c'(M) \geq 0$ telle que les applications $\text{Tor}_i^A(M, A/I^{n'}) \to \text{Tor}_i(M, A/I^{n' - c'})$ soient nulles pour $i = 1, 2, \ldots, t$ et $n' \geq c'$. Il résulte alors de Catégories dérivées, lemme \ref{derived-lemma-trick-vanishing-composition}, que l'application
$$
\tau_{\geq -t}(M \otimes_A^\mathbf{L} A/I^{n + tc'})
\to
\tau_{\geq -t}(M \otimes_A^\mathbf{L} A/I^n)
$$
se factorise par $M \otimes_A^\mathbf{L}A/I^{n + tc'} \to M/I^{n + tc'}M$. En combinant cela avec le résultat précédent, on obtient une factorisation
$$
M \otimes_A^\mathbf{L}A/I^{n + tc'} \to M/I^{n + tc'}M
\to M \otimes_A^\mathbf{L} A/I^{n - c}
$$
qui donne ce que l'on cherche. Si nous avons besoin de ce résultat, nous l'énoncerons soigneusement et en donnerons une démonstration détaillée.

\medskip\noindent
Pour (4), supposons donnés un anneau noethérien $P$, un homomorphisme d'anneaux $P \to A$ et un idéal $J \subset P$ tels que $I = JA$. D'après Compléments d'algèbre, section \ref{more-algebra-section-derived-base-change}, on obtient un foncteur $M \otimes_P^\mathbf{L} - : D(P) \to D(A)$ et un système projectif $M \otimes_P^\mathbf{L} P/J^n$ dans $D(A)$, comme dans (4). Si $P$ est noethérien, le système (4) est pro-isomorphe au système (1), par comparaison avec les complexes de Koszul. Si $P \to A$ est fini, le système (4) est strictement pro-isomorphe au système (2), car le système projectif $A \otimes_P^\mathbf{L} P/J^n$ est strictement pro-isomorphe au système projectif $A/I^n$ (d'après la discussion précédente) et l'on a
$$
M \otimes_P^\mathbf{L} P/J^n = M \otimes_A^\mathbf{L}
(A \otimes_P^\mathbf{L} P/J^n)
$$
d'après Compléments d'algèbre, lemme \ref{more-algebra-lemma-derived-base-change}.

\medskip\noindent
Un exemple usuel dans (4) consiste à prendre $P = \mathbf{Z}[x_1, \ldots, x_r]$, l'application $P \to A$ envoyant $x_i$ sur $f_i$, et $J = (x_1, \ldots, x_r)$. Dans ce cas, on montre que
$$
M \otimes_P^\mathbf{L} P/J^n =
M \otimes_{A[x_1, \ldots, x_r]}^\mathbf{L}
A[x_1, \ldots, x_r]/(x_1, \ldots, x_r)^n
$$
et l'on se ramène à l'un des cas étudiés ci-dessus (ce cas est toutefois strictement plus simple, car $A[x_1, \ldots, x_r]/(x_1, \ldots, x_r)^n$ est de dimension Tor au plus $r$ pour tout $n$ ; on peut donc éviter l'étape utilisant la proposition \ref{proposition-uniform-artin-rees}). Ce cas est étudié dans la démonstration de \cite[Proposition 3.5.1]{BS}.
\end{remark}








\section{Un peu d'uniformité, III}
\label{section-uniform}

\noindent
Dans cette section, fixons un anneau noethérien $A$ et un idéal $I \subset A$. Notre objectif est de démontrer le lemme \ref{lemma-bound-two-term-complex}, que nous utiliserons dans un chapitre ultérieur pour résoudre un problème de relèvement ; voir Algébrisation des espaces formels, lemme \ref{restricted-lemma-get-morphism-general-better}.

\medskip\noindent
Dans toute cette section, nous notons
$$
p : X \to \Spec(A)
$$
l'éclatement de $\Spec(A)$ en l'idéal $I$. Autrement dit, $X$ est le $\text{Proj}$ de l'algèbre de Rees $\bigoplus_{n \geq 0} I^n$. Nous considérons aussi le produit fibré
$$
\xymatrix{
Y \ar[r] \ar[d] & X \ar[d]^p \\
\Spec(A/I) \ar[r] & \Spec(A)
}
$$
Alors $Y$ est le diviseur exceptionnel de l'éclatement, donc un diviseur de Cartier effectif sur $X$ tel que $\mathcal{O}_X(-1) = \mathcal{O}_X(Y)$. Puisque la formation de $\text{Proj}$ commute au changement de base, on a
$$
Y = \text{Proj}(\bigoplus\nolimits_{n \geq 0} I^n/I^{n + 1}) = \text{Proj}(S)
$$
où $S = \text{Gr}_I(A) = \bigoplus_{n \geq 0} I^n/I^{n + 1}$.

\medskip\noindent
Nous notons $d = d(S) = d(\text{Gr}_I(A)) = d(\bigoplus_{n \geq 0} I^n/I^{n + 1})$ le maximum des dimensions des fibres de $p$ (et nous lui attribuons la valeur $0$ si $X = \emptyset$). Cette définition a bien un sens. En effet, on a
\begin{enumerate}
\item $d \leq t - 1$ si $I = (a_1, \ldots, a_t)$, puisqu'alors $X \subset \mathbf{P}^{t - 1}_A$,
\item $d$ est aussi la dimension maximale des fibres de $\text{Proj}(S) \to \Spec(S_0)$ lorsque $Y$ est non vide, et $d = 0$ si $Y = \emptyset$ (ou, de manière équivalente, $S = 0$, ou encore $I = A$).
\end{enumerate}
Ainsi $d$ ne dépend que de la classe d'isomorphisme de $S = \text{Gr}_I(A)$. Observons que $H^i(X, \mathcal{F}) = 0$ pour tout $\mathcal{O}_X$-module cohérent $\mathcal{F}$ et tout $i > d$, d'après Cohomologie des schémas, lemmes \ref{coherent-lemma-higher-direct-images-zero-above-dimension-fibre} et \ref{coherent-lemma-quasi-coherence-higher-direct-images-application}. Bien entendu, il en est de même pour les modules cohérents sur $Y$.

\medskip\noindent
Nous notons $q = q(S) = q(\text{Gr}_I(A)) = q(\bigoplus_{n \geq 0} I^n/I^{n + 1})$ l'entier défini comme suit. Notons que l'algèbre $S = \bigoplus_{n \geq 0} I^n/I^{n + 1}$ est un anneau gradué noethérien engendré en degré $1$ sur le degré $0$. Ainsi, d'après Cohomologie des schémas, lemmes \ref{coherent-lemma-coherent-on-proj} et \ref{coherent-lemma-recover-tail-graded-module}, on peut définir $q(S)$ comme le plus petit entier $q(S) \geq 0$ tel que, pour tout $q \geq q(S)$, on ait $H^i(Y, \mathcal{O}_Y(q)) = 0$ pour $1 \leq i \leq d$ et $H^0(Y, \mathcal{O}_Y(q)) = I^q/I^{q + 1}$. (Si $S = 0$, alors $q(S) = 0$.)

\medskip\noindent
Pour $n \geq 1$, on peut considérer le diviseur de Cartier effectif $nY$, que nous noterons $Y_n$.

\begin{lemma}
\label{lemma-bound-q-and-d}
Avec $q_0 = q(S)$ et $d = d(S)$ comme ci-dessus, on a
\begin{enumerate}
\item pour $n \geq 1$, $q \geq q_0$ et $i > 0$, on a $H^i(X, \mathcal{O}_{Y_n}(q)) = 0$,
\item pour $n \geq 1$ et $q \geq q_0$, on a $H^0(X, \mathcal{O}_{Y_n}(q)) = I^q/I^{q + n}$,
\item pour $q \geq q_0$ et $i > 0$, on a $H^i(X, \mathcal{O}_X(q)) = 0$,
\item pour $q \geq q_0$, on a $H^0(X, \mathcal{O}_X(q)) = I^q$.
\end{enumerate}
\end{lemma}

\begin{proof}
Si $I = A$, alors $X$ est affine et les assertions sont triviales. On peut donc supposer, et l'on suppose, $I \not = A$. Ainsi $Y$ et $X$ sont des schémas non vides.

\medskip\noindent
Démontrons (1) et (2) par récurrence sur $n$. Le cas initial $n = 1$ est la définition de $q_0$, puisque $Y_1 = Y$. Rappelons que $\mathcal{O}_X(1) = \mathcal{O}_X(-Y)$. On a donc une suite exacte courte
$$
0 \to \mathcal{O}_{Y_n}(1) \to \mathcal{O}_{Y_{n + 1}} \to \mathcal{O}_Y \to 0
$$
Ainsi, pour $i > 0$, on trouve
$$
H^i(X, \mathcal{O}_{Y_n}(q + 1)) \to
H^i(X, \mathcal{O}_{Y_{n + 1}}(q)) \to
H^i(X, \mathcal{O}_{Y}(q))
$$
et l'annulation donnée des termes extrêmes entraîne l'annulation voulue du terme central. Pour $i = 0$, on obtient un diagramme commutatif
$$
\xymatrix{
0 \ar[r] &
I^{q + 1}/I^{q + 1 + n} \ar[d] \ar[r] &
I^q/I^{q + 1 + n} \ar[d] \ar[r] &
I^q/I^{q + 1} \ar[d] \ar[r] &
0 \\
0 \ar[r] &
H^0(X, \mathcal{O}_{Y_n}(q + 1)) \ar[r] &
H^0(X, \mathcal{O}_{Y_{n + 1}}(q)) \ar[r] &
H^0(Y, \mathcal{O}_Y(q)) \ar[r] &
0
}
$$
à lignes exactes pour $q \geq q_0$ (pour la ligne du bas, observer que le terme suivant de la suite exacte longue de cohomologie est nul pour $q \geq q_0$). Puisque $q \geq q_0$, les flèches verticales de gauche et de droite sont des isomorphismes ; celle du milieu l'est donc aussi.

\medskip\noindent
Nous omettons les démonstrations de (3) et (4), qui sont analogues. On peut en fait déduire (3) et (4) de (1) et (2) à l'aide du théorème des fonctions formelles, mais ce serait un outil disproportionné.
\end{proof}

\noindent
Introduisons une notation : étant donnés $n \geq c \geq 0$, un {\it $(A, n, c)$-module} est un $A$-module de type fini $M$ annulé par $I^n$ et qui, en tant que $A/I^n$-module, est $I^c/I^n$-projectif ; voir Compléments d'algèbre, section \ref{more-algebra-section-near-projective}.

\medskip\noindent
Nous ferons l'abus de notation suivant : étant donné un $A$-module $M$, nous noterons $p^*M$ le module quasi-cohérent obtenu par image réciproque par $p$ du module quasi-cohérent $\widetilde{M}$ sur $\Spec(A)$ associé à $M$. On a par exemple $\mathcal{O}_{Y_n} = p^*(A/I^n)$. Pour une suite exacte courte $0 \to K \to L \to M \to 0$ de $A$-modules, on obtient une suite exacte
$$
p^*K \to p^*L \to p^*M \to 0
$$
puisque $\widetilde{\ }$ est un foncteur exact et $p^*$ un foncteur exact à droite.

\begin{lemma}
\label{lemma-almost-exactness}
Soit $0 \to K \to L \to M \to 0$ une suite exacte courte de $A$-modules telle que $K$ et $L$ soient annulés par $I^n$ et que $M$ soit un $(A, n, c)$-module. Alors le noyau de $p^*K \to p^*L$ est à support schématique dans $Y_c$.
\end{lemma}

\begin{proof}
Soit $\Spec(B) \subset X$ un ouvert affine. La restriction de la suite exacte à $\Spec(B)$ correspond à la suite de $B$-modules
$$
K \otimes_A B \to L \otimes_A B \to M \otimes_A B \to 0
$$
qui est isomorphe à la suite
$$
K \otimes_{A/I^n} B/I^nB \to
L \otimes_{A/I^n} B/I^nB \to
M \otimes_{A/I^n} B/I^nB \to 0
$$
{\emergencystretch=5em
Le noyau de la première application est donc l'image du module $\text{Tor}_1^{A/I^n}(M, B/I^nB)$. Rappelons que le diviseur exceptionnel $Y$ est défini par $I\mathcal{O}_X$. Il suffit donc de montrer que $\text{Tor}_1^{A/I^n}(M, B/I^nB)$ est annulé par $I^c$. Puisque la multiplication par $a \in I^c$ sur $M$ se factorise par un $A/I^n$-module libre de type fini, cela est clair.
\par}
\end{proof}

\noindent
On dispose de l'application canonique $\mathcal{O}_X \to \mathcal{O}_X(1)$, qui s'annule exactement le long de $Y$. Ainsi, pour tout $\mathcal{O}_X$-module cohérent $\mathcal{F}$, on a toujours des applications canoniques $\mathcal{F}(q) \to \mathcal{F}(q + n)$ pour tous $q \in \mathbf{Z}$ et $n \geq 0$.

\begin{lemma}
\label{lemma-annihilated}
Soit $\mathcal{F}$ un $\mathcal{O}_X$-module cohérent. Alors $\mathcal{F}$ est à support schématique dans $Y_c$ si et seulement si l'application canonique $\mathcal{F} \to \mathcal{F}(c)$ est nulle.
\end{lemma}

\begin{proof}
Cela vient du fait que $\mathcal{O}_X \to \mathcal{O}_X(1)$ s'annule exactement le long de $Y$.
\end{proof}

\begin{lemma}
\label{lemma-vanishing-coh-almost-projective}
Avec $q_0 = q(S)$ et $d = d(S)$ comme ci-dessus, supposons donnés des entiers $n \geq c \geq 0$, un $(A, n, c)$-module $M$, un indice $i \in \{0, 1, \ldots, d\}$ et un entier $q$. Distinguons les cas suivants :
\begin{enumerate}
\item Dans le cas $i = d \geq 1$ et $q \geq q_0$, on a $H^d(X, p^*M(q)) = 0$.
\item Dans le cas $i = d - 1 \geq 1$ et $q \geq q_0$, on a $H^{d - 1}(X, p^*M(q)) = 0$.
\item Dans le cas $d - 1 > i > 0$ et $q \geq q_0 + (d - 1 - i)c$, l'application $H^i(X, p^*M(q)) \to H^i(X, p^*M(q - (d - 1 - i)c))$ est nulle.
\item Dans le cas $i = 0$, $d \in \{0, 1\}$ et $q \geq q_0$, il existe une surjection
$$
I^qM \longrightarrow H^0(X, p^*M(q))
$$
\item Dans le cas $i = 0$, $d > 1$ et $q \geq q_0 + (d - 1)c$, l'application
$$
H^0(X, p^*M(q)) \to H^0(X, p^*M(q - (d - 1)c))
$$
a son image contenue dans celle de l'application canonique $I^{q - (d - 1)c}M \to H^0(X, p^*M(q - (d - 1)c))$.
\end{enumerate}
\end{lemma}

\begin{proof}
Soit $M$ un $(A, n, c)$-module. Choisissons une suite exacte courte
$$
0 \to K \to (A/I^n)^{\oplus r} \to M \to 0
$$
Nous utiliserons ci-dessous le fait que $K$ est un $(A, n, c)$-module ; voir Compléments d'algèbre, lemme \ref{more-algebra-lemma-ses-near-projective}. Considérons la suite exacte correspondante
$$
p^*K \to (\mathcal{O}_{Y_n})^{\oplus r} \to p^*M \to 0
$$
Décomposons-la en suites exactes courtes
$$
0 \to \mathcal{F} \to p^*K \to \mathcal{G} \to 0
\quad\text{et}\quad
0 \to \mathcal{G} \to (\mathcal{O}_{Y_n})^{\oplus r} \to p^*M \to 0
$$
D'après le lemme \ref{lemma-almost-exactness}, le module cohérent $\mathcal{F}$ est à support schématique dans $Y_c$.

\medskip\noindent
Démonstration de (1). Supposons $d > 0$. Il faut prouver que $H^d(X, p^*M(q)) = 0$ pour $q \geq q_0$. Par l'annulation de la cohomologie des tordus de $\mathcal{G}$ en degrés $> d$ et la suite exacte longue de cohomologie associée à la seconde suite exacte courte ci-dessus, il suffit de prouver que $H^d(X, \mathcal{O}_{Y_n}(q)) = 0$. Cela résulte du lemme \ref{lemma-bound-q-and-d}.

\medskip\noindent
Démonstration de (2). Supposons $d > 1$. Il faut prouver que $H^{d - 1}(X, p^*M(q)) = 0$ pour $q \geq q_0$. En raisonnant comme au paragraphe précédent, il suffit de montrer que $H^d(X, \mathcal{G}(q)) = 0$. À l'aide de la première suite exacte courte et de l'annulation de la cohomologie des tordus de $\mathcal{F}$ en degrés $> d$, on voit qu'il suffit de montrer que $H^d(X, p^*K(q))$ est nul, ce qui résulte de (1) et du fait que $K$ est un $(A, n, c)$-module (voir ci-dessus).

\medskip\noindent
Démonstration de (3). Soit $0 < i < d - 1$, et supposons l'assertion établie pour $i + 1$, sauf dans le cas $i = d - 2$, où l'on dispose de (2). La suite exacte longue de cohomologie associée à la seconde suite exacte courte ci-dessus fournit une injection
$$
H^i(X, p^*M(q - (d - 1 - i)c)) \subset
H^{i + 1}(X, \mathcal{G}(q - (d - 1 - i)c))
$$
puisque $q - (d - 1 - i)c \geq q_0$ entraîne l'annulation de $H^i(X, \mathcal{O}_{Y_n}(q - (d - 1 - i)c))$ (voir ci-dessus). Il suffit donc de montrer que l'application $H^{i + 1}(X, \mathcal{G}(q)) \to H^{i + 1}(X, \mathcal{G}(q - (d - 1 - i)c))$ est nulle. Pour l'étudier, considérons les morphismes de suites exactes
\begingroup\footnotesize
$$
\xymatrix{
H^{i + 1}(X, p^*K(q)) \ar[r] \ar[d] &
H^{i + 1}(X, \mathcal{G}(q)) \ar[r] \ar[d] \ar@{..>}[ld] &
H^{i + 2}(X, \mathcal{F}(q)) \ar[d] \\
H^{i + 1}(X, p^*K(q - c)) \ar[r] \ar[d] &
H^{i + 1}(X, \mathcal{G}(q - c)) \ar[r] \ar[d] &
H^{i + 2}(X, \mathcal{F}(q - c)) \\
H^{i + 1}(X, p^*K(q - (d - 1 - i)c)) \ar[r] &
H^{i + 1}(X, \mathcal{G}(q - (d - 1 - i)c))
}
$$
\endgroup
Puisque $\mathcal{F}$ est à support schématique dans $Y_c$, l'application canonique $\mathcal{G}(q) \to \mathcal{G}(q - c)$ se factorise par $p^*K(q - c)$ d'après le lemme \ref{lemma-annihilated}. Cela donne la flèche pointillée du diagramme. (En fait, pour la démonstration, il suffit de remarquer que la flèche verticale à l'extrême droite est nulle pour obtenir la flèche pointillée comme application d'ensembles.) Il suffit donc de montrer que $H^{i + 1}(X, p^*K(q - c)) \to
H^{i + 1}(X, p^*K(q - (d - 1 - i)c))$ est nulle. Si $i = d - 2$, la source de cette flèche est nulle d'après (2), puisque $q - c \geq q_0$ et que $K$ est un $(A, n, c)$-module. Si $i < d - 2$, alors, puisque $K$ est un $(A, n, c)$-module, l'hypothèse de récurrence montre que l'application est bien nulle, car $q - c - (q - (d - 1 - i)c) = (d - 2 - i)c = (d - 1 - (i + 1))c$ et $q - c \geq q_0 + (d - 1 - (i + 1))c$. Cela achève la démonstration de (3).

\medskip\noindent
Démonstration de (4). Supposons $d \in \{0, 1\}$ et $q \geq q_0$. La première suite exacte courte donne alors une surjection $H^1(X, p^*K(q)) \to H^1(X, \mathcal{G}(q))$, dont la source est nulle d'après le cas (1). Ainsi, pour tout $q \in \mathbf{Z}$, l'application
$$
H^0(X, (\mathcal{O}_{Y_n})^{\oplus r}(q))
\longrightarrow
H^0(X, p^*M(q))
$$
est surjective. Pour $q \geq q_0$, sa source est égale à $(I^q/I^{q + n})^{\oplus r}$ d'après le lemme \ref{lemma-bound-q-and-d}, ce qui prouve aisément l'assertion.

\medskip\noindent
Démonstration de (5). Supposons $d > 1$. En raisonnant comme dans la démonstration de (4), il suffit de montrer que l'image de
$$
H^0(X, p^*M(q))
\longrightarrow
H^0(X, p^*M(q - (d - 1)c))
$$
est contenue dans l'image de
$$
H^0(X, (\mathcal{O}_{Y_n})^{\oplus r}(q - (d - 1)c))
\longrightarrow
H^0(X, p^*M(q - (d - 1)c))
$$
Pour montrer cette inclusion, il suffit de montrer que, pour $\sigma \in H^0(X, p^*M(q))$ de bord $\xi \in H^1(X, \mathcal{G}(q))$, l'image de $\xi$ dans $H^1(X, \mathcal{G}(q - (d - 1)c))$ est nulle. Cela résulte des mêmes arguments que dans la démonstration de (3).
\end{proof}

\begin{remark}
\label{remark-duals}
Étant donné un couple $(M, n)$ constitué d'un entier $n \geq 0$ et d'un $A/I^n$-module de type fini $M$, posons $M^\vee = \Hom_{A/I^n}(M, A/I^n)$. Étant donné un couple $(\mathcal{F}, n)$ constitué d'un entier $n$ et d'un $\mathcal{O}_{Y_n}$-module cohérent $\mathcal{F}$, posons
$$
\mathcal{F}^\vee =
\SheafHom_{\mathcal{O}_{Y_n}}(\mathcal{F}, \mathcal{O}_{Y_n})
$$
Pour $(M, n)$ comme ci-dessus, il existe une application canonique
$$
can : p^*(M^\vee) \longrightarrow (p^*M)^\vee
$$
En effet, si l'on choisit une présentation $(A/I^n)^{\oplus s} \to (A/I^n)^{\oplus r} \to M \to 0$, on obtient une présentation $\mathcal{O}_{Y_n}^{\oplus s} \to \mathcal{O}_{Y_n}^{\oplus r} \to
p^*M \to 0$. En passant aux duaux, on obtient des suites exactes
$$
0 \to M^\vee \to (A/I^n)^{\oplus r} \to (A/I^n)^{\oplus s}
$$
et
$$
0 \to (p^*M)^\vee \to
\mathcal{O}_{Y_n}^{\oplus r} \to
\mathcal{O}_{Y_n}^{\oplus s}
$$
En prenant l'image réciproque de la première suite par $p$, on trouve l'application voulue $can$. La construction de cette application est fonctorielle en le $A/I^n$-module de type fini $M$. Le noyau et le conoyau de $can$ sont à support schématique dans $Y_c$ si $M$ est un $(A, n, c)$-module. En effet, dans ce cas, pour $a \in I^c$, l'application $a : M \to M$ se factorise par un $A/I^n$-module libre de type fini, pour lequel $can$ est un isomorphisme. Ainsi $a$ annule le noyau et le conoyau de $can$.
\end{remark}

\begin{lemma}
\label{lemma-factor-hom}
Avec $q_0 = q(S)$ et $d = d(S)$ comme ci-dessus, soient $M$ un $(A, n, c)$-module et $\varphi : M \to I^n/I^{2n}$ une application $A$-linéaire. Supposons $n \geq \max(q_0 + (1 + d)c, (2 + d)c)$ et, si $d = 0$, supposons $n \geq q_0 + 2c$. Alors le composé
$$
M \xrightarrow{\varphi} I^n/I^{2n} \to
I^{n - (1 + d)c}/I^{2n - (1 + d)c}
$$
est de la forme $\sum a_i \psi_i$, avec $a_i \in I^c$ et $\psi_i : M \to I^{n - (2 + d)c}/I^{2n - (2 + d)c}$.
\end{lemma}

\begin{proof}
Cas $d > 1$. Puisque l'on dispose d'un système compatible d'applications $p^*(I^q) \to \mathcal{O}_X(q)$ pour $q \geq 0$, il existe des applications canoniques $p^*(I^q/I^{q + \nu}) \to \mathcal{O}_{Y_\nu}(q)$ pour $\nu \geq 0$. En utilisant cela et en prenant l'image réciproque de $\varphi$, on obtient une application
$$
\chi : p^*M \longrightarrow \mathcal{O}_{Y_n}(n)
$$
telle que le composé $M \to H^0(X, p^*M) \to H^0(X, \mathcal{O}_{Y_n}(n))$ soit l'homomorphisme donné $\varphi$ suivi de l'application $I^n/I^{2n} \to H^0(X, \mathcal{O}_{Y_n}(n))$. Puisque $\mathcal{O}_{Y_n}(n)$ est inversible sur $Y_n$, l'application linéaire $\chi$ détermine une section
$$
\sigma \in \Gamma(X, (p^*M)^\vee(n))
$$
avec les notations de la remarque \ref{remark-duals}. La discussion de la remarque \ref{remark-duals} montre que le conoyau et le noyau de $can : p^*(M^\vee) \to (p^*M)^\vee$ sont à support schématique dans $Y_c$. D'après le lemme \ref{lemma-annihilated}, l'application $(p^*M)^\vee(n) \to (p^*M)^\vee(n - 2c)$ se factorise par $p^*(M^\vee)(n - 2c)$ ; nous omettons un petit détail. L'image de $\sigma$ dans $\Gamma(X, (p^*M)^\vee(n - 2c))$ provient donc d'un élément
$$
\sigma' \in \Gamma(X, p^*(M^\vee)(n - 2c))
$$
D'après le lemme \ref{lemma-vanishing-coh-almost-projective}, point (5), le fait que $M^\vee$ soit un $(A, n, c)$-module en vertu de Compléments d'algèbre, lemme \ref{more-algebra-lemma-dual-near-projective}, et le fait que $n \geq q_0 + (1 + d)c$, donc $n - 2c \geq q_0 + (d - 1)c$, on voit que l'image de $\sigma'$ dans $H^0(X, p^*M^\vee(n - (1 + d)c))$ est l'image d'un élément $\tau$ de $I^{n - (1 + d)c}M^\vee$. Écrivons $\tau = \sum a_i \tau_i$ avec $\tau_i \in I^{n - (2 + d)c}M^\vee$ ; cela a un sens puisque $n - (2 + d)c \geq 0$. Alors $\tau_i$ détermine un homomorphisme de modules $\psi_i : M \to I^{n - (2 + d)c}/I^{2n - (2 + d)c}$ à l'aide de l'application d'évaluation $M \otimes M^\vee \to A/I^n$.

\medskip\noindent
Montrons que cela convient\footnote{Nous espérons qu'un lecteur proposera une démonstration moins laborieuse de ce fait.}. Choisissons $z \in M$ et montrons que $\varphi(z)$ et $\sum a_i \psi_i(z)$ ont la même image dans $I^{n - (1 + d)c}/I^{2n - (1 + d)c}$. D'abord, l'élément $z$ détermine une application $p^*z : \mathcal{O}_{Y_n} \to p^*M$ dont le composé avec $\chi$ est l'application $\mathcal{O}_{Y_n} \to \mathcal{O}_{Y_n}(n)$ correspondant à $\varphi(z)$ par l'application $I^n/I^{2n} \to \Gamma(\mathcal{O}_{Y_n}(n))$. Ensuite, $z$ et $p^*z$ déterminent les applications d'évaluation $e_z : M^\vee \to A/I^n$ et $e_{p^*z} : (p^*M)^\vee \to \mathcal{O}_{Y_n}$. Puisque $\chi(p^*z)$ est la section correspondant à $\varphi(z)$, on voit que $e_{p^*z}(\sigma)$ est la section correspondant à $\varphi(z)$. Ici et dans la suite, nous faisons un abus de notation : pour une application $a : \mathcal{F} \to \mathcal{G}$ de modules sur $X$, nous notons aussi $a : \mathcal{F}(t) \to \mathcal{F}(t)$ l'application correspondante entre les modules tordus. Le diagramme
$$
\xymatrix{
p^*(M^\vee) \ar[d]_{can} \ar[r]_{p^*e_z} & \mathcal{O}_{Y_n} \ar@{=}[d] \\
(p^*M)^\vee \ar[r]^{e_{p^*z}} & \mathcal{O}_{Y_n}
}
$$
{\emergencystretch=2em
est commutatif par fonctorialité de la construction $can$. Ainsi $(p^*e_z)(\sigma')$, dans $\Gamma(Y_n, \mathcal{O}_{Y_n}(n - 2c))$, est la section correspondant à l'image de $\varphi(z)$ dans $I^{n - 2c}/I^{2n - 2c}$. L'étape suivante est que $\sigma'$ s'envoie sur l'image de $\sum a_i \tau_i$ dans $H^0(X, p^*M^\vee(n - (1 + d)c))$. Cela implique que $(p^*e_z)(\sum a_i \tau_i) = \sum a_i p^*e_z(\tau_i)$, dans $\Gamma(Y_n, \mathcal{O}_{Y_n}(n - (1 + d)c))$, est la section correspondant à l'image de $\varphi(z)$ dans $I^{n - (1 + d)c}/I^{2n - (1 + d)c}$. Rappelons que $\psi_i$ est défini à partir de $\tau_i$ par une application d'évaluation. Si l'on note
\par}
$$
\chi_i : p^*M \longrightarrow \mathcal{O}_{Y_n}(n - (2 + d)c)
$$
l'application obtenue à partir de $\psi_i$, le même raisonnement que ci-dessus montre que la section correspondant à $\psi_i(z)$ est $\chi_i(p^*z) = e_{p^*z}(\chi_i) = p^*e_z(\tau_i)$. On en déduit que l'image de $\varphi(z)$ dans $\Gamma(Y_n, \mathcal{O}_{Y_n}(n - (1 + d)c))$ est égale à celle de $\sum a_i\psi_i(z)$. Puisque $n - (1 + d)c \geq q_0$, on a $\Gamma(Y_n, \mathcal{O}_{Y_n}(n - (1 + d)c)) =
I^{n - (1 + d)c}/I^{2n - (1 + d)c}$ d'après le lemme \ref{lemma-bound-q-and-d}, ce qui prouve la compatibilité voulue.

\medskip\noindent
Cas $d = 1$. On raisonne comme ci-dessus pour obtenir
$$
\chi : p^*M \longrightarrow \mathcal{O}_{Y_n}(n),\quad
\sigma \in \Gamma(X, (p^*M)^\vee(n)),\quad
\sigma' \in \Gamma(X, p^*(M^\vee)(n - 2c)),
$$
puis on utilise le lemme \ref{lemma-vanishing-coh-almost-projective}, point (4), pour voir que $\sigma'$ est l'image d'un élément $\tau \in I^{n - 2c}M^\vee$. Le reste de l'argument est identique.

\medskip\noindent
Cas $d = 0$. L'argument est exactement le même que dans le cas $d = 1$.
\end{proof}

\begin{lemma}
\label{lemma-bound-two-term-complex}
Avec $d = d(S)$ et $q_0 = q(S)$ comme ci-dessus,
\begin{enumerate}
\item pour des entiers $n \geq c \geq 0$ tels que $n \geq \max(q_0 + (1 + d)c, (2 + d)c)$,
\item pour $K$ dans $D(A/I^n)$ tel que $H^i(K) = 0$ pour $i \not = -1, 0$ et $H^i(K)$ de type fini pour $i = -1, 0$, et tel que $\Ext^1_{A/I^c}(K, N)$ soit annulé par $I^c$ pour tous les $A/I^n$-modules de type fini $N$,
\end{enumerate}
l'application
$$
\Ext^1_{A/I^n}(K, I^n/I^{2n})
\longrightarrow
\Ext^1_{A/I^n}(K, I^{n - (1 + d)c}/I^{2n - 2(1 + d)c})
$$
est nulle.
\end{lemma}

\begin{proof}
Cas $d > 0$. Soit $K^{-1} \to K^0$ un complexe représentant $K$ comme dans Compléments d'algèbre, lemme \ref{more-algebra-lemma-ext-1-annihilated-definite}, point (5), relativement à l'idéal $I^c/I^n$ de l'anneau $A/I^n$. En particulier, $K^{-1}$ est $I^c/I^n$-projectif, puisque la multiplication par les éléments de $I^c/I^n$ se factorise même par $K^0$. D'après Compléments d'algèbre, lemme \ref{more-algebra-lemma-map-out-of-almost-free}, point (1), on a
$$
\Ext^1_{A/I^n}(K, I^n/I^{2n}) =
\Coker(\Hom_{A/I^n}(K^0, I^n/I^{2n}) \to \Hom_{A/I^n}(K^{-1}, I^n/I^{2n}))
$$
{\emergencystretch=8em
et de même pour les autres groupes Ext. Toute classe $\xi$ dans $\Ext^1_{A/I^n}(K, I^n/I^{2n})$ provient donc d'un élément $\varphi \in \Hom_{A/I^n}(K^{-1}, I^n/I^{2n})$. Notons $\varphi'$ l'image de $\varphi$ dans $\Hom_{A/I^n}(K^{-1}, I^{n - (1 + d)c}/I^{2n - (1 + d)c})$. D'après le lemme \ref{lemma-factor-hom}, on peut écrire $\varphi' = \sum a_i \psi_i$ avec $a_i \in I^c$ et $\psi_i \in \Hom_{A/I^n}(M, I^{n - (2 + d)c}/I^{2n - (2 + d)c})$. Choisissons $h_i : K^0 \to K^{-1}$ tel que $a_i \text{id}_{K^{-1}} = h_i \circ d_K^{-1}$. Posons $\psi = \sum \psi_i \circ h_i : K^0 \to I^{n - (2 + d)c}/I^{2n - (2 + d)c}$. Alors $\varphi' = \psi \circ \text{d}_K^{-1}$, et l'on en déduit que $\xi$ a déjà une image nulle dans $\Ext^1_{A/I^n}(K, I^{n - (1 + d)c}/I^{2n - (1 + d)c})$, et a fortiori dans $\Ext^1_{A/I^n}(K, I^{n - (1 + d)c}/I^{2n - 2(1 + d)c})$.
\par}

\medskip\noindent
Cas $d = 0$\footnote{L'argument donné pour $d > 0$ convient, mais fournit un résultat légèrement plus faible.}. Soient $\xi$ et $\varphi$ comme ci-dessus. Considérons le diagramme
$$
\xymatrix{
K^0 \\
K^{-1} \ar[u] \ar[r]^\varphi & I^n/I^{2n} \ar[r] & I^{n - c}/I^{2n - c}
}
$$
En prenant l'image réciproque sur $X$ et en utilisant l'application $p^*(I^n/I^{2n}) \to \mathcal{O}_{Y_n}(n)$, on obtient un diagramme à flèches pleines
$$
\xymatrix{
p^*K^0 \ar@{..>}[rrd] \\
p^*K^{-1} \ar[u] \ar[r] &
\mathcal{O}_{Y_n}(n) \ar[r] &
\mathcal{O}_{Y_n}(n - c)
}
$$
On peut recouvrir $X$ par des ouverts affines $U = \Spec(B)$ tels qu'il existe $a \in I$ vérifiant la propriété suivante : $IB = aB$ et $a$ est un non-diviseur de zéro sur $B$. En effet, on peut recouvrir $X$ par les spectres d'algèbres d'éclatement affines ; voir Diviseurs, lemme \ref{divisors-lemma-blowing-up-affine}. La restriction de $\mathcal{O}_{Y_n}(n) \to \mathcal{O}_{Y_n}(n - c)$ à $U$ est isomorphe à l'application de $\mathcal{O}_U$-modules quasi-cohérents correspondant à l'homomorphisme de $B$-modules $a^c : B/a^nB \to B/a^nB$. Puisque $a^c : K^{-1} \to K^{-1}$ se factorise par $K^0$, la flèche pointillée existe sur $U$. Autrement dit, on peut trouver la flèche pointillée localement sur $X$ ! Or le faisceau des flèches pointillées complétant le diagramme est un espace principal homogène sous
$$
\mathcal{F} = \SheafHom_{\mathcal{O}_X}(
\Coker(p^*K^{-1} \to p^*K^0), \mathcal{O}_{Y_n}(n - c))
$$
qui est un $\mathcal{O}_X$-module cohérent. L'obstruction à l'existence de la flèche pointillée est donc un élément de $H^1(X, \mathcal{F})$. Ce groupe de cohomologie est nul puisque $1 > d = 0$ ; voir la discussion qui suit la définition de $d = d(S)$. On peut donc trouver une flèche pointillée $\psi : p^*K^0 \to \mathcal{O}_{Y_n}(n - c)$ complétant le diagramme. Puisque $n - c \geq q_0$, on voit que $\psi$ induit une application $K^0 \to I^{n - c}/I^{2n - c}$. Une chasse au diagramme donne $\varphi' = \psi \circ \text{d}_K^{-1}$, et la démonstration s'achève comme précédemment.
\end{proof}






\clearpage
\enlargethispage{\smallskipamount}
\begin{multicols}{2}[\section{Autres chapitres}]
\noindent
Préliminaires
\begin{enumerate}
\item \hyperref[introduction-section-phantom]{Introduction}
\item \hyperref[conventions-section-phantom]{Conventions}
\item \hyperref[sets-section-phantom]{Théorie des ensembles}
\item \hyperref[categories-section-phantom]{Catégories}
\item \hyperref[topology-section-phantom]{Topologie}
\item \hyperref[sheaves-section-phantom]{Faisceaux sur les espaces}
\item \hyperref[sites-section-phantom]{Sites et faisceaux}
\item \hyperref[stacks-section-phantom]{Champs}
\item \hyperref[fields-section-phantom]{Corps}
\item \hyperref[algebra-section-phantom]{Algèbre commutative}
\item \hyperref[brauer-section-phantom]{Groupes de Brauer}
\item \hyperref[homology-section-phantom]{Algèbre homologique}
\item \hyperref[derived-section-phantom]{Catégories dérivées}
\item \hyperref[simplicial-section-phantom]{Méthodes simpliciales}
\item \hyperref[more-algebra-section-phantom]{Compléments d'algèbre}
\item \hyperref[smoothing-section-phantom]{Lissage des morphismes d'anneaux}
\item \hyperref[modules-section-phantom]{Faisceaux de modules}
\item \hyperref[sites-modules-section-phantom]{Modules sur les sites}
\item \hyperref[injectives-section-phantom]{Objets injectifs}
\item \hyperref[cohomology-section-phantom]{Cohomologie des faisceaux}
\item \hyperref[sites-cohomology-section-phantom]{Cohomologie sur les sites}
\item \hyperref[dga-section-phantom]{Algèbre différentielle graduée}
\item \hyperref[dpa-section-phantom]{Algèbre à puissances divisées}
\item \hyperref[sdga-section-phantom]{Faisceaux différentiels gradués}
\item \hyperref[hypercovering-section-phantom]{Hyperrecouvrements}
\end{enumerate}
Schémas
\begin{enumerate}
\setcounter{enumi}{25}
\item \hyperref[schemes-section-phantom]{Schémas}
\item \hyperref[constructions-section-phantom]{Constructions de schémas}
\item \hyperref[properties-section-phantom]{Propriétés des schémas}
\item \hyperref[morphisms-section-phantom]{Morphismes de schémas}
\item \hyperref[coherent-section-phantom]{Cohomologie des schémas}
\item \hyperref[divisors-section-phantom]{Diviseurs}
\item \hyperref[limits-section-phantom]{Limites de schémas}
\item \hyperref[varieties-section-phantom]{Variétés}
\item \hyperref[topologies-section-phantom]{Topologies sur les schémas}
\item \hyperref[descent-section-phantom]{Descente}
\item \hyperref[perfect-section-phantom]{Catégories dérivées de schémas}
\item \hyperref[more-morphisms-section-phantom]{Compléments sur les morphismes}
\item \hyperref[flat-section-phantom]{Compléments sur la platitude}
\item \hyperref[groupoids-section-phantom]{Schémas en groupoïdes}
\item \hyperref[more-groupoids-section-phantom]{Compléments sur les schémas en groupoïdes}
\item \hyperref[etale-section-phantom]{Morphismes étales de schémas}
\end{enumerate}
Thèmes de théorie des schémas
\begin{enumerate}
\setcounter{enumi}{41}
\item \hyperref[chow-section-phantom]{Homologie de Chow}
\item \hyperref[intersection-section-phantom]{Théorie de l'intersection}
\item \hyperref[pic-section-phantom]{Schémas de Picard des courbes}
\item \hyperref[weil-section-phantom]{Théories cohomologiques de Weil}
\item \hyperref[adequate-section-phantom]{Modules adéquats}
\item \hyperref[dualizing-section-phantom]{Complexes dualisants}
\item \hyperref[duality-section-phantom]{Dualité pour les schémas}
\item \hyperref[discriminant-section-phantom]{Discriminants et différentes}
\item \hyperref[derham-section-phantom]{Cohomologie de de Rham}
\item \hyperref[local-cohomology-section-phantom]{Cohomologie locale}
\item \hyperref[algebraization-section-phantom]{Géométrie algébrique et formelle}
\item \hyperref[curves-section-phantom]{Courbes algébriques}
\item \hyperref[resolve-section-phantom]{Résolution des surfaces}
\item \hyperref[models-section-phantom]{Réduction semi-stable}
\item \hyperref[functors-section-phantom]{Foncteurs et morphismes}
\item \hyperref[equiv-section-phantom]{Catégories dérivées de variétés}
\item \hyperref[pione-section-phantom]{Groupes fondamentaux de schémas}
\item \hyperref[etale-cohomology-section-phantom]{Cohomologie étale}
\item \hyperref[crystalline-section-phantom]{Cohomologie cristalline}
\item \hyperref[proetale-section-phantom]{Cohomologie pro-étale}
\item \hyperref[relative-cycles-section-phantom]{Cycles relatifs}
\item \hyperref[more-etale-section-phantom]{Compléments de cohomologie étale}
\item \hyperref[trace-section-phantom]{La formule des traces}
\end{enumerate}
Espaces algébriques
\begin{enumerate}
\setcounter{enumi}{64}
\item \hyperref[spaces-section-phantom]{Espaces algébriques}
\item \hyperref[spaces-properties-section-phantom]{Propriétés des espaces algébriques}
\item \hyperref[spaces-morphisms-section-phantom]{Morphismes d'espaces algébriques}
\item \hyperref[decent-spaces-section-phantom]{Espaces algébriques décents}
\item \hyperref[spaces-cohomology-section-phantom]{Cohomologie des espaces algébriques}
\item \hyperref[spaces-limits-section-phantom]{Limites d'espaces algébriques}
\item \hyperref[spaces-divisors-section-phantom]{Diviseurs sur les espaces algébriques}
\item \hyperref[spaces-over-fields-section-phantom]{Espaces algébriques sur des corps}
\item \hyperref[spaces-topologies-section-phantom]{Topologies sur les espaces algébriques}
\item \hyperref[spaces-descent-section-phantom]{Descente et espaces algébriques}
\item \hyperref[spaces-perfect-section-phantom]{Catégories dérivées d'espaces}
\item \hyperref[spaces-more-morphisms-section-phantom]{Compléments sur les morphismes d'espaces}
\item \hyperref[spaces-flat-section-phantom]{Platitude sur les espaces algébriques}
\item \hyperref[spaces-groupoids-section-phantom]{Groupoïdes dans les espaces algébriques}
\item \hyperref[spaces-more-groupoids-section-phantom]{Compléments sur les groupoïdes dans les espaces}
\item \hyperref[bootstrap-section-phantom]{Amorçage}
\item \hyperref[spaces-pushouts-section-phantom]{Sommes amalgamées d'espaces algébriques}
\end{enumerate}
Thèmes de géométrie
\begin{enumerate}
\setcounter{enumi}{81}
\item \hyperref[spaces-chow-section-phantom]{Groupes de Chow des espaces}
\item \hyperref[groupoids-quotients-section-phantom]{Quotients de groupoïdes}
\item \hyperref[spaces-more-cohomology-section-phantom]{Compléments sur la cohomologie des espaces}
\item \hyperref[spaces-simplicial-section-phantom]{Espaces simpliciaux}
\item \hyperref[spaces-duality-section-phantom]{Dualité pour les espaces}
\item \hyperref[formal-spaces-section-phantom]{Espaces algébriques formels}
\item \hyperref[restricted-section-phantom]{Algébrisation des espaces formels}
\item \hyperref[spaces-resolve-section-phantom]{Résolution des surfaces revisitée}
\end{enumerate}
Théorie des déformations
\begin{enumerate}
\setcounter{enumi}{89}
\item \hyperref[formal-defos-section-phantom]{Théorie formelle des déformations}
\item \hyperref[defos-section-phantom]{Théorie des déformations}
\item \hyperref[cotangent-section-phantom]{Le complexe cotangent}
\item \hyperref[examples-defos-section-phantom]{Problèmes de déformation}
\end{enumerate}
Champs algébriques
\begin{enumerate}
\setcounter{enumi}{93}
\item \hyperref[algebraic-section-phantom]{Champs algébriques}
\item \hyperref[examples-stacks-section-phantom]{Exemples de champs}
\item \hyperref[stacks-sheaves-section-phantom]{Faisceaux sur les champs algébriques}
\item \hyperref[criteria-section-phantom]{Critères de représentabilité}
\item \hyperref[artin-section-phantom]{Axiomes d'Artin}
\item \hyperref[quot-section-phantom]{Espaces de Quot et de Hilbert}
\item \hyperref[stacks-properties-section-phantom]{Propriétés des champs algébriques}
\item \hyperref[stacks-morphisms-section-phantom]{Morphismes de champs algébriques}
\item \hyperref[stacks-limits-section-phantom]{Limites de champs algébriques}
\item \hyperref[stacks-cohomology-section-phantom]{Cohomologie des champs algébriques}
\item \hyperref[stacks-perfect-section-phantom]{Catégories dérivées de champs}
\item \hyperref[stacks-introduction-section-phantom]{Introduction aux champs algébriques}
\item \hyperref[stacks-more-morphisms-section-phantom]{Compléments sur les morphismes de champs}
\item \hyperref[stacks-geometry-section-phantom]{La géométrie des champs}
\end{enumerate}
Thèmes de théorie des espaces de modules
\begin{enumerate}
\setcounter{enumi}{107}
\item \hyperref[moduli-section-phantom]{Champs de modules}
\item \hyperref[moduli-curves-section-phantom]{Espaces de modules de courbes}
\end{enumerate}
Divers
\begin{enumerate}
\setcounter{enumi}{109}
\item \hyperref[examples-section-phantom]{Exemples}
\item \hyperref[exercises-section-phantom]{Exercices}
\item \hyperref[guide-section-phantom]{Guide bibliographique}
\item \hyperref[desirables-section-phantom]{Desiderata}
\item \hyperref[coding-section-phantom]{Style de programmation}
\item \hyperref[obsolete-section-phantom]{Obsolète}
\item \hyperref[fdl-section-phantom]{Licence de documentation libre GNU}
\item \hyperref[index-section-phantom]{Index généré automatiquement}
\end{enumerate}
\end{multicols}


\bibliography{my}
\bibliographystyle{amsalpha}

\end{document}
